% !TeX program = xelatex
% Standalone English volume. Use XeLaTeX until the references stabilize.
\documentclass[a4paper,11pt]{article}
\usepackage[T1]{fontenc}
\usepackage{lmodern}
\usepackage[english]{babel}
\usepackage[left=26mm,right=26mm,top=25mm,bottom=25mm,headheight=15pt,headsep=7mm,footskip=11mm]{geometry}
\usepackage{amsmath,amssymb,mathtools,amsthm}
\newtheorem{proposition}{Proposition}[section]
\newtheorem{lemma}[proposition]{Lemma}
\newtheorem{theorem}[proposition]{Theorem}
\newtheorem{corollary}[proposition]{Corollary}
\newtheorem{maintheorem}[proposition]{Main Theorem}
\theoremstyle{definition}
\newtheorem{definition}[proposition]{Definition}
\theoremstyle{remark}
\newtheorem{remark}[proposition]{Remark}
\makeatletter
\renewcommand*\l@section{\@dottedtocline{1}{0em}{2.4em}}
\renewcommand*\l@subsection{\@dottedtocline{2}{1.45em}{3.45em}}
\renewcommand*\l@subsubsection{\@dottedtocline{3}{3.8em}{4.55em}}
\makeatother
\usepackage{xcolor,fancyhdr,titlesec,needspace,xurl}
\usepackage{longtable,multicol}
\usepackage{hyperref}
\definecolor{ink}{HTML}{1A3542}
\definecolor{muted}{HTML}{4F606A}
\hypersetup{unicode=true,colorlinks=true,linkcolor=ink,filecolor=ink,urlcolor=ink,citecolor=ink,bookmarksnumbered=true,
 pdfdisplaydoctitle=true,pdflang={en},
 pdftitle={The generalized Riemann hypothesis for Dirichlet L-functions. Edition 930, Volume VII},
 pdfsubject={A natural-language reconstruction of the proposed proof},
 pdfauthor={A mathematical reconstruction}}
\titleformat{\section}{\Large\bfseries\color{ink}}{\thesection}{.65em}{}
\titlespacing*{\section}{0pt}{.8em}{.65em}
\setlength{\parindent}{1.2em}
\linespread{1.045}
\setlength{\parskip}{.22em plus .05em minus .03em}
\setlength{\emergencystretch}{2em}
\setcounter{topnumber}{2}
\setcounter{tocdepth}{3}\setcounter{secnumdepth}{3}
\widowpenalty=4500
\clubpenalty=4500
\pagestyle{fancy}\fancyhf{}
\fancyhead[L]{\small\color{muted}Dirichlet GRH \textperiodcentered\ Volume I}
\fancyhead[R]{\small\color{muted}Dirichlet zeros and power intervals}
\fancyfoot[C]{\small\thepage}
\renewcommand{\headrulewidth}{.3pt}
\newcommand{\lead}[1]{\par\smallskip\noindent{\bfseries\color{ink}#1.}\enspace}
\newcommand{\sour}[2]{\unskip\penalty150\hspace{.35em}{\small\hyperlink{ref#1}{[#2]}}}
\titleformat{\subsection}{\large\bfseries}{\thesubsection}{.6em}{}
\titlespacing*{\subsection}{0pt}{1.15em plus .15em}{.4em}
\setlength{\abovedisplayskip}{9pt plus 2pt minus 3pt}
\setlength{\belowdisplayskip}{9pt plus 2pt minus 3pt}
\newcommand{\RR}{\mathbb R}\newcommand{\CC}{\mathbb C}\newcommand{\QQ}{\mathbb Q}
\newcommand{\ZZ}{\mathbb Z}\newcommand{\ord}{\operatorname{ord}}
\newcommand{\Frob}{\operatorname{Frob}}\newcommand{\Gal}{\operatorname{Gal}}
\newcommand{\len}{\operatorname{length}}\newcommand{\Res}{\operatorname*{Res}}
\newcommand{\wmean}{\operatorname{av}}\newcommand{\Ker}{\operatorname{ker}}
\newcommand{\Lh}{\mathcal L}\newcommand{\Gg}{\mathcal G}
\newcommand{\real}{\operatorname{Re}}\newcommand{\imag}{\operatorname{Im}}
\newcommand{\abs}[1]{\left|#1\right|}\newcommand{\eps}{\varepsilon}
\newcommand{\anchoritem}[2]{\par\Needspace{2.8\baselineskip}\hypertarget{ref#1}{}\noindent{\bfseries\color{ink}#2.}\enspace}
\newcommand{\srcloc}[2]{{\small\nolinkurl{#1.lean}}, line~#2}
\numberwithin{equation}{section}
\raggedbottom

\newcommand{\Sel}{\mathcal S}
\newcommand{\Cl}{\operatorname{Cl}}
\newcommand{\Aut}{\operatorname{Aut}}
\newcommand{\GL}{\operatorname{GL}}
\newcommand{\OO}{\mathcal O}
\hypersetup{bookmarksdepth=3}


\addto\captionsenglish{\renewcommand{\partname}{Volume}}


\makeatletter
\expandafter\def\csname r@cor:interval238\endcsname{{21.8}{45}{}{proposition.21.8}{Dirichlet-GRH-v930-Volume-IV.pdf}}
\expandafter\def\csname r@cor:present116\endcsname{{32.20}{45}{}{proposition.32.20}{Dirichlet-GRH-v930-Volume-VI.pdf}}
\expandafter\def\csname r@cor:tangent119\endcsname{{33.6}{52}{}{proposition.33.6}{Dirichlet-GRH-v930-Volume-VI.pdf}}
\expandafter\def\csname r@eq:word68\endcsname{{32.1}{31}{}{equation.32.1}{Dirichlet-GRH-v930-Volume-VI.pdf}}
\expandafter\def\csname r@lem:component68\endcsname{{32.7}{32}{}{proposition.32.7}{Dirichlet-GRH-v930-Volume-VI.pdf}}
\expandafter\def\csname r@lem:cusps240\endcsname{{21.19}{56}{}{proposition.21.19}{Dirichlet-GRH-v930-Volume-IV.pdf}}
\expandafter\def\csname r@lem:kummer72\endcsname{{32.14}{39}{}{proposition.32.14}{Dirichlet-GRH-v930-Volume-VI.pdf}}
\expandafter\def\csname r@lem:spin122\endcsname{{33.14}{60}{}{proposition.33.14}{Dirichlet-GRH-v930-Volume-VI.pdf}}
\expandafter\def\csname r@nav:III-00\endcsname{{30.4}{8}{}{subsection.30.4}{Dirichlet-GRH-v930-Volume-VI.pdf}}
\expandafter\def\csname r@prop:finite15\endcsname{{11.4}{6}{}{proposition.11.4}{Dirichlet-GRH-v930-Volume-II.pdf}}
\expandafter\def\csname r@prop:res23\endcsname{{6.11}{32}{}{proposition.6.11}{Dirichlet-GRH-v930-Volume-I.pdf}}
\expandafter\def\csname r@sec:chart69\endcsname{{32.5}{34}{}{subsection.32.5}{Dirichlet-GRH-v930-Volume-VI.pdf}}
\expandafter\def\csname r@sec:levelstrict345\endcsname{{22.13}{74}{}{subsection.22.13}{Dirichlet-GRH-v930-Volume-IV.pdf}}
\expandafter\def\csname r@sec:missing\endcsname{{6.6}{32}{}{subsection.6.6}{Dirichlet-GRH-v930-Volume-I.pdf}}
\expandafter\def\csname r@sec:ordinaryorth122\endcsname{{33.5}{59}{}{subsection.33.5}{Dirichlet-GRH-v930-Volume-VI.pdf}}
\expandafter\def\csname r@sec:original68\endcsname{{32.4}{31}{}{subsection.32.4}{Dirichlet-GRH-v930-Volume-VI.pdf}}
\expandafter\def\csname r@seriesguide\endcsname{{1.1}{4}{}{subsection.1.1}{Dirichlet-GRH-v930-Volume-I.pdf}}
\expandafter\def\csname r@thm:conductor232\endcsname{{19.15}{12}{}{proposition.19.15}{Dirichlet-GRH-v930-Volume-IV.pdf}}
\expandafter\def\csname r@thm:cot116\endcsname{{32.19}{44}{}{proposition.32.19}{Dirichlet-GRH-v930-Volume-VI.pdf}}
\expandafter\def\csname r@thm:criticalmarks339\endcsname{{22.12}{71}{}{proposition.22.12}{Dirichlet-GRH-v930-Volume-IV.pdf}}
\expandafter\def\csname r@thm:crystalline121\endcsname{{33.11}{57}{}{proposition.33.11}{Dirichlet-GRH-v930-Volume-VI.pdf}}
\expandafter\def\csname r@thm:cyclo119\endcsname{{33.5}{52}{}{proposition.33.5}{Dirichlet-GRH-v930-Volume-VI.pdf}}
\expandafter\def\csname r@thm:eigen240\endcsname{{21.20}{57}{}{proposition.21.20}{Dirichlet-GRH-v930-Volume-IV.pdf}}
\expandafter\def\csname r@thm:first72\endcsname{{32.15}{40}{}{proposition.32.15}{Dirichlet-GRH-v930-Volume-VI.pdf}}
\expandafter\def\csname r@thm:fullad120\endcsname{{33.8}{54}{}{proposition.33.8}{Dirichlet-GRH-v930-Volume-VI.pdf}}
\expandafter\def\csname r@thm:integer\endcsname{{30.11}{10}{}{proposition.30.11}{Dirichlet-GRH-v930-Volume-VI.pdf}}
\expandafter\def\csname r@thm:levelstrict345\endcsname{{22.17}{75}{}{proposition.22.17}{Dirichlet-GRH-v930-Volume-IV.pdf}}
\expandafter\def\csname r@thm:markedstrict337\endcsname{{22.11}{70}{}{proposition.22.11}{Dirichlet-GRH-v930-Volume-IV.pdf}}
\expandafter\def\csname r@thm:oddtrace123\endcsname{{34.9}{71}{}{proposition.34.9}{Dirichlet-GRH-v930-Volume-VI.pdf}}
\expandafter\def\csname r@thm:orth122\endcsname{{33.13}{59}{}{proposition.33.13}{Dirichlet-GRH-v930-Volume-VI.pdf}}
\expandafter\def\csname r@thm:present31\endcsname{{30.4}{5}{}{proposition.30.4}{Dirichlet-GRH-v930-Volume-VI.pdf}}
\expandafter\def\csname r@thm:smooth69\endcsname{{32.11}{35}{}{proposition.32.11}{Dirichlet-GRH-v930-Volume-VI.pdf}}
\expandafter\def\csname r@thm:traces\endcsname{{30.8}{8}{}{proposition.30.8}{Dirichlet-GRH-v930-Volume-VI.pdf}}
\expandafter\def\csname r@thm:wild98\endcsname{{25.1}{24}{}{proposition.25.1}{Dirichlet-GRH-v930-Volume-V.pdf}}
\makeatother
\begin{document}
\setcounter{part}{6}
\fancyhead[L]{\small\color{muted}Dirichlet GRH \textperiodcentered\ Volume VII}
\fancyhead[R]{\small\color{muted}Auxiliary fields and cuspidal families}
\thispagestyle{plain}\begin{center}
{\LARGE\bfseries The generalized Riemann hypothesis\\[2pt]for Dirichlet \(L\)-functions}\\[6pt]
{\small 27 November 2026 \quad\textperiodcentered\quad Edition 930}\end{center}
\part{Auxiliary fields and ordinary cuspidal families}\label{part:VII}
\begin{abstract}\noindent An auxiliary quadratic field supplies the characters for ordinary cuspidal congruences. We follow these congruences along cyclotomic layers and recover strict classes from the traces of the three-character deformation.\end{abstract}
\setcounter{section}{35}
\noindent The \hyperref[seriesguide]{guide in Volume~I} lists all twenty-seven volumes.
\begingroup\small\setlength{\parskip}{0pt}\tableofcontents\endgroup
\clearpage

\section{Auxiliary fields and cuspidal congruences}\label{nav:g-IV-5}\label{nav:volume-VII}

\subsection{One auxiliary level and a dual cohomology class}\label{nav:III-37}
\begin{proposition}[An ordinary cuspidal lift with one new prime]
\label{prop:level241}
Use \(h,k,N,\chi,\ell,a\) from Lemma~\ref{lem:cusps240}.
Suppose \(k>2\) or \(\chi\ne1\). Let \(s\equiv1\bmod h\ell^a\)
be prime, and let \(\xi\) be any character of order \(\ell\)
modulo \(s\). For every \(m\ge0\) there is an ordinary
cuspidal newform of weight \(k'=k+m(\ell-1)\), level \(Ns\)
and character \(\chi\xi\), with
\(a_p\equiv1+\chi(p)p^{k-1}\pmod{\mathfrak m}\) for
every \(p\nmid hs\), including \(p=\ell\).
On the nonzero large-prime strict summands in
Corollary~\ref{cor:interval238} these alternatives exhaust
the cases on every sufficiently distant interval.
\end{proposition}
\begin{proof}
The difference \(E_k(1,\chi)(z)-E_k(1,\chi)(sz)\) in
Lemma~\ref{lem:cusps240} already gives a residual cusp form
at the single auxiliary prime \(s\). The constant-term
calculation there uses only that prime. The coefficient at the
first power of the Fourier variable is one. Multiply by the
\(m\)-th power of the normalized level-one Eisenstein series of
weight \(\ell-1\). Its constant term is one, and its positive
coefficients are integral multiples of \(2(\ell-1)/B_{\ell-1}\).
The von Staudt--Clausen theorem gives \(v_\ell(B_{\ell-1})=-1\).
Thus the residual form of weight \(k'\) is cuspidal with the
same Fourier expansion. It has the stated Hecke eigenvalues
because \(p^{k'-1}\equiv p^{k-1}\bmod\ell\) for \(p\ne\ell\),
and both powers vanish modulo \(\ell\) at \(p=\ell\). Work on the
compact connected modular curve of level
\(\Gamma(h)\cap\Gamma_0(s)\) over the coefficient DVR. With
the notation of Theorem~\ref{thm:eigen240}, let
\(\mathcal B=\omega^{k'}(-C)\otimes L_\xi\), where \(C\) is the
cusp divisor. The same width computation for this squarefree
level gives the torsor defining \(L_\xi\). Its reduction is
trivial, equivariantly for the level-\(h\) changes represented
by matrices congruent to the identity modulo \(s\).

Let \(G\) be the finite group of level-\(h\) changes coming from
\(\Gamma_0(Ns)\). The \(\chi\)-idempotent exists over the DVR
because \(|G|<\ell\). With \(\pi\) a uniformizer, the obstruction
to lifting a residual section of \(\mathcal B\) lies in
\(H^1(\mathcal B)[\pi]\), by the cohomology sequence of
multiplication by \(\pi\). Also \(H^1(\mathcal B)/\pi\) embeds
equivariantly in \(H^1(\mathcal B_{\rm sp})\). The latter is
dual to \(H^0(\omega^{2-k'})\) on the special fiber by
Kodaira--Spencer and Serre duality. For \(k'>2\) it vanishes.
For \(k'=2\) it is the trivial representation, by the action
on constants on the connected fiber. If \(\chi\ne1\), it is
nontrivial on \(G\) even after reduction. Indeed the
lower-right entries realize every unit modulo \(N\) by lifting
diagonal matrices modulo \(h\) and the identity modulo \(s\),
and \(\chi\) has order prime to \(\ell\). In both cases
Nakayama's lemma gives \(H^1(\mathcal B)^\chi=0\), since this
module is finite over the DVR. Our residual eigenvector of
character \(\chi\) therefore lifts in that character space.
The torsion-free Hecke algebra argument of
Theorem~\ref{thm:eigen240} supplies a characteristic-zero
newform. The primitive conductor of \(\chi\xi\) gives exact
level \(Ns\), and the congruence at \(\ell\) gives ordinarity.

In the arithmetic application \(k=r+1\) and \(\chi\) induces
\(\eta\alpha\), with \(\eta\) the character of \(\QQ(i)\).
The omitted case is \(r=1,\alpha=\eta\). Its strict group
injects after duality into the \(\ell\)-primary part of the
fixed finite group \(K_2(\ZZ[i])\), as in the proof of
Corollary~\ref{cor:present116}. It vanishes above a fixed
prime bound, whereas the large-prime cutoff tends to infinity.
\end{proof}

The requirement of arbitrarily large weights in this residue
class occurs in the supplied source at
\srcloc{Development/AutomorphicArithmetic/AutomorphicVarieties/ModularityArgument/AutomorphyLifting/PrimeIndexConstruction/MixSeedPAlt}{55}.
The congruences here keep the auxiliary prime and the tame
level fixed as the weight increases.

\begin{lemma}[An integral class for the dual character]
\label{lem:dual241}
For a datum in Proposition~\ref{prop:level241} from the common
interval, put \(T=O(\kappa_\ell^{r+1}\alpha)\), where \(O\)
is the minimal unramified coefficient ring and \(k_\ell\) its
residue field. There exists \(d\in H^1(\QQ,T)\), unramified
at every finite prime other than \(\ell\), whose reduction
in \(H^1(\QQ,T/\ell)\) is nonzero.
\end{lemma}
\begin{proof}
Let \(S\) contain the infinite place, \(\ell\), and the prime
divisors of the conductor of \(\alpha\). Write \(G_S\) for the
Galois group unramified outside \(S\) and \(V=T\otimes_O
\operatorname{Frac}(O)\). Complex conjugation acts by
\((-1)^{r+1}(-1)^r=-1\). The global Euler characteristic
formula over \(\QQ\) gives
\(\dim H^1(G_S,V)=1+\dim H^0(G_S,V)+\dim H^2(G_S,V)\).
Here the correction one is the dimension of the minus space
at the real place. The groups \(H^i(G_S,T)\) are finitely
generated over \(O\) and give the cohomology of \(V\) upon
inverting \(\ell\). These are the usual local and global
Euler characteristic and finiteness theorems for number
fields. \cite[I, Sections~2, 4 and 5]{VII-V-V-milne98}
Equivalently the rank formula follows by taking the
Euler characteristics of \(T/\ell^nT\), dividing their
logarithmic lengths by \(n\), and letting \(n\) tend to
infinity. The long exact sequence for multiplication by
\(\ell^n\) gives the ranks in this limit.
Choose a class nonzero modulo \(\ell H^1(G_S,T)\).
Its reduction is nonzero by the same long exact sequence,
and inflation from \(G_S\) preserves that fact.

At a finite \(p\in S\) other than \(\ell\), both \(V\)
and \(V^*(1)\) have zero invariants. If \(\alpha\) is ramified
this follows on inertia. Otherwise the Frobenius eigenvalues
are \(p^{r+1}\alpha(p)\) and \(p^{-r}\alpha(p)^{-1}\),
neither of which is one. Local duality and the local Euler
characteristic give \(H^1(\QQ_p,V)=0\). On inertia the
action on \(T\) has finite order prime to \(\ell\).
Averaging lifts an invariant of \(V/T\) to an invariant
of \(V\), so \(H^1(I_p,T)\to H^1(I_p,V)\) is injective.
The integral class is thus unramified at \(p\), and outside
\(S\) this follows from its construction.
\end{proof}

\subsection{One auxiliary quadratic field and the common Hecke divisor}\label{sec:heckedivisor345}

Keep the polynomial-conductor intervals of
Theorem~\ref{thm:conductor232} and their actual fields $E_x$.
Write $N_*$ for the modulus of $E_*$ and set
$t_*=2+\lceil V^{602}\rceil$. The rank-three comparison uses an
imaginary quadratic field whose discriminant is coprime to the
level. The modulus $N_*$ is even. We first give one choice that
works for an entire sequence of intervals.

\begin{proposition}[The quadratic field for the large coefficient primes]\leavevmode\label{prop:quadratic345}
Fix a prime $s\equiv3\bmod4$ prime to $2F$, where $F$ is the
conductor of $\psi$, and put $K=\QQ(\sqrt{-s})$. Choose two
distinct primes $u_1,u_2\ge7$ split in $K$ and prime to $30Fs$.
These three choices can be kept fixed. On every sufficiently
distant interval, $\gcd(su_1u_2,N_*)=1$. Put
$P(x)=\exp(x^\varepsilon/12)$. For each $x$ in the interval,
$1\le r<t(x)$, $\ell>P(x)$ and each complementary character
$\alpha$ of $E_x$, the following assertions hold.

The order of $\alpha$ is prime to $\ell$, its local character at
$\ell$ on inertia is trivial, and $\alpha(-1)=(-1)^r$.
We have $\ell\nmid s|\Cl(K)|$ and $\ell-1\nmid 2r$. The
intersection of the kernels of $\bar\kappa_\ell$ and of the
finite Galois character $\alpha$ contains an element acting
nontrivially on $K$. In particular these data satisfy the
arithmetic hypotheses of the generic rank-three comparison in
\srcloc{Development/AutomorphicArithmetic/AutomorphicVarieties/HeightAndVanishing/Counts}{9}.
\end{proposition}
\begin{proof}
Dirichlet's theorem gives the two auxiliary primes, for example in
the progression $1\bmod 30Fs$. They split because the quadratic
character of $K$ has conductor $s$. The primes $q,d_0$ in the
modulus $N_*$ tend to infinity with the initial scale. Its other
prime factors divide $2F$, proving the assertion about coprimality.
The bound $N_*\le V^{C_{\psi,\varepsilon}}$ holds on these intervals.
Thus $P(x)>\max(N_*,2t_*+1,s|\Cl(K)|)$ uniformly for large $V$.
The character group of $E_x$ has order at most $N_*$, its
characters are unramified at $\ell$, and the parity is the
definition of the complementary component. The inequality
$0<2r<\ell-1$ gives the claimed separation of the inertia
characters after cyclotomic twisting. Finally inertia at $s$
acts trivially on both $\mu_\ell$ and $E_x$, since
$s\nmid\ell N_*$. It acts nontrivially on $K$, whose
discriminant is $-s$. This supplies the required element,
and also shows that $K\cap E_*=\QQ$.
\end{proof}

Regard each character of $E_*$ as a character modulo $N_*$.
For $1\le r<t_*$ let $H$ range over the distinct kernels of
its geometric characters of sign $(-1)^r$, regarded as subgroups
of $(\ZZ/N_*\ZZ)^\times$. The characters of kernel $H$ form
one rational orbit $\mathcal A_H$. Indeed they are the faithful
characters of the same cyclic quotient, permuted transitively by
raising their values to exponents prime to its order.
Every character in such an orbit factors through $E_*$.
Let $\mathcal P$ consist of the primes $u\ge7$ split in $K$
and prime to $30sN_*$. Fix a place above each $u$ and use
Theorem~\ref{thm:integer} at level $N_*$. Denote the absolute
value of its nonzero integer by $I_{u,H,r}$. Its diagonal
trace is denoted by $d_\alpha(u)$ and its labelled multiset of
automorphic traces by $\mathcal T^*(u)$.

\begin{theorem}[The common divisor of the good Frobenius tests]\leavevmode\label{thm:heckedivisor345}
For each $H,r$ there is a positive integer $J_{H,r}$ with
\[
 v_\ell(J_{H,r})=\min_{\substack{u\in\mathcal P\\u\ne\ell}}
                         v_\ell(I_{u,H,r}) .
\]
One finite subset of $\mathcal P$ realizes all these minima,
simultaneously for the entire interval and every prime $\ell$.
The product $J_*=\prod_{r,H}J_{H,r}$ satisfies
\[
 \log J_* \ \ll_{K,u_1,u_2}\
      [E_*:\QQ]\,t_*\, [N_*(t_*+2)]^{A'}
        \ \ll_{\psi,\varepsilon} V^{C}
\]
for a constant $C$ depending on the fixed choices and $\varepsilon$.

If $\ell\nmid J_{H,r}$ and $\alpha\in\mathcal A_H$, there is no
characteristic-zero Galois specialization whose traces are congruent
to the diagonal modulo the maximal ideal at every good split
Frobenius away from $\ell$ and belong there to the specified
automorphic lists. In particular this excludes all such
specializations of the split residual representation, including
when several labels have identical eigenvalues at one operator.
\end{theorem}
\begin{proof}
The minimum exists in the nonnegative integers, since $\mathcal P$
contains the two distinct fixed primes. For every $\ell$ it is at
most $v_\ell(I_{u_1,H,r}I_{u_2,H,r})$, since at least one of
the two primes differs from $\ell$. Thus it has finite support. Forming
the product of the corresponding prime powers gives $J_{H,r}$.
Start the finite test set with $u_1,u_2$. Only the finitely many
primes dividing $I_{u_1,H,r}I_{u_2,H,r}$ can require any further
test. For each of them add a prime $u\ne\ell$ attaining the
minimum. The union over $r,H$ is finite. The set may depend on the
interval. Its largest prime enters none of the displayed estimates.

Each $J_{H,r}$ divides the product at the two fixed primes.
Sum the logarithmic bound of Theorem~\ref{thm:integer} at these
two primes. For fixed $r$, the sum of the cardinalities of the
orbits is at most $[E_*:\QQ]$, since they partition the selected
characters. Summing up to $t_*$ proves the asserted size.

For the exclusion, choose a minimizing $u\ne\ell$. The valuation
of $I_{u,H,r}$ is zero. By construction every local difference
$b-d_\alpha(u)$ for a label $b\in\mathcal T^*(u)$ is integral
at $\ell$, and their product over the orbit and labels differs
from $I_{u,H,r}$ by an $\ell$-adic unit. Consequently each
difference is a unit. Any congruence at every split Frobenius
would give a positive valuation at this particular $u$, a
contradiction. An integral lift of the split residual representation
satisfies those congruences wherever its Frobenius trace is on
the list, by reduction of the trace.
\end{proof}

\begin{corollary}[The length inequality on the period interval]\leavevmode\label{cor:heckebudget345}
Suppose that for each contributing triple $\ell,r,H$ with
$\ell>P(x)$ for some $x$ in the interval the following
comparison holds. For every $u\in\mathcal P$, $u\ne\ell$, and
every $\alpha\in\mathcal A_H$, there is a submultiset
$T_{\ell,r,\alpha,u}\subseteq\mathcal T^*(u)$ such that, with
one embedding into $\overline{\QQ}_\ell$ for all factors at $u$,
\[
 \sum_{\alpha\in\mathcal A_H}\frac{\lambda^{N_*}_{\ell,r,\alpha}}{e_\alpha}
 \ \le\ 2\sum_{\alpha\in\mathcal A_H}
               \sum_{b\in T_{\ell,r,\alpha,u}}
                          v_\ell\bigl(b-d_\alpha(u)\bigr).
\]
Here a contributing triple has $1\le r<t(x)$ and the characters
of $H$ occur on at least one such $E_x$. Then, uniformly on the
original interval,
\[
 \sum_{\ell>P(x)}(\log\ell)\sum_{1\le r<t(x)}
             \len_{\ZZ_\ell} M_{\ell,r}(E_x)
       \ \le\ 2\log J_* .
\]
Independently of the comparison hypothesis, both auxiliary primes
and the terminal prime of Theorem~\ref{thm:criticalmarks339} can
be chosen prime to $J_*$, on the same fields and with the same
lower count of candidates.
\end{corollary}
\begin{proof}
Theorem~\ref{thm:integer} bounds the right side of the assumed
orbit inequality by $2v_\ell(I_{u,H,r})$ for every allowed $u$.
Taking the minimum bounds the orbit length by $2v_\ell(J_{H,r})$.
At any particular $x$, the characters from $E_x$ form a union of
whole rational orbits when inflated to $E_*$. Their scalar groups
are evaluated at the single modulus $N_*$. Since $\ell>P(x)$ is
prime to $[E_x:\QQ]$, Theorem~\ref{thm:levelstrict345} gives their
exact total normalized length on $E_x$. Adding the nonnegative
terms over $\ell,r,H$ counts each orbit and twist at most once.
For a positive integer $J$, the sum of $v_\ell(J)\log\ell$ over
all primes is $\log J$, which proves the bound.

There are at most $\log J_*/\log2$ distinct prime divisors of
$J_*$. This is polynomial in $V$. The proof of
Theorem~\ref{thm:markedstrict337} permits deletion of this many
primes before making any of the three choices. Its candidate
sets have sizes $U^{1-o(1)}$ or $Y^{1-o(1)}$, with both
$\log U$ and $\log Y$ growing faster than $\log V$. The same
deletion is valid after the first auxiliary choice under the
pairwise exclusions in that proof. It preserves the critical
conditions in Theorem~\ref{thm:criticalmarks339}.
\end{proof}

The single-character comparison sought here is formulated with
Galois traces in
\srcloc{Development/AutomorphicArithmetic/AutomorphicVarieties/UnitaryThreeGroups/KatoUse}{55}.
The exclusion theorem determines where the listed systems can be
congruent. The inequality in Corollary~\ref{cor:heckebudget345}
further requires that those congruences bound every strict
extension class with its full length. This implication remains
open in the present argument.


\subsection{Lower triangular orthogonal lifts and strict cohomology}\label{sec:evenroot358}

Fix a period field \(E_x\) on the intervals of Section~\ref{sec:missing}.
Write \(N_x\) for its cyclotomic modulus. Throughout this subsection
\(\ell\nmid 2N_x[E_x:\QQ]\), \(1\le r<\lfloor x^{401}\rfloor\), and
\(\alpha\) is a character of \(\Gal(E_x/\QQ)\) with
\(\alpha(c)=(-1)^r\). Put \(\Psi=\alpha\kappa_\ell^r\),
\(\vartheta=\Psi^{-1}\), and \(\eta=\chi_{-4}\). Take the unramified
coefficient ring \(O\) generated over \(\ZZ_\ell\) by the values
of \(\alpha\). Write \(k=O/\ell O\) and \(D=\operatorname{Frac}(O)/O\).
Assume that \(\bar\Psi(I_\ell)\) has order greater than two, and choose
\(g_*\in I_\ell\) with \(\lambda=\Psi(g_*)\) of residual order greater
than two. Every complementary triple with \(\ell>\exp(x^\eps/12)\)
satisfies these hypotheses on the sufficiently distant intervals.

Let \(S_0=S(\vartheta)\subset H^1(\QQ,D(\vartheta))\) be the
group unramified away from \(\ell\) and zero at \(\ell\).
Use indices \(0,1,2\). Denote the antidiagonal matrix with entries
one by \(J_3\), and put
\[
 d(g)=\operatorname{diag}(\Psi(g),1,\Psi(g)^{-1}),\qquad
 X=E_{10}-E_{21},\qquad n(z)=1+zX+\tfrac12 z^2X^2 .
\]
For a complete Noetherian local \(O\)-algebra \(A\) with residue
field \(k\), define \(\mathcal B^-(A)\) to be the continuous lower
triangular lifts \(M:G_\QQ\to\mathrm{SO}(J_3,A)\) of \(\bar d\) with
diagonal \(d\), equal to \(d\) on \(G_{\QQ_\ell}\). At each finite
\(p\ne\ell\) require \(M\) to kill \(I_p\cap\ker\alpha\cap\ker\eta\).
Matrices are in the marked basis. These conditions fix a
representative of each root cocycle, as the proof shows.

\begin{theorem}[The integral ring of lower triangular orthogonal lifts]
\label{thm:evenroot358}
If \(S_0\simeq\bigoplus_{i=1}^s O/\ell^{a_i}O\), with all
\(a_i>0\), the functor \(\mathcal B^-\) is represented by
\[
 R^- \simeq
 O[[T_1,\ldots,T_s]]/
       (\ell^{a_1}T_1,\ldots,\ell^{a_s}T_s).
\]
For \(\ell>\exp(x^\eps/12)\) it is the closed subfunctor of the
ordinary marked ring of Theorem~\ref{thm:cot116} on which the
matrix part is lower triangular and preserves \(J_3\).
The characteristic polynomial on \(G_{\QQ(i)}\) is the split
characteristic polynomial over the whole ring \(R^-\).
\end{theorem}
\begin{proof}
For a lower triangular matrix preserving \(J_3\) with determinant
one, write \(a\) for its upper left entry and \(az\) for the entry
immediately below. Pairing its first and last columns gives last
diagonal entry \(a^{-1}\), and then its determinant gives middle
entry one. Pairing its first column with its second and with itself
gives the other two lower entries \(-z\) and \(-az^2/2\).
Thus the matrix is \(n(z)\operatorname{diag}(a,1,a^{-1})\).
The identities \(X^3=0\) and \(d(g)Xd(g)^{-1}=\vartheta(g)X\)
give \(n(z)n(z')=n(z+z')\). With fixed diagonal, the group law is
exactly
\[
 z(gh)=z(g)+\vartheta(g)z(h),\qquad z(g)\in\mathfrak m_A .
\]
Since \(\vartheta(c)=1\) and two is a unit, \(z(c)=0\).
At \(p\ne\ell\) the specified inertia quotient has order prime to
\(\ell\). Vanishing on its kernel is equivalent to being a
coboundary on inertia by finite averaging, first over finite
coefficients. At \(\ell\) a coboundary has the form
\((\vartheta(g)-1)t\). Its value at \(g_*\) determines \(t\)
since \(\vartheta(g_*)-1\) is a unit. A global strict class therefore
has a unique cocycle zero at \(g_*\), and that cocycle is zero
throughout \(G_{\QQ_\ell}\).

We give the integral representing module. For any finite
\(O\)-module \(W\), cocycles with these conditions factor,
together with \(\vartheta\), through one finitely generated
profinite quotient \(\Gamma\). Indeed, take the finite ramification
set consisting of the primes over \(2N_x\ell\) and infinity.
In the group unramified outside it replace \(\ker\bar\vartheta\)
by its maximal pro-\(\ell\) quotient. The kernel removed is
characteristic in that subgroup. The pairs \((z,\vartheta)\)
send it to the identity, since the additive finite module and the
principal units are pro-\(\ell\). The finite generation argument
of Theorem~\ref{thm:present31} applies to this quotient.
For generators \(\gamma_1,\ldots,\gamma_b\), map the free profinite
group on their lifts to \(O^b\rtimes O^\times\) by sending the
\(j\)-th generator to \((e_j,\vartheta(\gamma_j))\). Write \(L(w)\)
for the first coordinate. Form \(C\) from \(O^b\) by dividing out
the submodule generated by \(L(w)\) for the relators and for the
preimages of the local subgroups on which \(z\) must vanish.
These submodules of \(O^b\) are closed. By composition, maps
\(C\to W\) are precisely the normalized cocycles with values in \(W\).
Conversely a cocycle gives such a map by its values on the generators.

For \(W=D[\ell^j]\) these normalized cocycles are exactly
\(S_0[\ell^j]\). Indeed the invariants of \(D(\vartheta)\) globally
and at \(\ell\) are zero by evaluation at \(g_*\). On the other
inertia groups the invariants are zero or divisible, since the
inertial character has order prime to \(\ell\). The multiplication
sequence for \(\ell^j\) therefore identifies the finite and
divisible local conditions. Taking the union gives
\(\operatorname{Hom}_O(C,D)=S_0\). This group is finite by
Theorem~\ref{thm:wild98} and character descent. The structure
theorem over \(O\) now gives precisely the elementary divisors
\(\ell^{a_i}\) for \(C\). Applying its representing property to
\(\mathfrak m_A/\mathfrak m_A^j\) and taking inverse limits gives
\(\operatorname{Hom}_O(C,\mathfrak m_A)\). Choosing generators
of the cyclic factors of \(C\) gives the asserted power series
presentation, including the case \(s=0\).

It remains to identify the closed subfunctor in the marked ring.
The conditions of lower triangularity and preservation of \(J_3\)
are closed matrix equations. The inverse-transpose involution
\(M\mapsto J_3 M^{-\mathsf T}J_3\) fixes every such matrix, so
the matrix part is an ordinary homomorphism on \(G_\QQ\).
Its first diagonal entry \(a\) is a character reducing to
\(\bar\Psi\). The characteristic polynomial on inertia gives
\(a(g_*)=\lambda\), since its other two possible root factors
are units there. Use the ordered product in \eqref{eq:word68} at
\((s,g_*,g_*)\), with \(s\in I_\ell\). Its first diagonal entry is
\((a(s)-\Psi(s))(\lambda-1)(\lambda-\lambda^{-1})\). Thus
\(a=\Psi\) on inertia at \(\ell\). On every other inertia group
\(a/\Psi\) factors through the prescribed prime-to-\(\ell\)
quotient and takes values in principal units, so it is trivial.
Also \(a(c)=\Psi(c)=1\). Class field theory and the trivial ideal
class group of \(\QQ\) give \(a=\Psi\) globally. This argument
works on all finite quotients of \(A\) and hence on \(A\).

In the resulting expression \(M(g)=n(z(g))d(g)\), the lower
adjacent mark is \((\lambda+1)z(g_*)=0\), so \(M(g_*)=d(g_*)\).
Use \eqref{eq:word68} now at \((g_*,s,g_*)\). Its \(10\)-entry
is \((1-\lambda)(\lambda-\lambda^{-1})M(s)_{10}\).
Both factors are units, giving \(z(s)=0\). For
\(g\in G_{\QQ_\ell}\), evaluating the cocycle on
\(g_*g=g(g^{-1}g_*g)\) gives \(z(g)=0\). In the other direction
every lift in \(\mathcal B^-(A)\) has exactly the split local
matrix at \(\ell\), which lies on the specified component by
Lemma~\ref{lem:component68}. It obeys the inertia kernels and
the marks, including \(M(c)=1\). Pairing it with \(\eta\) gives
the required marked polarized lift. Finally triangularity gives
the characteristic polynomial directly over \(R^-\).
\end{proof}

\begin{corollary}[Orders detected by lifting through finite coefficients]
\label{cor:rootcounts358}
Let \(A_j=O/\ell^{j+1}O\). For integers \(n\ge1,\ m\ge0\),
\[
 \#\operatorname{im}\bigl(\mathcal B^-(A_{n+m})\to
                              \mathcal B^-(A_n)\bigr)
 = |k|^{\sum_i\min(n,\max(a_i-m,0))}
 = \#\bigl((\ell^m S_0)[\ell^n]\bigr).
\]
For a character summand retained in \(U_{\ell,r}\) of
Section~\ref{sec:missing}, take \(m=b_0\) when \(\ell\le P\)
and \(m=0\) when \(\ell>P\). Under the hypotheses of this
subsection, the stabilized logarithm on the left is the full
contribution of that residue character orbit to the logarithmic
order of \(R_{\ell,r}\).
\end{corollary}
\begin{proof}
Each \(T_i\) maps to the maximal ideal of \(A_{n+m}\) and is
killed by \(\ell^{a_i}\). Its possible reductions in \(A_n\)
form the image of the ideal
\(\ell^{\max(1,n+m+1-a_i)}O\). Counting residues gives the
displayed exponent, independently in every coordinate. Prime-to-degree
descent for the character orbit identifies its strict group with
the additive group of \(S_0\) over the minimal ring \(O\), as in
Theorem~\ref{thm:levelstrict345}. The averaging idempotent in
Section~\ref{sec:missing} is zero on each retained orbit.
Here \(\ell\nmid[E_x:\QQ]\) also ensures that the scalar local
conditions at the conductor agree by the inertia calculation in
Section~\ref{sec:levelstrict345}. The quotient by the subgroup killed by \(\ell^m\) is isomorphic
to the image of multiplication by \(\ell^m\). Letting \(n\)
exceed every \(a_i\) proves the assertion, with the residue degree
already included in \(\log |k|\).
\end{proof}

In particular the special fiber of \(R^-\) is the formal power
series ring in \(s\) variables over \(k\). Every map into an
\(\ell\)-torsion-free coefficient ring sends all \(T_i\) to zero.
Thus even all split Frobenius characteristic polynomials on this
closed family are independent of the integers \(a_i\). A bound for
the orders requires arithmetic information on the integral
relations or on their interaction with the other deformations.


\subsection{Parity of orthogonal points after CM base change}\label{sec:paritydescent358}

Return to the large primes of Section~\ref{sec:ordinaryorth122}.
Write \(H=G_{\QQ(i)}\). For a characteristic-zero marked point,
let \(U\) be its representation on \(H\). In the orthogonal case
Lemma~\ref{lem:spin122} gives a two-dimensional representation
\(\rho\) of \(G_\QQ\) with \(U\simeq(\operatorname{Ad}^0\rho)|_H\).
At complex conjugation \(c\), its determinant is \(1\) or \(-1\).
These are called the even and odd cases, respectively.

\begin{theorem}[Potential automorphy on the orthogonal locus]
\label{thm:paritydescent358}
Let \(U\) be absolutely irreducible with orthogonal monodromy
at a characteristic-zero marked point. The following are equivalent.
\begin{enumerate}
\item There is a CM field \(F\supset\QQ(i)\) over which
      \((U\otimes\kappa_\ell^{-1})|_{G_F}\) is the representation
      associated to a regular algebraic conjugate self-dual
      cuspidal representation of \(\mathrm{GL}_3(\mathbb A_F)\),
      with the normalization of Corollary~\ref{cor:tangent119}.
\item The two-dimensional lift \(\rho\) is odd.
\end{enumerate}
In the odd case one may take \(F=\QQ(i)\).
\end{theorem}
\begin{proof}
Put \(V=\operatorname{Ad}^0\rho\), a representation over
\(\overline{\QQ}_\ell\) with determinant one and \(V\simeq V^\vee\).
Theorem~\ref{thm:orth122} gives connected monodromy
\(\mathrm{SO}_3\) on \(H\). For every finite-index subgroup of \(H\)
its image is still Zariski dense in this group. Indeed the group is a
finite union of translates of the closure of the smaller image, and
a connected algebraic group is irreducible. Thus \(V|_{G_F}\) is
absolutely irreducible for every \(F\) in the first assertion.

Suppose such a field and cuspidal representation \(\Pi\) exist.
Write \(F^+\) for its maximal totally real subfield and \(\sigma\)
for the nontrivial element of \(\Gal(F/F^+)\). In the given
normalization the unramified Satake multiset of \(\Pi_w\) equals
the eigenvalue multiset of \(V\) at geometric Frobenius, after
the fixed coefficient embedding. Indeed local reciprocity for
the associated representation introduces
\(\lvert\det\rvert_w^{-1}\), multiplying the three roots by
\(\#k(w)\), which is exactly the factor from \(\kappa_\ell^{-1}\).
Since \(V\) extends to \(G_{F^+}\), these multisets agree at
\(w\) and \(\sigma w\). They are stable under inversion and
their products are one. Strong multiplicity one, including its
rank-one case for the central character, therefore gives
\[
          \Pi^\sigma\simeq\Pi,\qquad
          \Pi^\vee\simeq\Pi,\qquad \omega_\Pi=1.
\]

By quadratic cyclic descent there is a cuspidal \(\Sigma\) on
\(\mathrm{GL}_3(\mathbb A_{F^+})\) with base change \(\Pi\).
The descents differ by twisting with the quadratic character
\(\delta\) of \(F/F^+\). In particular \(\Sigma^\vee\) is either
\(\Sigma\) or \(\Sigma\otimes\delta\). The pullback of
\(\omega_\Sigma\) under the idelic norm is \(\omega_\Pi=1\).
Embedding an idele of \(F^+\) into \(F\) gives its square under
this norm, so \(\omega_\Sigma^2=1\). The second possibility would
give \(\omega_\Sigma^{-1}=\omega_\Sigma\delta^3\), contradicting
the nontriviality of \(\delta\). Hence \(\Sigma\) is self-dual.
Here we use cyclic base change and its description of the
fibers, with the descent argument of Clozel and
Rajan.~\cite{VII-V-I-ac123,VII-V-clozelrajan358}

At a real place of \(F^+\), base change restricts the real Weil
parameter to \(\CC^\times\). Its infinitesimal exponents at the
chosen complex embedding are therefore those of \(\Pi\). They are
distinct and integral since \(\Pi\) is regular algebraic and the
half-sum of positive roots for \(\mathrm{GL}_3\) is \((1,0,-1)\).
Consequently \(\Sigma\) is regular algebraic at every real place.
Its associated semisimple Galois representation exists. After the
same Tate normalization write it as \(V_\Sigma=r_{\ell,\iota}(\Sigma)
\otimes\kappa_\ell\). Unramified local compatibility and base change
give the characteristic polynomials of \(V|_{G_F}\) at all good
places of \(F\). Chebotarev density and the characteristic-zero
semisimplification theorem identify \(V_\Sigma|_{G_F}\) with
\(V|_{G_F}\), because the latter is irreducible. In particular
\(r_{\ell,\iota}(\Sigma)\) is irreducible as well.

Identify the two representations on \(G_F\). At any
\(g\in G_{F^+}\) the two matrices implement the same conjugation
on this subgroup's image. Their quotient centralizes that image
and is scalar by Schur's lemma. These scalar ratios give a
character of the group of order two \(\Gal(F/F^+)\).
At a real complex conjugation \(c_v\) the traces of \(V_\Sigma\)
and \(V\) thus agree up to sign. Taylor's Proposition~1 applies
to \(\Sigma\) over the totally real field \(F^+\) in odd rank
three, with the irreducibility just proved. It gives
\(\lvert\operatorname{tr} r_{\ell,\iota}(\Sigma)(c_v)\rvert\le1\).
The Tate twist only changes its sign there.~\cite{VII-V-taylor358}

If \(\rho\) were even, its matrix at \(c_v\) would be scalar,
since \(c_v\) is conjugate to complex conjugation in \(G_\QQ\),
its square is one and its two eigenvalues have product one.
Conjugation by this scalar is the identity on the three-dimensional
space of trace-zero matrices. Thus \(\operatorname{tr}V(c_v)=3\),
contradicting the inequality. This proves the necessity of oddness
even when \(F/\QQ\) is not solvable.

In the odd case the proof of Theorem~\ref{thm:oddtrace123}
identifies \(\rho\) with the Galois representation of a
weight-\((r+1)\) newform with nebentypus \(\eta\alpha\).
Its adjoint lift is cuspidal and stays cuspidal over \(\QQ(i)\),
as proved there using the connected adjoint monodromy. Its
normalized unramified parameters are \(a/b,1,b/a\) for the two
Frobenius eigenvalues \(a,b\) of \(\rho\). The ordinary Tate
normalization gives exactly \(U\otimes\kappa_\ell^{-1}\).
This supplies the required realization.
\end{proof}

Theorem~\ref{thm:paritydescent358} concerns irreducible
characteristic-zero points. The lower triangular lifts of
Theorem~\ref{thm:evenroot358} give every finite layer of the
complementary strict group at the split point. Controlling its
length by congruences of irreducible systems requires information
on the ambient integral relations at that point.


\subsection{Recovering the strict group from integral Frobenius traces}\label{sec:mixedtrace359}

Use the fields, marked ordinary ring, coefficient ring \(O\), residue
field \(k\), and characters of Section~\ref{sec:evenroot358}. In this
subsection \(\ell>P=\exp(x^\eps/12)\), the interval is sufficiently
distant, and \(1\le r<\lfloor x^{401}\rfloor\). Take \(K=\QQ(i)\)
or the fixed imaginary quadratic field of
Proposition~\ref{prop:quadratic345}. In the latter choice all
Frobenius tests below are on the field of
Theorem~\ref{thm:heckedivisor345}. Write \(\eta\) for the
character of \(K\), \(H=G_K\), \(\beta=\eta\Psi\), and
\(D=\operatorname{Frac}(O)/O\). Use the same marking, with
this sign in the polarized group and this sign's kernel in the
inertia conditions. At \(\ell\) use the ordinary component through
the split point defined in Section~\ref{sec:original68}.
Let \(S\) consist of infinity and the primes dividing
\(2|\operatorname{disc} K|N_x\ell\).
Write \(G_S\) for the Galois group unramified outside \(S\).
Let \(\mathbb T\) be the closed \(O\)-algebra generated by the
traces on \(H\) of the universal marked representation. Its evaluation
at the split diagonal has kernel \(\mathfrak a\). Since \(K\)
is fixed, \(\ell\) is prime to its discriminant, class number
and number of roots of unity on the indicated intervals.

\begin{lemma}[An integral upper extension]\label{lem:upperclass359}
If \(r\ge2\), or if \(r=1\) and \(\eta\alpha\ne1\), there is an
\(O(\beta)\)-valued continuous cocycle \(b\) on \(G_S\), unramified
as a class at every finite \(p\ne\ell\), whose class modulo \(\ell\)
is nonzero. When \(r=1\), it can be chosen to come from local units
after the unramified extension at \(\ell\) killing \(\eta\alpha\).
In the remaining case \(r=1,\ \alpha=\eta\), the group
\(S_0=S(\Psi^{-1})\) is zero at all these primes on all sufficiently
distant intervals.
\end{lemma}
\begin{proof}
First take \(r\ge2\) and set \(V=\operatorname{Frac}(O)(\beta)\).
The global Euler characteristic over \(\QQ\) gives
\(\dim H^1(G_S,V)-\dim H^2(G_S,V)=1\), since
\(\beta(c)=-1\) and \(H^0(G_S,V)=0\). This is the rational form
of the finite-coefficient Euler characteristic used in
Section~\ref{sec:chart69}, obtained by applying it to the quotients
of the lattice by \(\ell^j\) and taking ranks. Integral cohomology
is finitely generated and rationalizes to this cohomology. Its
degree-one group is torsion-free because
\(H^0(G_S,k(\bar\beta))=0\). Choose a primitive element in this
positive-rank \(O\)-module. Its reduction is nonzero by the
multiplication sequence for \(\ell\).

At \(p\ne\ell\), pass to a finite local extension killing the
finite part of \(\beta\). On inertia there a cocycle becomes
a homomorphism to \(O\). Its only possible contribution is from
the pro-\(\ell\) tame inertia quotient. For residue degree \(f\),
arithmetic Frobenius acts there by \(p^f\) and on the coefficients
by \(p^{fr}\). These scalars differ, so the homomorphism is zero.
One may take the local extension of degree prime to \(\ell\),
since the character has prime-to-\(\ell\) order. Restriction on
inertial cohomology is injective by corestriction, proving the
required unramified condition.

If \(r=1\), put \(\gamma=\eta\alpha\). It is even. For nontrivial
\(\gamma\), its character field in the sense of the fixed field
of its kernel is a totally real abelian field \(L\) of degree
prime to \(\ell\). Dirichlet's unit theorem identifies the real
logarithms of its units, after tensoring with \(\RR\), with the
augmentation subspace of the regular permutation representation.
Hence the \(\gamma^{-1}\)-part of
\(\mathcal O_L^\times\otimes_{\ZZ}O\) has rank one. The character
multiplicity follows over \(\operatorname{Frac}(O)\) by extending
scalars in the rational representation. This part is saturated
in the same summand of the \(S\)-units, since the quotient of
units of these two kinds embeds by valuations in a free divisor
group. Kummer theory identifies the completed \(S\)-units with
\(H^1(G_{L,S},\ZZ_\ell(1))\). In the inverse limit the class
group contributes zero, since it is finite and its transition
maps on torsion are multiplication by \(\ell\). Twisting and
descending by the integral character idempotent gives a
primitive global class with all the stated local unit
conditions. Units at \(p\ne\ell\) give unramified Kummer classes,
since extracting roots of \(\ell\)-power order from a unit
is finite étale over the local integers. For \(\gamma=1\),
\(S_0\) is a character summand of the strict group over \(K\) at the fixed
twist \(r=1\). Proposition~\ref{prop:finite15} gives finite
prime support for its strict groups, and \(P\) eventually
exceeds every prime in this support.
\end{proof}

\begin{lemma}[Global diagonal corrections]\label{lem:diagonalpatch359}
For every finite \(O\)-module \(W\), give \(W(\eta)\) the sign
action. In degree two the localization map for \(G_S\) is
injective. In degree one localization is surjective onto
\[
 \bigoplus_{p\in S,\ p<\infty}
 H^1(\QQ_p,W(\eta))/H^1_{\rm ur}(\QQ_p,W(\eta)),
\]
where the unramified group is the kernel of restriction to inertia.
\end{lemma}
\begin{proof}
Decompose the underlying finite abelian group of \(W\) into
cyclic groups of order \(\ell^j\). The Tate dual of one sign
summand is \(\mu_{\ell^j}(\eta)\). Consider a global dual class
orthogonal to the unramified subgroups at the places in \(S\).
Restriction to \(K\) is injective by corestriction.
It gives a Kummer class \(u\) in
\(K^\times/(K^\times)^{\ell^j}\) with \(c(u)=u^{-1}\).
At every nonsplit rational prime the unique prime above it
therefore gives \(2v(u)=0\bmod\ell^j\). At a split prime in
\(S\), local reciprocity pairs an unramified character with
the valuation of a Kummer element. Orthogonality gives
\(v(u)=0\bmod\ell^j\) there. At split primes outside \(S\)
the Kummer class is unramified and the residue characteristic
differs from \(\ell\), giving the same divisibility.
Thus \((u)=\mathfrak b^{\ell^j}\) for a fractional ideal
\(\mathfrak b\). Since \(\ell\nmid |\Cl(K)|\), this ideal is
principal, so \(u\) is an \(\ell^j\)-th power times a unit.
The units of \(K\) are roots of unity of order prime to \(\ell\). The Kummer
class is therefore zero. Poitou--Tate duality now gives the two
assertions. In degree two the kernel is dual to the everywhere
locally trivial subgroup of this dual group. In degree one
the orthogonal complement of global restrictions in the sum
of the local groups is the image of global dual classes.
Those annihilating all the unramified subgroups have just been
shown to vanish. Modified cohomology at infinity is zero since
\(\ell\) is odd.~\cite{VII-V-V-milne98}
\end{proof}

\begin{theorem}[The full complementary group in the trace cotangent module]
\label{thm:mixedtrace359}
Put \(C=\operatorname{Hom}_O(S_0,D)\). For every triple in
the indicated large-prime range there is a surjective augmented
\(O\)-algebra map
\[
          \mathbb T\longrightarrow O\oplus C,\qquad C^2=0 .
\]
In particular \(\mathfrak a/\mathfrak a^2\) surjects onto \(C\),
with all its elementary divisors. Let \(\mathfrak d\) be the
ideal of \(\mathbb T\) generated by the differences of the traces
of \(h\) and \(h^{-1}\), for \(h\in H\). The map already sends
\(\mathfrak d/\mathfrak a\mathfrak d\) onto \(C\).
Its image is generated by these trace differences at good
split rational primes, with either prescribed place over each.
\end{theorem}
\begin{proof}
For \(C=0\) take split evaluation. Otherwise use the universal
normalized cocycle \(z:G_S\to C(\Psi^{-1})\) in the construction
of Theorem~\ref{thm:evenroot358}. Its representing module is
canonically \(C\) by finite double duality. In particular
\(z(g_*)=0\), \(z|_{G_{\QQ_\ell}}=0\), and its values generate
\(C\) by the generator presentation there. This construction
depends only on \(\Psi\) and its strict local conditions. At a
prime \(p\ne\ell\) the normalized cocycle vanishes on the kernel
of \(\alpha\) on inertia, so changing \(K\) imposes no new
condition on \(z\). Use \(b\) from
Lemma~\ref{lem:upperclass359}, subtracting a coboundary to
arrange \(b(c)=0\).

Form the local coefficient algebra \(A=(O\oplus C)[y]/(y^2)\),
where \(C^2=0\). This is complete and finite over \(O\).
For the two root matrices and their commutator use
\[
 X=E_{10}-E_{21},\quad Y=E_{01}+E_{12},\quad
 L=[X,Y]=\operatorname{diag}(-1,2,-1).
\]
The differential of the sign involution is
\(Z\mapsto-J_3 Z^{\mathsf T}J_3\). Combined with \(d(g)\), it
acts on \(X,Y,L\) by \(\Psi^{-1},\beta,\eta\), respectively.
Seek the matrix part of the polarized representation in the form
\[
 M(g)=\exp\bigl(z(g)X+y b(g)Y+y w(g)L\bigr)d(g),\qquad
                           w(g)\in C .
\]
The augmentation ideal \(C+Oy+yC\) has cube zero, so
\(\exp T=1+T+T^2/2\). Multiplication of the two polynomials
for \(g,h\) shows that the required correction satisfies
\[
 w(gh)-w(g)-\eta(g)w(h)
 =\tfrac12\bigl(z(g)\beta(g)b(h)-b(g)\Psi(g)^{-1}z(h)\bigr).
\]
The right side is a continuous two-cocycle with sign coefficients.
At \(\ell\) it is zero. At each \(p\in S\) away from \(\ell\)
both cocycles and their actions factor locally through the quotient
by \(I_p^0=I_p\cap\ker\alpha\cap\ker\eta\). Indeed a coboundary
on inertia vanishes on this kernel. The quotient has inertia
of prime-to-\(\ell\) order and procyclic residue quotient, so
its second cohomology with these coefficients vanishes.
There is consequently a local solution \(w_p\) inflated from
that quotient, and we take \(w_\ell=0\).

Lemma~\ref{lem:diagonalpatch359} gives a global solution \(w\)
and then a global one-cocycle to subtract so that \(w-w_p\)
is an unramified class at each finite \(p\in S\). In particular
its restriction to inertia is a coboundary. It vanishes on
\(I_p^0\) away from \(\ell\), and on \(I_\ell\), since the
sign character is unramified at \(\ell\). A global coboundary
in the sign line sets \(w(c)=0\) without changing any of these
restrictions. The matrices now kill the prescribed inertia
kernels and have \(M(c)=1\).

At \(\ell\) the matrices are upper triangular with the required
inertial diagonal, since \(z=0\) locally. For \(r\ge2\)
Theorem~\ref{thm:smooth69} identifies this ordinary condition
with the selected local component. The local proof applies for
either \(K\). The sign is unramified, and the inertia characters
on the Tate dual of the strictly upper entries are
\(\bar\kappa_\ell^{1-r}\) and \(\bar\kappa_\ell^{1-2r}\), both
nontrivial in this range. For \(r=1\),
Theorem~\ref{thm:first72} describes the component as the locus
where the Kummer valuation coordinates of the two adjacent
classes vanish. The adjacent cocycle here is \(y bY\).
Use the coefficient submodule \(Ay=Oy\oplus Cy\) of the
adjacent module. The diagonal correction lies in \(yC\) and
acts trivially on \(Ay\) since \(y^2=0\), so the action on the
two entries there is the original diagonal action.
After unramified restriction killing the finite characters the
two coordinates are images of unit Kummer classes. Valuation is
functorial in the coefficients by Lemma~\ref{lem:kummer72}.
The local description there and in Theorem~\ref{thm:first72}
uses separate coordinates for the split sign and the single
unramified quadratic extension for the nonsplit sign. These
are also the two possibilities for this \(K\).
It also remains detectable after this restriction, since the
extension degree is prime to \(\ell\) and it preserves the
valuation of a uniformizer. Thus both valuations vanish.

The remaining marks set the two adjacent entries \(01,10\) in
\(M(g_*)-J_3M(g_*)^{\mathsf T}J_3\) to zero. The lower mark
already holds because \(z(g_*)=0\). Conjugate the polarized
representation by \(1+yt(E_{01}-E_{12})\), with \(t\in O\).
This unipotent matrix is fixed by the sign involution and
preserves the mark at \(c\). The change in the upper mark is
\(yt(1-\lambda)(1+\lambda^{-1})\); at \(g_*\) the only other
nondiagonal term is already a multiple of \(y\). The coefficient
of \(yt\) is a unit. This imposes the upper mark while preserving
all local conditions. Traces on \(H\) are unchanged by this
conjugation.

Every trace on \(H\) has its split constant term and a possible
term in \(yC\), since the two linear root matrices have zero
diagonal. We prove that the trace differences span \(yC\).
Choose \(a\ge1\) with \(\ell^aC=0\) and let
\[
         N=H\cap\ker(\Psi\bmod\ell^a).
\]
On \(N\) the two cocycles with coefficients reduced modulo
\(\ell^a\) are homomorphisms. The values \(z(N)\) span \(C\).
Indeed, after quotienting \(C\) by their span, the cocycle
factors through the abelian quotient \(G_S/N\), where \(N\)
denotes its image in \(G_S\). Commuting any element with \(g_*\)
and using \(z(g_*)=0\) shows that the resulting cocycle is zero,
since \(\Psi(g_*)^{-1}-1\) is a unit. Likewise some \(k_0\in N\)
has \(b(k_0)\) a unit. Otherwise \(\bar b\) would factor through
the same abelian quotient. Its cohomology class there vanishes
by commuting with the central element \(g_*\) acting by
\(\bar\beta(g_*)\ne1\), contrary to its choice.

For \(h\in N\) take \(h'=h(chc^{-1})\) and
\(k'=k_0(ck_0c^{-1})^{-1}\). On \(N\), conjugation by \(c\)
fixes \(z\) and negates \(b\bmod\ell^a\). The two pairs of
values at \(h'\) and \(k'\) are thus \((2z(h),0)\) and
\((0,2b(k_0))\), in the indicated coefficients. In \(A/\ell^aA\)
the split diagonal is one on \(N\). For the commutator convention
\([g,h]=ghg^{-1}h^{-1}\), multiplying the two matrices and their
inverses therefore gives
\[
 M([h',k'])=1+4y\,z(h)b(k_0)L\pmod{\ell^a A}.
\]
The terms in \(yC\) of the two factors are central among these
unipotent matrices and cancel. Multiplying by \(M(g_*)\) and
subtracting \(\operatorname{tr}M(g_*)\) gives
\[
 \operatorname{tr}M([h',k']g_*)-\operatorname{tr}M(g_*)
        =4y\,z(h)b(k_0)(2-\lambda-\lambda^{-1})
\]
in \(A\) itself. Indeed the split constants cancel exactly
on the commutator, and reduction modulo \(\ell^a A\) is
injective on \(yC\). The final scalar \(2-\lambda-\lambda^{-1}
=-(\lambda-1)^2/\lambda\) is a unit. This proves that the
image of \(\mathbb T\) is \(O\oplus yC\).

Set \(f(h)=\operatorname{tr}M(h)-\operatorname{tr}d(h)\) on \(H\).
Before the final change of frame, the action of the marked \(c\)
on these matrices fixes \(X,d(h)\) and negates \(Y,L\) in the
exponential. It therefore acts by the substitution \(y\mapsto-y\).
The same conjugation takes the trace of \(M(h)\) to the trace of
its inverse by the polarization. Since the split trace is invariant
under inversion, this gives
\[
       f(chc^{-1})=-f(h),\qquad
       \operatorname{tr}(M(h))-\operatorname{tr}(M(h)^{-1})=2f(h).
\]
These identities persist after the change of frame. The values
of \(f\) span \(yC\) by the commutator calculation. The ideal
\(\mathfrak a\) maps into \(yC\), whose square is zero. This
proves both asserted module surjections.

Finally \(f\) is continuous and conjugacy invariant on \(H\)
with finite target. Its fibers are unions of conjugacy classes
in a single finite quotient of \(H\), by compactness. Chebotarev
over \(K\) realizes each such class at good prime ideals of
positive density, even after omitting any specified finite set.
One may take ideals of degree one, since those of higher degree
have density zero. The rational primes beneath them split in
\(K\). Choosing the other prime above one of them changes
\(f(h)\) to its negative, so every choice gives the same
\(O\)-span.
\end{proof}

For a residue orbit \([\alpha]\), the quotient in
Theorem~\ref{thm:mixedtrace359} has logarithmic order
\([k:\mathbf F_\ell]\log\ell\,\len_O S_0\). Thus the comparison
uses the entire large-prime contribution of that orbit on \(E_x\).
Every self-dual characteristic-zero system kills \(\mathfrak d\),
since its trace agrees with the trace of its contragredient at
every element. Theorem~\ref{thm:mixedtrace359} locates the
complementary length in the integral ideal of these differences.
To use the characteristic-zero congruences one still needs to
bound this quotient of the ideal, including its torsion at \(\ell\).


\subsection{The part of the strict order bounded by classical traces}\label{sec:classicaldefect360}

Take the fixed field \(K=\QQ(\sqrt{-s})\), the modulus \(N_*\)
and the split primes \(\mathcal P\) of Section~\ref{sec:heckedivisor345}.
At \(x\in[V,V^{3/2}]\) use the actual period field \(E_x\).
For each large-prime complementary triple take \(\mathbb T\),
\(\mathfrak a,\mathfrak d\), and \(C\) as in
Theorem~\ref{thm:mixedtrace359}, and fix one of the resulting maps
\(\Phi:\mathbb T\to O\oplus C\). All length estimates below hold
for every such choice. Write \(J=\operatorname{Frac}(O)\).

Use one isomorphism \(\overline{\QQ}_\ell\simeq\CC\) at this prime.
Consider the continuous geometric integral trace specializations
of \(\mathbb T\) with values in the integers of finite extensions
of \(J\). Retain those whose good Frobenius traces at degree-one
places of \(K\) agree with one fixed classical label of
Section~\ref{nav:III-00}, at level \(N_*\) and weight parameter
\(r\). Agreement includes every \(u\in\mathcal P\) away from
\(\ell\), at both places over \(u\), in the normalization of
that section. A label is a global cuspidal representation or a
newform with its indicated sign. Equalities at different primes
use this same label. Call the resulting set of distinct trace
specializations \(X\). Including the split evaluation, form
\[
 B=\operatorname{im}\left(\mathbb T\longrightarrow
                 O\times\prod_{\xi\in X}\overline{\QQ}_\ell\right).
\]
Put \(I=\operatorname{im}(\mathfrak d\to B)\) and
\(\mathfrak n=\operatorname{im}(\mathfrak a\to B)\).
For \(u\in\mathcal P\), denote the trace at its fixed place
by \(T_u\) and its value at the split point by \(d_\alpha(u)\).

\begin{proposition}[The determinant on the image of the trace differences]
\label{prop:traceorder360}
The algebra \(B\) is finite and torsion-free over \(O\). Let
\(X'\subset X\) consist of the specializations at which some
trace differs from the trace of the inverse. For \(u\ne\ell\)
in \(\mathcal P\), using \(v_\ell(\ell)=1\), one has
\[
 \len_O(I/\mathfrak n I)
 \ \le\ \len_O\bigl(I/(T_u-d_\alpha(u))I\bigr)
 \ =\sum_{\xi\in X'} v_\ell\bigl(\xi(T_u)-d_\alpha(u)\bigr).
\]
Each summand uses a distinct global label, including when two
labels have the same value at \(u\).
\end{proposition}
\begin{proof}
Two specializations assigned the same label agree on the
traces of degree-one primes outside a finite set, and hence on
all traces on \(G_K\) by Chebotarev and continuity. Such traces
topologically generate \(\mathbb T\). Thus a fixed label gives
at most one specialization, and \(X\) is finite by
Theorem~\ref{thm:traces}. The set \(X\) is stable under
\(\operatorname{Gal}(\overline{\QQ}_\ell/J)\). Indeed at rank
three the rationality theorem for regular algebraic cuspidal
representations carries a conjugate of the finite part to the
finite part of another cuspidal representation. The weights
are permuted among embeddings, and the fixed vectors for the
given compact level are preserved. Contragredience and the
automorphism of the base field commute with conjugation of
scalars. The condition of conjugate self-duality is therefore
preserved as well. The central restriction at the connected
real split center is still trivial, since the algebraic
central exponents there have the same sum. The normalization
of the two spherical operators uses integral powers of \(u\)
in rank three. In the newform case the conjugate Fourier
coefficients give another primitive eigenform of the same
weight and level and the normalization uses \(u^{-r}\).
These observations give the asserted stability for both
signs.~\cite{VII-V-clozel360,VII-V-V-newforms}

Each of the finitely many specializations takes values in a
finite local coefficient field. Its image is integral, so
\(B\) is a submodule of a finite free \(O\)-module. Hence
\(B\) is finite and torsion-free. Its generic fiber is a
finite reduced \(J\)-algebra. After scalar extension to
\(\overline{\QQ}_\ell\), its geometric coordinates are exactly
the distinguished point and the distinct elements of \(X\).
In detail, the displayed evaluation is faithful, so its
coordinates meet every field factor of the generic fiber.
Every embedding of each such factor is obtained by conjugating
one that occurs. The stability just proved includes all those
conjugates. Distinct coordinate characters then give distinct
factors of the split algebra. The distinguished point differs
from every \(\xi\) because the split trace has complex absolute
value greater than three at an allowed prime, whereas the
classical traces have absolute value at most three.

In a product of fields, an ideal is the product of the factors
on which it has a nonzero projection. It follows that \(I\),
after these scalar extensions, has precisely one
one-dimensional summand per element of \(X'\). Multiplication
by \(T_u-d_\alpha(u)\) acts there with eigenvalue
\(\xi(T_u)-d_\alpha(u)\), which is nonzero by the same absolute
value argument. Since \(I\) is free over the unramified DVR
\(O\), the length of its cokernel is the valuation of this
determinant, by Smith normal form. Finally \(T_u-d_\alpha(u)\)
belongs to \(\mathfrak n\). This proves the inequality, including
the case \(I=0\).
\end{proof}

The kernel of classical evaluation contributes the submodule
\[
 C_{\rm ker}=\Phi\bigl(\ker(\mathfrak d\longrightarrow I)\bigr)
       \ \subseteq\ C .
\]
It consists of the contributions from integral trace differences
whose evaluations vanish at all the selected classical systems,
including sums and products by arbitrary elements of the trace
algebra before evaluation.

\begin{proposition}[Independence of the mixed extension]\label{prop:canonicaldefect360}
For fixed marked data the submodule \(C_{\rm ker}\subseteq
S_0^\vee\) is independent of the choices of cocycles in the
construction of Theorem~\ref{thm:mixedtrace359}.
\end{proposition}
\begin{proof}
In its representing module \(C=S_0^\vee\) the lower cocycle is
the universal normalized cocycle. Its value uniquely represents
the identity of \(C\). At \(r\ge2\) the allowed integral upper
classes form a free \(O\)-module of rank one. Indeed, after
restriction to the field killing the finite character
\(\eta\alpha\), rational second cohomology for twist \(r\)
vanishes. With only the primes over \(\ell\) inverted the
integral second cohomology is finite by the comparison with
the finite group \(K_{2r-2}\) used in Theorem~\ref{thm:wild98}.
Adding an inverted place of norm \(q\), prime to \(\ell\),
adds a localization term in degree one of its residue field
with twist \(r-1\). Its rational cohomology vanishes since
\(q^{r-1}-1\ne0\). Corestriction gives the desired rational
vanishing over \(\QQ\). The Euler calculation and torsion
calculation in Lemma~\ref{lem:upperclass359} give the rank and
freeness. At \(r=1\) the construction uses the rank-one global
unit eigenspace there, except when \(C=0\).

Consequently two primitive upper classes differ by a unit
\(u\in O^\times\). The normalized representatives satisfy
\(b'=ub\), since a coboundary at \(c\) is multiplication of its
defining value by \(-2\), a unit. Compare their mixed corrections
\(w'\) and \(uw\) before the final changes of frame. The
two-cocycle equations have the same right side, so their
difference is a one-cocycle in \(C(\eta)\). The local construction
makes it unramified at \(\ell\). Away from \(\ell\) it kills the
inertia kernel \(I_p^0\); finite averaging on the quotient
shows that its class is unramified. Over \(K\) its restriction
is therefore an everywhere unramified homomorphism into \(C\).
It factors through \(\Cl(K)\) and is zero since
\(\ell\nmid|\Cl(K)|\). Corestriction shows that the global
class is zero as well. The normalization at \(c\) again forces
its representing coboundary to vanish. It follows that the
two trace maps on the augmentation ideal differ by multiplication
by \(u\), which preserves every \(O\)-submodule of \(C\).
Changes of frame preserve the traces. This proves the assertion.
\end{proof}

For the whole complementary sign
at the large primes put
\[
 \mathcal L_{>}(x)=
 \sum_{\ell>P}\log\ell\sum_{1\le r<t(x)}
    \len_{\ZZ_\ell}\bigl(\mathcal S_{\ell,r}(E_x)^{c=(-1)^r}\bigr).
\]

\begin{theorem}[A uniform bound after classical evaluation]
\label{thm:classicaldefect360}
Write \(b_\alpha=[O/\ell O:\mathbf F_\ell]\). On every
sufficiently distant polynomial-conductor interval, uniformly in
\(x\) on that interval, one has
\[
 0\ \le\ \mathcal L_{>}(x)-
 \sum_{\ell>P}\log\ell\sum_{1\le r<t(x)}\ \sum_{[\alpha]}
       b_\alpha\,\len_O C_{\rm ker}
 \ \le\ \log J_* \ \ll_{\psi,\eps} V^{C_1}.
\]
The innermost sum runs over residue character orbits of the
complementary sign on \(E_x\), and \(J_*\) is the single positive
integer from Theorem~\ref{thm:heckedivisor345}. In particular the
upper bound is \(o(\exp(x^\eps/2))\) uniformly on these intervals.
\end{theorem}
\begin{proof}
Theorem~\ref{thm:mixedtrace359} gives \(\Phi(\mathfrak d)=C\)
and \(\Phi(\mathfrak a\mathfrak d)=0\). Hence \(I\to C/C_{\rm ker}\)
is a well-defined surjection, formed by lifting to \(\mathfrak d\)
and applying \(\Phi\) modulo \(C_{\rm ker}\). Its kernel contains
\(\mathfrak nI\). By Proposition~\ref{prop:traceorder360} the
length of \(C/C_{\rm ker}\) is at most the sum of the indicated
classical difference valuations at every \(u\ne\ell\).

Inflate the characters to \(E_*\). For one rational orbit
\(\mathcal A_L\), indexed by its kernel \(L\) in the notation
of Section~\ref{sec:heckedivisor345}, all or none of the orbit
occurs on \(E_x\). Each residue orbit in it has \(b_\alpha\)
geometric characters because \(\ell\) is prime to the degree
of \(E_x\). For its representative the preceding inequality
uses a submultiset of \(\mathcal T^*(u)\), with distinct labels.
Conjugate that inequality by the \(b_\alpha\) local embeddings.
The valuations are preserved, while \(d_\alpha(u)\) ranges over
the split values of these geometric characters with multiplicity.
This repetition does not require the maps \(\Phi\) to be chosen
equivariantly, since only the representative's numerical upper
bound is conjugated. The enlarged labelled multiset
\(\mathcal T^*(u)\) contains all the conjugate factors.
The rational integer estimate of Theorem~\ref{thm:integer} now
gives
\[
 \sum_{[\alpha]\ \mathrm{in}\ \mathcal A_L}
           b_\alpha\,\len_O(C/C_{\rm ker})
      \le v_\ell(I_{u,L,r}).
\]
Taking the minimum over \(u\ne\ell\) bounds this by
\(v_\ell(J_{L,r})\). Add over \(\ell,r,L\) as in
Theorem~\ref{thm:heckedivisor345}. The identity
\(\sum_\ell v_\ell(J)\log\ell=\log J\) bounds the total by
\(\log J_*\). At every \(x\) the sum of
\(b_\alpha\log\ell\,\len_O C\) is exactly \(\mathcal L_{>}(x)\)
by prime-to-degree character descent and finite duality. This
proves the two inequalities. The polynomial bound for
\(\log J_*\) and \(x\ge V\) give the final estimate.
\end{proof}

The quantity \(C_{\rm ker}\) can equal \(C\) if the available
classical evaluations annihilate \(\mathfrak d\). The remaining
large-prime estimate therefore requires an upper bound for its total
length in
Section~\ref{sec:missing}. The smaller-prime tails and the
comparison for the signed integral period also enter that section.
The source takes its upper estimate on the graded model in
\srcloc{Development/AutomorphicArithmetic/AutomorphicVarieties/UnitaryThreeGroups/OneCharOf}{137}
and compares the full cohomological length with that model at
line~140. In the trace calculation above the unresolved integral
comparison is concentrated in \(C_{\rm ker}\).


\subsection{An auxiliary CM field with ordinary cuspidal representations}\label{sec:seed361}

Fix one polynomial-conductor interval and its modulus \(N_*\) as in
Theorem~\ref{thm:heckedivisor345}. Keep \(K=\QQ(\sqrt{-s})\) from
Proposition~\ref{prop:quadratic345}, with quadratic character \(\eta\).
Fix a coefficient prime \(\ell>\exp(V^\eps/12)\). For large \(V\)
we have \(\ell>sN_*\) and \(\ell>4t_*+5\).
The constructions below use one auxiliary field for every \(E_x\)
on the interval, every \(1\le r<t(x)\), and every character
\(\alpha\) of \(E_x\) with \(\alpha(-1)=(-1)^r\). We view finite
characters in \(\overline{\QQ}_\ell\), and use arithmetic Frobenius
in the classical modular form normalization.
Write \(\Psi=\alpha\kappa_\ell^r\).

\begin{proposition}[The fields for ordinary base change]\label{prop:auxfield361}
For each integer \(a\ge0\) there is a totally real abelian field
\(F^+\) such that \(L=KF^+\) is abelian and CM, is disjoint from
\(E_*(\mu_\ell)\) over \(\QQ\), and has the following properties.
The quadratic extension \(L/F^+\) is unramified at every finite
place and split at every place over \(sN_*\ell\).
The degree \([L:K]\) divides \(2\ell^{a+1}\).
At \(\ell\) the ramification index of \(L/\QQ\) is \(\ell^a\),
and \(\bar\kappa_\ell\) maps inertia of \(G_L\) onto
\(\mathbf F_\ell^\times\). There is a character \(\theta\) of
\(\QQ\) of order \(\ell\), of prime conductor \(w\equiv1\bmod sN_*\ell\),
which is trivial on \(G_L\). There exists one \(g\in G_L\) with
\(\bar\kappa_\ell(g)=-1\) and \(\bar\Psi(g)=1\) simultaneously
for every pair \(r,\alpha\) above.
\end{proposition}
\begin{proof}
By Dirichlet's theorem choose distinct primes \(v,w\) prime to
\(2sN_*\ell\), with \(v\equiv-1\bmod 8sN_*\ell\) and
\(w\equiv1\bmod sN_*\ell\). Let \(F_w\) be the unique degree-\(\ell\)
subfield of \(\QQ(\zeta_w)\), and \(F_a\) the unique degree-\(\ell^a\)
subfield of \(\QQ(\zeta_{\ell^{a+1}})\). Both are real since their
degrees are odd. Set \(F^+=\QQ(\sqrt{sv})F_wF_a\).
We have \(sv\equiv1\bmod4\), so the first factor has discriminant
\(sv\). At every \(p\mid sN_*\ell\), the number \(-v\) is a square
in \(\QQ_p\), by reduction and Hensel's lemma at odd \(p\) and by
\(-v\equiv1\bmod8\) at two. Thus adjoining \(\sqrt{sv}\) locally
already supplies \(\sqrt{-s}\). At other places the base change
of the unramified quadratic algebra of \(K\) is unramified or split.
This proves both assertions on \(L/F^+\). Complex conjugation
fixes \(F^+\) and acts nontrivially on \(K\), so \(L\) is CM with
maximal real subfield \(F^+\).

The only possible ramified primes of \(L\) are \(s,v,w,\ell\),
all prime to \(N_*\). At \(\ell\) only \(F_a\) ramifies and it
is totally ramified there. The other factors give an unramified
local extension, so the ramification index is exactly \(\ell^a\).
The field \(E_*(\mu_\ell)\) has ramification index dividing
\(\ell-1\) there and is unramified at \(s,v,w\). Its intersection
with \(L\) is therefore unramified at every finite prime, hence
trivial by Minkowski's discriminant bound. Disjointness supplies
an element of \(G_L\) acting as complex conjugation on
\(E_*(\mu_\ell)\). Its residual cyclotomic value is \(-1\), and
its value under every \(\bar\Psi\) is
\(\alpha(-1)(-1)^r=1\). The index of the local inertia subgroup in rational
inertia is therefore prime to \(\ell-1\), proving the residual
cyclotomic assertion. Multiplying the three real degrees gives
the displayed divisibility. A faithful character of
\(\operatorname{Gal}(F_w/\QQ)\) has the asserted conductor and
restriction. These choices are independent of \(x,r,\alpha\).
\end{proof}

\begin{lemma}[An ordinary newform with the prescribed character at the auxiliary prime]
\label{lem:cusp361}
Put \(\chi=\eta\alpha\) and let \(D\) be its primitive conductor.
For each \(m\ge r\) with \(m\equiv r\bmod\ell-1\), there exists
an \(\ell\)-ordinary cuspidal newform \(g\) of weight \(m+1\),
level exactly \(Dw\), and character \(\chi\theta\). Its residual
Galois semisimplification is \(1\oplus\chi\bar\kappa_\ell^r\),
where we reduce the finite character in the second summand.
\end{lemma}
\begin{proof}
The coprime conductors of \(\eta\) and \(\alpha\) give
\(D=s\,\operatorname{cond}(\alpha)\mid sN_*\). In particular \(\chi\)
is nontrivial and has order prime to \(\ell\). Put \(k=r+1\), so that \(2\le k<\ell-1\) and \(\chi(-1)=(-1)^k\).
At level \(D\) the holomorphic Eisenstein series
\(E=E_k(1,\chi)\) has constant term \(-B_{k,\chi}/(2k)\) and
coefficient \(\sum_{j\mid n}\chi(j)j^{k-1}\) at \(q^n\) for
\(n>0\). The weight-two formula applies because \(\chi\ne1\).
The constant term is integral at \(\ell\). Indeed
\(B_{k,\chi}=D^{k-1}\sum_{u=1}^D\chi(u)B_k(u/D)\), and the
Bernoulli numbers through degree \(k\) are integral at \(\ell\)
by the von Staudt--Clausen formula. All other coefficients are
integral there as well.~\cite{VII-V-V-newforms,VII-V-billerey361}

Work over a sufficiently large finite coefficient DVR at \(\ell\)
containing the character values and the roots of unity for the
cusps. The good-reduction modular curve and the \(q\)-expansion
principle make \(E\) integral at every cusp. One can check this
after adding full auxiliary level prime to \(D\ell\). On each
geometric component the expansions at a lift of infinity are
integral. A form with a vertical pole would, after multiplication
by the smallest power of the uniformizer clearing the pole,
reduce to a nonzero section with zero expansion on that
component. This contradicts the expansion principle.~\cite{VII-V-katz361}

At level \(M=Dw\) consider \(E(z)-E(wz)\). Its reduction is cuspidal.
To check all cusps, take
\(\gamma=\left(\begin{smallmatrix}b&b'\\c&c'\end{smallmatrix}\right)
\in\mathrm{SL}_2(\ZZ)\) and \(h=\gcd(w,c)\).
Choose \(\gamma'\in\mathrm{SL}_2(\ZZ)\) with first column
\((wb/h,c/h)^{\mathsf T}\). Then
\(\operatorname{diag}(w,1)\gamma
=\gamma'\left(\begin{smallmatrix}h&t\\0&w/h\end{smallmatrix}\right)\)
for some \(t\in\ZZ\). Since \(h,w\equiv1\bmod D\), the two
first columns agree modulo \(D\). Their unimodularity shows that
the full matrices agree modulo \(D\) after right translating one
by an integral upper unitriangular matrix. The constant terms
of \(E|_k\gamma\) and \(E|_k\gamma'\) are therefore equal.
The slash formula gives for the difference the constant term
\((1-(h/w)^k)a_0(E|_k\gamma)\). It vanishes in the residue
field because \(w\equiv1\bmod\ell\). The coefficient of \(q\)
in the difference is one. It is an eigenform for the operators
\(T_p\) with \(p\nmid M\), since both degeneracy maps commute
with these operators. Its eigenvalue at \(\ell\) reduces to one.

We give the lifting argument because \(\theta\) has order \(\ell\).
Write \(\Delta=(\ZZ/M\ZZ)^\times=P\times H\), with \(P\) its
\(\ell\)-Sylow subgroup. The group \(P\) acts freely by diamond
operators on the smooth compact modular curve \(X_1(M)\) in
both geometric fibers of the coefficient DVR. At a smooth
elliptic curve a stabilizing diamond has order dividing the
order of an automorphism, with prime factors at most three in
these characteristics. At a cusp its order divides the order
of an automorphism of the generalized elliptic \(n\)-gon with
\(n\mid M\). Its automorphisms preserving the group structure
and the polygon action form \(\mu_n\rtimes\{\pm1\}\).
In both cases the stabilizer has order prime to \(\ell\).
Here \(M\ge5\), so the full point of order \(M\) gives a fine
modular curve.~\cite{VII-V-katz361,VII-V-rapoport361}

Consequently \(X_1(M)\to Y=X_1(M)/P\) is a finite étale
torsor. The Hodge line \(\omega\) and the reduced cusp divisor
descend to \(Y\). Denote the cusp divisor by \(D_Y\). Descent with the
character of \(P\) prescribed by \(\theta\) gives a line
\(\mathcal L\) such that its sections tensored with a modular
line on \(Y\) are the corresponding \(\theta\)-eigenvectors
on the cover. This is descent for a line with the inverse
character used in its equivariant structure. It needs no
division by \(|P|\). The reduction of \(\mathcal L\) is
canonically trivial, with its natural \(H\)-action, since
\(\theta\) reduces to one and is trivial on \(H\).
The character \(\chi\) is trivial on \(P\). Its projector on
\(H\) is integral because \(\ell\nmid |H|\).

For \(k'=m+1\), use the cusp line
\(\mathcal F=\omega^{k'}(-D_Y)\otimes\mathcal L\).
On the geometric residue fiber the Serre dual of \(H^1(\mathcal F)\)
is \(H^0(\omega^{2-k'})\), since Kodaira--Spencer gives
\(\Omega^1_Y=\omega^2(-D_Y)\), by étale descent.
For \(k'>2\) it vanishes by positivity of \(\deg\omega\).
For \(k'=2\) it consists of the constants on the geometrically
connected curve and \(H\) acts trivially. In both cases the
\(\chi\)-part of \(H^1(\mathcal F)\) on that fiber is zero,
as \(\chi\) remains nontrivial after reduction. Proper coherent
cohomology in relative dimension one gives surjectivity in
degree one of the map to the residue fiber and identifies
\(H^1(Y,\mathcal F)/\varpi\) with
\(H^1(Y_{\rm sp},\mathcal F_{\rm sp})\), where the subscript denotes
base change to the residue field. The degree-two group
over the affine DVR vanishes. Nakayama's lemma and the integral
projector give \(H^1(Y,\mathcal F)^\chi=0\).
The multiplication sequence for the uniformizer now lifts
every residual \(\chi\)-cusp section to the prescribed
character \(\chi\theta\) in weight \(k'\).

Multiply the residual cusp form just constructed by the
\((m-r)/(\ell-1)\)-st power of the Hasse invariant. This keeps
its expansion, since the Hasse invariant has expansion one,
and gives the required weight and residual character.
The \(q\)-expansion formula for \(T_p\) preserves its residual
eigenvalues, since \(p^{k'-1}=p^{k-1}\) in the residue field for
\(p\ne\ell\), and both powers vanish there for \(p=\ell\).
The resulting eigensystem lifts to a characteristic-zero
cuspidal eigensystem in the space just described. Indeed the
commuting good Hecke operators, including \(T_\ell\), generate
a finite torsion-free algebra of endomorphisms of the integral
cusp lattice. The residual eigensystem gives a maximal ideal
of this algebra. A minimal prime below it has characteristic
zero, since the uniformizer is a nonzero divisor. On its
generic factor the system occurs in the cusp space by
faithfulness and simultaneous triangularization after scalar
extension. The central character \(\chi\theta\) is primitive
of conductor \(M\). Newform theory therefore puts the system
in exact level \(M\). Its \(\ell\)-th eigenvalue is a unit.
The good Hecke polynomials and Chebotarev give the stated
residual semisimplification.~\cite{VII-V-V-newforms,VII-V-V-deligne,VII-V-katz361}
\end{proof}

\begin{theorem}[Cuspidal ordinary adjoints on the common auxiliary field]
\label{thm:adjointseed361}
For every \(x,r,\alpha,m\) above, the same \(L\) carries a
determinant-one representation \(U_m\) of dimension three which
is absolutely irreducible on every open subgroup of \(G_L\).
Its residual semisimplification is
\(\bar\Psi\oplus1\oplus\bar\Psi^{-1}\).
On inertia away from \(\ell\) it is the direct sum of
\(\alpha,1,\alpha^{-1}\).
At each place above \(\ell\) it is ordinary crystalline with ordered
inertial characters \(\kappa_\ell^m,1,\kappa_\ell^{-m}\).
The three ordinary characters reduce, on the full decomposition
group, to \(\bar\Psi,1,\bar\Psi^{-1}\), in that order.
In the normalization of Section~\ref{nav:III-00},
\(U_m\otimes\kappa_\ell^{-1}\) is attached to a regular algebraic
conjugate self-dual cuspidal representation over \(L\).
In particular the construction includes the original weight \(m=r\).
\end{theorem}
\begin{proof}
Let \(W\) be the two-dimensional representation of the newform of
Lemma~\ref{lem:cusp361} and put
\(U_m=\operatorname{Ad}^0(W)|_{G_L}\), acting on trace-zero
endomorphisms. The rational representation \(W\) is absolutely
irreducible, with determinant \(\chi\theta\kappa_\ell^m\).
At \(\ell\) the ordinary newform theorem gives an unramified
quotient whose Frobenius eigenvalue is the unit root of
\(X^2-a_\ell(g)X+\chi\theta(\ell)\ell^m\). It reduces to one.
The other character follows from the determinant. The
representation is crystalline by good reduction of the level
in the modular Galois representation theorem. In the adjoint,
the upper root line, the upper triangular trace-zero plane and the
whole space give an ordered stable flag. The three quotients are
the upper-to-lower character ratio, one, and its inverse.
Both \(\eta\) and \(\theta\) restrict trivially to \(G_L\).
These facts prove the assertions at \(\ell\), including the
residual ordering.~\cite{VII-V-ordinary361,VII-V-VI-saito40}

At \(p\mid Dw\) the conductor exponent of \(W\) equals that
of its determinant, by local--global compatibility and the newform
conductor theorem. The determinant is ramified. Each wild
ramification group contributes at least as much to the Artin
conductor of \(W\) as to that of its determinant. A group acting
trivially on \(W\) acts trivially on the determinant, whose
codimension of invariants is at most one. The same inequality
holds for the tame codimension. Equality of the totals forces
\(\dim W^{I_p}=1\). In the Weil--Deligne parameter, finite
inertia therefore has a trivial summand, and its other character
is the ramified determinant. The monodromy operator commutes
with inertia. Its two distinct lines preclude any nonzero
nilpotent operator commuting with that action. Thus the actual
inertial characters of \(W\) are \(1,\chi\theta\), also at \(w\).
Outside \(Dw\ell\) the representation is unramified. Restrict
to \(G_L\) and take the two ratios and the middle character to
obtain the asserted inertial type.~\cite{VII-V-I-compat123,VII-V-V-newforms}

The residual characters of \(W\) have ratio of order greater than
two on inertia at \(\ell\), because \(2r<\ell-1\).
A nontrivial quadratic self-twist would have to permute the two
residual characters. It cannot fix either, and interchanging
them would force their ratio to have square one. A quadratic
character stays nontrivial on reduction at odd \(\ell\). Thus
\(W\) has no quadratic self-twist. Its projective image is infinite. Otherwise the representation
would become scalar on an open subgroup, contradicting its two
distinct Hodge--Tate weights. Over \(\overline{\QQ}_\ell\) every
positive-dimensional proper connected subgroup of
\(\mathrm{PGL}_2\) fixes a point or is a torus preserving two
points. Indeed its Lie algebra has dimension at most two and is
solvable, so it is triangularizable. Unless it is a torus its
nontrivial unipotent subgroup fixes a unique point, which must
also be fixed by its normalizer. A unique fixed point contradicts
irreducibility. In the torus case the full group must interchange
the two lines, and the matrix acting as \(1,-1\) on the two lines
then gives the excluded quadratic self-twist. It follows that
the projective image is Zariski dense in \(\mathrm{PGL}_2\) and
stays dense on every open subgroup. Its adjoint on
\(\mathfrak{sl}_2\) is irreducible, since an invariant subspace
is an ideal of the simple Lie algebra. It preserves the
symmetric form \(\operatorname{tr}(XY)\) and has determinant one.

Reduction of adjoint characteristic polynomials gives the three
stated characters using Lemma~\ref{lem:cusp361}. One can realize a
split residual matrix lattice after a finite coefficient extension.
In a stable basis with upper triangular residual \(W\), pass to an
extension of ramification index at least two and multiply the
second basis vector by its new uniformizer. The upper entry then
vanishes on reduction and the originally divisible lower entry
still vanishes. The resulting adjoint lattice gives the three
residual diagonal characters. To write a conjugate polarization
for \(U_m\), pair \(X,Y\) by
\(\operatorname{tr}(X\operatorname{Ad}(W(c))Y)\), where \(c\) is
complex conjugation. The involution \(\operatorname{Ad}(W(c))\)
is orthogonal, hence self-adjoint, for the trace form. This gives
a symmetric nonsingular form with the required conjugate invariance.

Let \(\pi_g\) be the unitary representation of the newform.
A nontrivial self-twist of \(\pi_g\) is quadratic by its central
character and would give one for \(W\) by equality of good
Frobenius polynomials. Gelbart--Jacquet therefore gives a
cuspidal adjoint \(\operatorname{Sym}^2\pi_g\otimes\omega_g^{-1}\).
It is self-dual of central character one. Base change it up a
tower of cyclic prime degrees to \(L\). All these degrees are
two or \(\ell\). Loss of cuspidality at a step of degree \(d\)
would require a self-twist of order \(d\) by cyclic base change.
Equality of central characters would give \(d\mid3\), which is
impossible. At each complex place the normalized parameters
are \((z/\bar z)^m,1,(z/\bar z)^{-m}\), so they are regular
algebraic. The representation over \(L\) is invariant under
complex conjugation and is self-dual. At good places its Satake
parameters are the two ratios of the Frobenius eigenvalues of
\(W\) and one. The rank-three algebraic normalization multiplies
them at geometric Frobenius by the residue cardinality, as in
Section~\ref{nav:III-00}. This identifies the attached
representation with \(U_m\otimes\kappa_\ell^{-1}\) by
Chebotarev.~\cite{VII-V-I-gj123,VII-V-I-ac123,VII-V-V-shin}
\end{proof}

\begin{theorem}[Strict classes and automorphic descent across the same extension]
\label{thm:auxdescent361}
Let \(\Phi,\mathbb T,\mathfrak a,C\) be as in
Theorem~\ref{thm:mixedtrace359} on \(K\).
Let \(\mathbb T_L\) be the closed \(O\)-subalgebra generated by
the traces on \(G_L\) and \(\mathfrak d_L\) its ideal generated
by \(t_h-t_{h^{-1}}\), \(h\in G_L\).
Then \(\Phi(\mathbb T_L)=O\oplus C\) and
\(\Phi(\mathfrak d_L)=C\). Good rational primes splitting
completely in \(L\) suffice to generate the latter image.

Let \(U\) be a strongly absolutely irreducible characteristic-zero
point of the marked ordinary problem on \(K\) in the original
weight \(r\). If \(U|_{G_L}\otimes\kappa_\ell^{-1}\) is attached
to a regular algebraic conjugate self-dual cuspidal representation
of \(\mathrm{GL}_3(\mathbb A_L)\), then \(U\) is counted by the
rank-three labels of Section~\ref{nav:III-00} on \(K\) at level
\(N_*\). The cuspidal descent of central character one is unique.
\end{theorem}
\begin{proof}
Take \(j\ge1\) such that \(\ell^jC=0\), and put
\(N'=G_L\cap\ker(\Psi\bmod\ell^j)\). Work on the absolute group \(G_\QQ\) by inflating the cocycles.
The quotient by \(N'\) is abelian. On \(N'\) the lower cocycle
\(z\) and the upper cocycle \(b\bmod\ell^j\) from
Theorem~\ref{thm:mixedtrace359} are homomorphisms. Assume \(C\ne0\).
The \(O\)-span of \(z(N')\) is all of \(C\). This span is stable
under the Galois action by the cocycle identity. After dividing it
out, \(z\) factors through the abelian quotient. Commuting with
the marked \(g_*\), where \(z(g_*)=0\) and \(\Psi(g_*)-1\) is
a unit, makes this quotient cocycle zero. All values of \(z\)
generate \(C\), which proves the assertion. The nonzero residual
class of \(b\) also stays nonzero on \(N'\). Otherwise it factors
through the same abelian quotient, where commuting with \(g_*\)
makes it a coboundary, since \(\eta(g_*)\Psi(g_*)-1\) is a unit.
Choose \(k_0\in N'\) with \(b(k_0)\) a unit.

Choose \(g\) in inertia of \(G_L\) at \(\ell\) with
\(\Psi(g)-1\) a unit, using Proposition~\ref{prop:auxfield361}.
In the mixed matrices, for \(h\in N'\) form
\(h'=h(chc^{-1})\) and \(k'=k_0(ck_0c^{-1})^{-1}\). Both belong
to \(N'\) because \(L/\QQ\) is Galois. Their reduced pairs
of root coordinates are \((2z(h),0)\) and \((0,2b(k_0))\).
The same matrix commutator calculation as in
Theorem~\ref{thm:mixedtrace359} gives
\[
 t_{[h',k']g}-t_g\ \longmapsto\
   4z(h)b(k_0)\bigl(2-\Psi(g)-\Psi(g)^{-1}\bigr)\quad\text{in }C.
\]
Indeed the central diagonal correction on the commutator already
has degree two, so multiplication by the matrix at \(g\) uses
only its split diagonal, whether or not \(g\) is marked.
The split constants cancel, and reduction modulo \(\ell^j\)
is injective on \(C\). The final factor is a unit. Inversion
negates the augmentation part of each evaluated trace on
\(G_L\) by Theorem~\ref{thm:mixedtrace359}. Since two is a unit,
the differences from inverse traces thus span \(C\). If \(C=0\)
the assertions follow by split evaluation. Finally apply
Chebotarev over \(L\) to the finite-valued trace correction.
Prime ideals of higher residue degree have density zero, and an
unramified rational prime beneath a degree-one prime of a
Galois field splits completely. Conjugation above that prime
preserves the correction or changes its sign, by the polarization
on \(K\). This proves the assertion for primes.

On \(G_K\), \(\det U=1\). Its determinant is unramified at
\(\ell\) by the ordinary inertial diagonal. Away from \(\ell\)
inertia in the marked problem on \(G_K\) factors through the
finite image of \(\alpha\), of order prime to \(\ell\).
The residual determinant is trivial, so the determinant on this
finite inertia is trivial. The resulting everywhere unramified
character of \(K\) has image in the principal units of the
coefficient DVR. It factors through the class group of \(K\),
of order prime to \(\ell\), and is therefore trivial. At each
good place the unitary normalization of the assumed cuspidal
representation \(\Pi_L\) has central character one by this
determinant computation and the rank-three Tate adjustment.
Equality on good primes gives equality of the Hecke characters.

Descend one cyclic step \(F/F_0\) between \(L\) and \(K\), of
prime degree \(d=2\) or \(\ell\). Its fields are stable under
complex conjugation and are CM since they are abelian over
\(\QQ\) and contain \(K\). On \(F\) the good Frobenius parameters
of the cuspidal representation are invariant under
\(\operatorname{Gal}(F/F_0)\), since they come from \(U\) on
\(K\). Strong multiplicity one and cyclic base change give
cuspidal descent. The descents differ by twists by characters
of \(\operatorname{Gal}(F/F_0)\), via class field theory.
The central character of a descent base changes to one and is
itself such a character by the norm theorem. Cubing is an
automorphism of this character group. Thus exactly one descent
\(\Pi_0\) has trivial central character. Its conjugate
contragredient is another such descent, so is isomorphic to
\(\Pi_0\). At a complex place there is no local extension and
the infinitesimal character is inherited, so the descent is regular algebraic.~\cite{VII-V-I-ac123,VII-V-clozelrajan358}

The attached Galois representation of \(\Pi_0\) restricts to
\(U|_{G_F}\otimes\kappa_\ell^{-1}\) after semisimplification
by good Frobenius compatibility. This restriction is irreducible
by the hypothesis on \(U\), hence the identification holds for
the representation itself. Over \(G_{F_0}\) the two extensions
differ by a character of \(\operatorname{Gal}(F/F_0)\). After
identifying their restrictions, Schur's lemma makes the ratio
of their matrices at each element scalar. Both determinants are
\(\kappa_\ell^{-3}\), by the unit central character, so the
character has cube one. Its order divides \(d\) as well, and
it is trivial. Repeating the argument gives \(\Pi_K\) over \(K\)
with the prescribed Galois representation.~\cite{VII-V-V-shin}

We check the level on the original field. At every place above
\(\ell\), the ordinary component is crystalline with labeled
exponents \(r,0,-r\). The local proof of
Theorem~\ref{thm:crystalline121} applies here because \(K\)
at \(\ell\) is either split or unramified quadratic, the two
cases of Theorems~\ref{thm:smooth69} and~\ref{thm:first72}.
The conjugate place has the dual conjugate representation.
The rank-three weight is Shin-regular, since the rank is odd.
Local--global compatibility at \(\ell\) therefore identifies
the Frobenius-semisimplified Weil--Deligne parameter including
its monodromy. Crystallinity gives trivial inertia and zero
monodromy, hence spherical level there. The labeled
Hodge--Tate exponents and the common Tate twist give
infinitesimal exponents \(r,0,-r\) in the unitary normalization.
Subtracting the half-sum of positive roots \((1,0,-1)\)
gives coefficient weights bounded by \(4(r+2)\) at both
embeddings. The positive real split center acts trivially.
\cite{VII-V-lgp361,VII-V-V-shin}

At a finite place \(v\nmid\ell\), the image of inertia is finite of order
prime to \(\ell\), with characters \(\alpha,1,\alpha^{-1}\), since
the residual characters have these unique finite lifts on the
permitted cyclic quotient. Monodromy is zero. The characters
on inertia extend to the whole local Weil group since \(\alpha\)
is global. After Frobenius semisimplification one diagonalizes
Frobenius on each inertial eigenspace. Thus the parameter is
a sum of characters with those inertial restrictions. Away from
\(N_*\) the representation is spherical. For \(v\mid N_*\) write
\(b=v(N_*)\) in the normalized local valuation. The characters
are trivial on \(1+\mathfrak p_v^b\) by local reciprocity, since
\(K/\QQ\) is unramified at the rational prime and the norm
preserves this congruence. Local--global compatibility and the
Langlands classification realize the component as a quotient
of a normalized principal series of such characters, in
Langlands order.~\cite{VII-V-V-shin}

For completeness let \(K_b=1+\mathfrak p_v^bM_3(\mathcal O_v)\)
with \(b\ge1\), and let \(B\) be the upper triangular group.
The principal series is generated by translates of the
\(K_b\)-fixed section supported on \(BK_b\) with value one at
the identity. On the lower unipotent chart its restriction is
the indicator of the subgroup of lower unipotents with entries
below the diagonal in \(\mathfrak p_v^b\), by triangular
elimination. Diagonal translation contracts this subgroup to
arbitrarily small open subgroups, shrinking each negative root
coordinate by a positive power of the uniformizer. Up to a
nonzero scalar the corresponding sections are their indicators
on the same chart. A smooth section is fixed by some \(K_e\)
with \(e\ge b\). By Iwasawa decomposition finitely many disjoint
sets \(BK_e u\), \(u\in\mathrm{GL}_3(\mathcal O_v)\), cover the
flag variety. In the lower unipotent trivialization of each
such chart the section is constant, since \(K_e\) is normal in
the maximal compact group. Subdivide the lower unipotent
subgroup of level \(e\) into finitely many right cosets of
one of the contracted subgroups contained in it. Their
indicators are translates of the displayed section, so they
span the prescribed section. Any nonzero quotient must
therefore retain a nonzero \(K_b\)-fixed vector. This proves
the required principal congruence level at \(v\).
Taking the tensor product of the local fixed vectors puts
\(\Pi_K\) on the stated list, with a consistent choice of the
Frobenius sign at every good place.
\end{proof}

The common element in Proposition~\ref{prop:auxfield361} shows that
\(\bar\kappa_\ell\bar\Psi^j\) is nontrivial on \(G_L\) for every
integer \(j\), including the four root characters used in the
residual separation conditions of
\srcloc{Development/AutomorphicArithmetic/AutomorphicVarieties/FrobeniusSmoothVectors/FrobeniusSplitPrimeDatum}{72}.
The auxiliary prime and its character realize the initial
character conditions in
\srcloc{Development/AutomorphicArithmetic/AutomorphicVarieties/CurveDVRSpecialization/HostFrameEx}{58}.
The ordinary adjoints above supply the representations described in
\srcloc{Development/AutomorphicArithmetic/AutomorphicVarieties/IdeleGenerators/ShapeFacts}{177}
for the data of this interval. Their orthogonal form makes every
trace difference \(t_h-t_{h^{-1}}\) zero on these individual
representations. Theorem~\ref{thm:auxdescent361} shows that the
mixed integral family over the same auxiliary field still
detects the whole module \(C\). A bound for \(C_{\rm ker}\) in
Theorem~\ref{thm:classicaldefect360} requires control of the
integral deformations beyond these adjoints. The smaller-prime
orders and the signed integral period in Section~\ref{sec:missing}
also remain to be estimated.


\subsection{The ordinary deformation ring at a cuspidal adjoint}
\label{sec:weightchart362}

Fix one of the forms \(g\) of Lemma~\ref{lem:cusp361} and the field
\(L\) of Proposition~\ref{prop:auxfield361}. Put \(F=L^+\) and
\(d=[F:\QQ]\). Choose a finite coefficient field \(J/\QQ_\ell\)
over which the two-dimensional representation \(W\) of \(g\)
and its ordinary filtration are defined. Enlarge it to contain
the images of the embeddings of the completions of \(F\) at
\(\ell\). We calculate at the characteristic-zero representation
\(U=\operatorname{Ad}^0(W)\) of Theorem~\ref{thm:adjointseed361}.
Here \(U\) on \(G_L\) is the restriction of this adjoint on
\(G_\QQ\). All claims in this subsection include \(m=r=1\).

On its three-dimensional space \(D\), choose a basis and let \(S\)
be the matrix of the symmetric trace form. Put
\(\iota(B)=S^{-1}B^{-\mathsf T}S\).
Use the group \(\mathrm{GL}(D)\rtimes\{1,-1\}\) in which the
negative sign acts by \(\iota\). Write \(\eta_F\) for the sign
of \(L/F\). The homomorphism at the point in question is
\[
        R_0(g)=(U(g),\eta_F(g)),\qquad g\in G_F .
\]
It is a homomorphism because the matrices \(U(g)\) preserve \(S\).
Let \(V=\operatorname{End}_J(D)\) with the adjoint action of
\(R_0\) on the Lie algebra of the identity component. Thus the
sign in this action operates by \(X\mapsto-S^{-1}X^{\mathsf T}S\).
Cohomology below is continuous, with the \(\ell\)-adic topology
on finite-dimensional \(J\)-spaces.

A deformation takes values in the same semidirect product over a
local Artinian \(J\)-algebra \(A\) of residue field \(J\). It
has prescribed reduction \(R_0\) and sign \(\eta_F\). Two lifts
are equivalent if conjugate by a matrix congruent to one. We impose the following
conditions. At each \(v\mid\ell\), choose one place of \(L\)
above \(v\). It is split over \(F\), so on this decomposition
group there is an ordinary three-dimensional representation.
Require a stable full flag by direct summands lifting the
ordered flag of \(U\). At a finite \(v\nmid\ell\) require the
kernel of \(R_0\) on inertia to act trivially. Thus only
finitely many finite places need to be added to the places above
\(\ell\) in a global ramification set. At the real places keep
the prescribed sign. Denote the resulting deformation functor
by \(\mathcal D\).

Its weight at \(v\mid\ell\) is the ordered triple
\((\xi_{v,1},\xi_{v,2},\xi_{v,3})\) of inertial characters on the
flag quotients. Via local reciprocity these are characters on
\(\mathcal O_{F_v}^{\times}\). We use one reciprocity convention
for the base and the lifts. Their base values are
\(\kappa_\ell^m,1,\kappa_\ell^{-m}\) on inertia. Denote by
\(\mathcal W\) the functor of all continuous lifts of these
inertial characters, independently at the chosen places.
An inertia matrix at the base has three distinct eigenvalues,
since the cyclotomic character has open image on inertia and
\(m>0\). In an ordinary lift the three corresponding eigenvalues
have unit pairwise differences in \(A\). The spectral projectors
of this matrix show that the flag reducing to the prescribed
ordered flag is unique. In particular taking its characters
defines a map on strict conjugacy classes.

\begin{lemma}[The local condition when the weights are fixed]
\label{lem:localweight362}
For \(v\mid\ell\) let \(\mathfrak b_v\) be the endomorphisms
preserving the base flag and \(\mathfrak n_v\) the strictly upper
triangular endomorphisms. Write \(L_v^0\subset H^1(F_v,V)\) for
the tangent condition with fixed inertial characters.
Then
\[
 L_v^0\subseteq H^1_f(F_v,V),\qquad
 \dim_J L_v^0=h^0(F_v,V)+3[F_v:\QQ_\ell].
\]
Across any small extension of Artinian \(J\)-algebras an ordinary
local lift can be lifted with any prescribed lift of its weights.
Here a small extension is a surjection whose kernel is killed
by the maximal ideal of its source.
\end{lemma}
\begin{proof}
Let \(\delta\) be the ratio of the upper and lower ordinary
characters of \(W|_{G_{F_v}}\). The adjoint quotients have
characters \(\delta,1,\delta^{-1}\). The upper entries have
characters \(\delta,\delta,\delta^2\) in their associated graded
space, and \(V/\mathfrak b_v\) has the inverse characters.
Write \(u\) for the unit root of the good Hecke polynomial
of \(g\) at \(\ell\). Before restriction to \(F_v\),
\(\delta/\kappa_\ell^m\) is unramified with Frobenius value
\((\chi\theta)(\ell)u^{-2}\), by the ordinary formula in
Theorem~\ref{thm:adjointseed361}. Every complex conjugate of
\(u\) has absolute value \(\ell^{m/2}\) by purity for the
newform of weight \(m+1\). Hence this unramified value stays
nontrivial after any positive power. In particular when \(m=1\)
the restriction of \(\delta\) still differs from
\(\kappa_\ell\). For \(m>1\) the two characters already differ
on inertia. For every \(m\), \(\delta^2\) has inertial
exponent \(2m\ne1\).~\cite{VII-V-VI-deligne40,VII-V-ordinary361}

The quotient \(V/\mathfrak b_v\) has no invariants, by its
negative exponents. Neither
\(\mathfrak b_v^\vee(1)\) nor \(\mathfrak n_v^\vee(1)\) has
invariants, since no upper character is \(\kappa_\ell\) and the
diagonal characters are trivial. This proves
\(H^2(F_v,\mathfrak n_v)=H^2(F_v,\mathfrak b_v)=0\) by local
duality. It also gives \(H^0(F_v,\mathfrak b_v)=H^0(F_v,V)\)
and an injection \(H^1(F_v,\mathfrak b_v)\to H^1(F_v,V)\).
One may fix the flag by a change of basis reducing to the identity.
Upper triangular tangent cocycles modulo changes of basis then
give precisely the image of this injection. The quotient
\(\mathfrak t_v=\mathfrak b_v/\mathfrak n_v\) has trivial action
and dimension three. The map to \(H^1(F_v,\mathfrak t_v)\) is
onto by the vanishing of \(H^2(F_v,\mathfrak n_v)\).
Local reciprocity gives surjectivity from there to the
homomorphisms on the unit group, since one can prescribe the
value on a uniformizer independently. The resulting space of
inertial variations has dimension \(3[F_v:\QQ_\ell]\) by the
logarithm on an open subgroup of the principal units. The local
Euler characteristic gives
\(\dim H^1(F_v,\mathfrak b_v)=h^0(F_v,V)+6[F_v:\QQ_\ell]\).
The kernel of the inertial variations therefore has the stated
dimension.~\cite{VII-V-V-milne98}

All modules in
\(0\to\mathfrak n_v\to\mathfrak b_v\to\mathfrak t_v\to0\)
are crystalline, as subquotients of the endomorphisms of a
crystalline representation. In the cyclotomic convention,
the three positive upper exponents contribute their full
dimension to the de Rham quotient by \(\mathrm{Fil}^0\) at
each embedding. The Bloch--Kato local dimension formula and
the local Euler characteristic therefore give
\[
 \dim H^1_f(F_v,\mathfrak n_v)
   =h^0(F_v,\mathfrak n_v)+3[F_v:\QQ_\ell]
   =\dim H^1(F_v,\mathfrak n_v).
\]
For a trivial coefficient, \(H^1_f\) consists of unramified
classes. Also \(H^1_f(F_v,\mathfrak b_v)\to
H^1_f(F_v,\mathfrak t_v)\) is surjective. Indeed the
two-term crystalline complex computing \(H^0,H^1_f\) is
\(D_{\rm cris}\to D_{\rm cris}\oplus(D_{\rm dR}/\mathrm{Fil}^0)\).
The short exact sequence above gives a short exact sequence of
these complexes by exactness on crystalline representations
and strictness of their de Rham filtrations. The rightmost
cokernels therefore map surjectively. A class with unramified
diagonal differs from a crystalline lift of its diagonal by
the image of a class in \(H^1(F_v,\mathfrak n_v)\), which is
crystalline by the dimension equality. Its image in
\(H^1(F_v,V)\) belongs to \(H^1_f\).~\cite{VII-V-VI-bk117}

Finally lift the flag basis across a small extension with
kernel \(I\). On its three quotients lift the diagonal
characters with the specified unit characters by local
reciprocity. The independent choices on a uniformizer give
continuous lifts of the unramified parts, since an infinitesimal
change there extends through \(\ZZ_\ell\) via the logarithm.
With this diagonal fixed the upper triangular lifting
obstruction lies in \(H^2(F_v,\mathfrak n_v)\otimes_J I\).
In detail, lift the entries continuously and form the defect of
multiplication. It is a two-cocycle in the strictly upper
entries tensored with \(I\), where the action is the base
adjoint action because \(\mathfrak m_{A'} I=0\) in the
extension \(A'\). Subtract a one-cochain with this coboundary.
This gives the required local homomorphism.
\end{proof}

Use \(L_v^0=H^1_{\rm ur}(F_v,V)\) at finite \(v\nmid\ell\)
and \(L_v^0=0\) at infinity. Let \(H^1_{\mathcal L^0}(F,V)\)
be the group of global classes satisfying these and the
conditions of Lemma~\ref{lem:localweight362}. For the dual
\(V^\vee(1)\) use their annihilators for the local Tate pairings,
denoted \((\mathcal L^0)^\perp\).

\begin{proposition}[Vanishing of both Selmer groups at the cuspidal adjoints]
\label{prop:dualselmer362}
At every adjoint of Theorem~\ref{thm:adjointseed361} one has
\[
 H^1_{\mathcal L^0}(F,V)=0,\qquad
 H^1_{(\mathcal L^0)^\perp}(F,V^\vee(1))=0 .
\]
\end{proposition}
\begin{proof}
Theorem~\ref{thm:adjointseed361} gives connected monodromy
\(\mathrm{PGL}_2\) in its three-dimensional adjoint action on
restriction to \(G_L\). The closure of the image of the
commutator of \(G_L\) contains the algebraic derived group,
by density and continuity of the algebraic commutator map.
This derived group is all of \(\mathrm{PGL}_2\).
The image on \(G_{L(\mu_{\ell^\infty})}\) is therefore Zariski
dense in this group. This gives the enormous-image condition
of Newton and Thorne. Indeed elements with three distinct
eigenvalues form a nonempty Zariski open subset and the group
acts irreducibly. Its span is the full endomorphism algebra
by Burnside's theorem. For any nonzero endomorphism \(X\) the
open subset where \(\operatorname{tr}(hX)\ne0\) consequently
meets the open subset of regular elements, and the cyclotomic
image meets their intersection. Some eigenprojector of \(h\)
then pairs nontrivially with \(X\). This applies in particular
to an element of every simple submodule in the enormous-image
test. One can enlarge \(J\) to split the eigenvalues, since
a local field has finitely many extensions of degree at most
three. Restriction of scalars back to \(J\) preserves vanishing.

The polarized regular algebraic cuspidal representation from
Theorem~\ref{thm:adjointseed361} has Galois representation
\(U|_{G_L}\otimes\kappa_\ell^{-1}\). The twist has trivial effect
on its adjoint action. The conjugate pairing on the underlying
untwisted space has matrix \(S U(c)\) for any complex
conjugation \(c\in G_F\). It is symmetric and gives exactly
the negative-transpose adjoint action defined above. This
identifies the polarized adjoint in Newton and Thorne's
Theorem~A with \(V\), by Schur's lemma on \(G_L\). Their theorem
now yields \(H^1_f(F,V)=0\). Lemma~\ref{lem:localweight362}
and the definition at the other places give the first
vanishing.~\cite{VII-V-VI-nt119}

We spell out the dimension calculation for the second.
The invariants of \(\operatorname{End}(D)\) on \(G_L\) are
scalars and complex conjugation negates scalars in \(V\).
Thus \(h^0(F,V)=0\). On the cyclotomic kernel the invariant
functionals on \(\operatorname{End}(D)\) are multiples of the
trace, by the same irreducibility and the nondegenerate trace
pairing. The cyclotomic action on this line in \(V^\vee(1)\)
over \(G_L\) is nontrivial. Hence \(h^0(F,V^\vee(1))=0\).
For a finite ramification set \(\Sigma\) containing the
archimedean places, the places above \(\ell\) and those of
ramification, the Poitou--Tate dimension formula is
\[
 \begin{split}
 h^1_{\mathcal L^0}(F,V)-h^1_{(\mathcal L^0)^\perp}(F,V^\vee(1))
 ={}&h^0(F,V)-h^0(F,V^\vee(1))\\
    &+\sum_{v\in\Sigma}\bigl(\dim L_v^0-h^0(F_v,V)\bigr).
 \end{split}
\]
This is the dimension identity obtained from the global duality
sequence and the local and global Euler characteristics. At each
finite \(v\nmid\ell\), unramified cohomology is the cokernel of
Frobenius minus one on \(V^{I_v}\). Its dimension equals that
of the kernel, namely \(h^0(F_v,V)\). The places above \(\ell\)
contribute \(3d\) in total by Lemma~\ref{lem:localweight362}.
At each real place the invariant Lie algebra consists of the
skew endomorphisms for the nonsingular symmetric pairing
\(S U(c_v)\) and has dimension three. The \(d\) real places
contribute \(-3d\). Thus the difference on the left is zero.
The first vanishing gives the second.~\cite{VII-V-V-milne98}
\end{proof}

\begin{theorem}[All ordinary deformations at the adjoint in characteristic zero]
\label{thm:weightchart362}
The map \(\mathcal D(A)\to\mathcal W(A)\) is a bijection for
every local Artinian coefficient algebra \(A\) as above.
The ordinary deformation functor has a power series ring
in \(3d\) variables over \(J\) as its complete local ring.
The subfunctor with determinant one on \(G_L\) has \(2d\)
variables and is given by \(\prod_i\xi_{v,i}=1\) for every
\(v\mid\ell\).
\end{theorem}
\begin{proof}
Consider a small extension \(A'\to A\) with kernel \(I\) and a
specified lift of the weights. At \(v\mid\ell\) use
Lemma~\ref{lem:localweight362}. Away from \(\ell\) we may
factor the local group by the given kernel on inertia. Its
remaining finite inertia has zero positive-degree cohomology
on \(J\)-spaces by averaging. Its residue quotient is procyclic
and has zero degree-two continuous cohomology on these
\(\ell\)-adic coefficients. For example one obtains the last
fact from the finite \(\ell\)-power quotients of a stable
lattice, the \(\ell\)-cohomological dimension of the procyclic
group, and the Mittag--Leffler property of finite groups.
The quotient therefore has zero degree-two obstruction and
its allowed local differences are unramified cohomology.
At a real place finite group averaging applies directly.

Lift the entries of a global representative continuously on the
group unramified outside \(\Sigma\). Such lifts of cochains
exist using a \(J\)-linear section of \(A'\to A\) and are
invertible because their reductions are. The multiplicative defect represents
a class in \(H^2(G_{F,\Sigma},V)\otimes_J I\). Local solvability
places it in the kernel of localization. Poitou--Tate identifies
this kernel with the dual of the group of everywhere locally
trivial dual classes (using localization at \(\Sigma\) for
restricted-ramification cohomology). Those classes are
contained in \(H^1_{(\mathcal L^0)^\perp}(F,V^\vee(1))=0\).
Thus there is a global homomorphism lifting the representative.

Compare this lift with the chosen admissible local lifts at
\(\Sigma\). Their differences give an element of
\(\bigoplus_{v\in\Sigma}H^1(F_v,V)/L_v^0\), tensored with \(I\).
Poitou--Tate gives surjectivity onto this sum from global
degree-one cohomology, because the dual Selmer group just
computed is zero. Adjust the global lift by a global
one-cocycle having those local differences modulo \(L_v^0\).
At an ordinary place a representative of \(L_v^0\) preserves
the flag and fixes the inertia of the diagonal, after a local
change of basis. At the other finite places its representative
factors through the allowed local quotient. This adjustment
therefore satisfies every local condition and has precisely
the specified weights. For two lifts with the same weights,
first lift a change of basis identifying their reductions.
Their difference class then belongs to
\(H^1_{\mathcal L^0}(F,V)\otimes_J I=0\).
A global strict change of basis identifies the two lifts.
Starting at \(A=J\) and using a composition series of small
extensions proves the bijection.~\cite{VII-V-V-milne98}

The ratio of a lifted character to its base value takes values
in \(1+\mathfrak m_A\). Logarithm and exponential give inverse
polynomial isomorphisms between this group and
\(\mathfrak m_A\). A continuous homomorphism from the units
of \(F_v\) into the latter kills torsion and is freely
specified by \([F_v:\QQ_\ell]\) elements of \(\mathfrak m_A\)
after choosing a \(\ZZ_\ell\)-basis of the maximal torsion-free
pro-\(\ell\) quotient. This proves the power series assertion.
If the weight product is one everywhere, the determinant on
\(G_L\) is unramified above \(\ell\). Polarization gives the
inverse product at the conjugate places. Elsewhere the
determinant has finite inertia and reduces to the trivial
character, so is unramified there as well. The determinant
then factors through the finite class group of \(L\), and
\(1+\mathfrak m_A\) is torsion-free. It is consequently one.
The converse follows by restricting the determinant to the
ordinary inertia groups. The independent weight characters
now number two per local unit generator.
\end{proof}

\begin{theorem}[The inverse-trace ideal at the cuspidal adjoint]\label{thm:traceideal362}
In the complete local ring of Theorem~\ref{thm:weightchart362}
the ideal generated by
\(\operatorname{tr}R(h)-\operatorname{tr}R(h^{-1})\),
\(h\in G_L\), equals the weight ideal defined by
\[
       \xi_{v,1}\xi_{v,3}=1,\qquad \xi_{v,2}=1
                  \quad(v\mid\ell).
\]
The quotient is a power series ring in \(d\) variables over
\(J\). In the determinant-one ring the same ideal is generated
by the \(d\) logarithmic coordinates of the middle characters.
Thus the vanishing on the symmetric weights holds over every
Artinian quotient, including all higher infinitesimal orders.
\end{theorem}
\begin{proof}
For any \(A\), apply \(\iota\) to the matrix part of a
deformation, retaining its sign. This is an automorphism of
the semidirect product and fixes \(R_0\). On the ordinary
groups the new flag is the dual flag in reverse order. Its
reduction is the original flag, since the ordered characters
at the base agree and the separating inertia matrix identifies
the three successive projectors. The dual weights are the
inverse weights in reverse order. If the two displayed equations
hold, the dual deformation has the original weights, all local
inertia conditions, and the original reduction.
Theorem~\ref{thm:weightchart362} identifies the two up to
conjugacy. Their traces on \(G_L\) are therefore equal.

Conversely, suppose that all the inverse traces agree in \(A\).
Fix \(v\mid\ell\) and an inertia element \(g_0\) separating
the three base characters. At \(g_0,g_0^2,g_0^3\) trace
equality identifies the two characteristic polynomials by
Newton's identities in characteristic zero. Their roots
\(a_i=\xi_{v,i}(g_0)\) coincide in the prescribed order with
the roots of the dual lift. Indeed both polynomials split
and differences of distinct residual roots are units, so
a lifted root is uniquely determined by its reduction.
For \(h\) in the same inertia group use the equality of traces
at \(g_0^j h\) for \(j=0,1,2\). Computing in each of the two
stable flags gives
\[
       \sum_{i=1}^3 a_i^j
       \bigl(\xi_{v,i}(h)-\xi_{v,4-i}(h)^{-1}\bigr)=0 .
\]
The Vandermonde determinant on the \(a_i\) is a unit.
Consequently each difference is zero. The middle character
has square one and reduces to one, so it is one since two is
invertible. This proves the two equations for every \(v\).

Both implications concern arbitrary Artinian \(J\)-algebras.
To see equality of the ideals themselves, apply each implication
to the quotient by one ideal and by successive powers of its
maximal ideal, and use Krull intersection. Universal traces are
well defined by the unique compatible strict classes in those
finite quotients. Finally the logarithmic description of the
weight ring in Theorem~\ref{thm:weightchart362} leaves one free
character per place, the first, with the third its inverse.
When the product is already one the middle character supplies
all the equations, yielding the stated \(d\) generators.
\end{proof}

At the split residual representation the mixed deformation of
Theorem~\ref{thm:auxdescent361} has the prescribed inertial
weights and its inverse-trace differences span \(C\).
The calculation here determines this same trace expression
in the characteristic-zero formal rings at the cuspidal
adjoints. The fixed-weight condition appearing in
\srcloc{Development/AutomorphicArithmetic/AutomorphicVarieties/ModularityArgument/LinkedUnramifiedChains/AbsGood}{30}
now has zero deformation space at these points, even with
arbitrary infinitesimal coefficients. To bound \(C_{\rm ker}\)
in Section~\ref{sec:missing} still requires an integral
comparison at the residual representation. The source's steps
relating different primes of its deformation space occur in
\srcloc{Development/AutomorphicArithmetic/AutomorphicVarieties/CurveDVRSpecialization/PropagationTheoremGL2}{32}
and at line~36 there.


\subsection{Cuspidal representations with unequal inverse traces at the prescribed weight}
\label{sec:ordinaryfamily363}

Fix the interval, \(K\), \(E_*\), \(N_*\) and the coefficient prime
\(\ell>\exp(V^\eps/12)\) of Proposition~\ref{prop:auxfield361}.
In that construction keep the auxiliary primes \(v,w\) fixed.
Write \(L_0=K\QQ(\sqrt{sv})F_w\), \(F_0=L_0^+\), and let
\(L_n=L_0F_n\), where \(F_n\) is the real degree-\(\ell^n\) field
in \(\QQ(\mu_{\ell^{n+1}})\). The index \(n\) here is a layer in this
tower. All the representations below have rank three. Fix for
the moment one of the pairs \(r,\alpha\) occurring on the interval,
and put \(\Psi=\alpha\kappa_\ell^r\).
The cuspidal adjoint \(U_0\) of
Theorem~\ref{thm:adjointseed361} in weight \(m=r\) on \(L_0\)
will be the starting point.

At each place of \(F_0\) over \(\ell\) choose one of the two
places of \(L_0\) above it. Such completions are unramified over
\(\QQ_\ell\). Take a sufficiently large finite coefficient DVR
\(O\) with fraction field of characteristic zero and residue
characteristic \(\ell\). Put \(A=O[[T]]\). For an inertia element
\(h\) at a chosen place define
\[
 u(h)=\frac{\log\langle\kappa_\ell(h)\rangle}{\log(1+\ell)},
 \qquad \nu_T(h)=(1+T)^{u(h)} .
\]
Here the brackets remove the Teichm\"uller part and the last power
is its binomial series with exponent in \(\ZZ_\ell\). The norm on
units of an unramified local extension is onto. This follows on
the residue fields from the finite-field norm, and on successive
principal unit quotients from surjectivity of the finite-field
trace, lifting successively and taking the limit. Local
reciprocity thus shows that \(u\) maps onto \(\ZZ_\ell\).
Consider the ordered inertial characters
\begin{equation}
   \xi_1=\kappa_\ell^r\nu_T^2,\qquad
   \xi_2=\nu_T,\qquad
   \xi_3=\kappa_\ell^{-r}\nu_T^{-3}.                 \label{eq:weights363}
\end{equation}
The product is one. Write \(t_n=\zeta_{\ell^n}-1\) for \(n\ge1\),
with a primitive root chosen for each \(n\), and
\(z_k=(1+\ell)^k-1\) for positive multiples \(k\) of \(\ell-1\).
At \(z_k\) the three exponents are \(r+2k,k,-r-3k\).
At \(t_n\) their algebraic exponents are \(r,0,-r\) and their
finite parts have exponents \(2,1,-3\) in the character
\(\nu_{t_n}\) of exact order \(\ell^n\). The three finite parts
are distinct because \(\ell>5\).

\begin{proposition}[A family of Hecke traces through the adjoint]
\label{prop:ordinaryfamily363}
There is a finite integral domain \(B_0\) over \(A\), containing
\(A\), with a point at \(T=0\) giving \(U_0\). For every \(t_n\)
and \(z_k\) it has a characteristic-zero point above that
parameter. Each such point \(y\) gives a semisimple
representation \(U_y\) of \(G_{L_0}\) attached, up to the fixed
twist by \(\kappa_\ell\), to a classical automorphic form on a
unitary group. For each \(h\in G_{L_0}\) there is an element
\(t(h)\in B_0\) evaluating to \(\operatorname{tr} U_y(h)\)
at all these points and at the adjoint. The residual
semisimplification at every one of them is
\(\bar\Psi\oplus1\oplus\bar\Psi^{-1}\).
\end{proposition}
\begin{proof}
Use the definite unitary group in three variables for \(L_0/F_0\)
which is quasi-split at every finite place. The existence condition
for this inner form holds in odd rank. The quadratic extension is
unramified at all finite places by Proposition~\ref{prop:auxfield361}.
Labesse's descent, in the form of Geraghty's Proposition~2.2.7,
puts the regular algebraic conjugate self-dual cuspidal
representation of Theorem~\ref{thm:adjointseed361} on this group.
At the split places the descent has the specified general linear
component. It has hyperspecial fixed vectors at the inert places,
where the original representation is unramified. Choose fixed
level away from \(\ell\) admitting these vectors and include at
one additional split place a sufficiently deep principal subgroup
to make the level sufficiently small. At that place the group
can be taken torsion-free, for instance inside the domain of
convergence of the matrix logarithm at an odd residue prime.
This removes finite stabilizers of the definite double cosets.
All tame level away from hyperspecial subgroups can thus be
confined to finitely many split places.~\cite{VII-V-geraghty363}

At \(\ell\) this descent is spherical. Its unramified
Frobenius slopes, after the algebraic normalization, are its
three ordered Hodge slopes. Indeed the two-dimensional ordinary
newform has Frobenius roots of valuations \(0,r\) over
\(\QQ_\ell\), and the adjoint has the two ratios and one.
The Tate twist shifts all slopes together. At a residue degree
\(f\) the valuations of Frobenius increase by the factor \(f\).
The formula for Iwahori Hecke operators on an unramified
principal series now gives an ordinary refinement. In that
formula the valuation of the \(j\)-th normalized operator on the
ordered refinement is the sum of the first \(j\) Frobenius slopes
minus the sum of the first \(j\) Hodge slopes, in the same
normalization. The local component is spherical and generic
since it comes from a cuspidal general linear representation.
Thus this refinement occurs in the component also when \(r=1\).
These are the spherical and Iwahori computations in Geraghty's
Lemmas~2.7.5 and~2.7.6.~\cite{VII-V-geraghty363}

We use the ordinary control theorem for this definite group.
We record the parts used here. For sufficiently small tame level
the dual ordinary module is finite free over the completed group
ring of the three principal-unit factors at each chosen place.
The prime-to-\(\ell\) diamond character can be fixed by an
integral idempotent, giving a free summand over this local ring.
At a dominant algebraic weight times a finite-order character,
its characteristic-zero specialization is dual to the space of
classical ordinary forms with that algebraic weight and diamond
character. One takes level at \(\ell\) deep enough for the
character. These assertions are Proposition~2.5.3, the
weight comparison in Section~2.6 and Lemma~2.6.4 of Geraghty.
In Proposition~2.5.3 one can take \(b_0=1\) because of our
choice of tame level. We use no localization hypothesis on the
residual representation.~\cite{VII-V-geraghty363}

Choose the prime-to-\(\ell\) summand of the adjoint.
Specialize the weight algebra to \(A\) along
\eqref{eq:weights363}. To specify this map, use the three
ordered Galois inertial characters as coordinates relative to
their values at the adjoint. At a purely algebraic weight
\(\lambda\) the \(j\)-th character of the normalized
representation, composed with local reciprocity on units, is
\(\kappa_\ell^{\,2-j}\prod_\tau\tau(\cdot)^{-\lambda_{\tau,4-j}}\),
where the cyclotomic factor is also read through reciprocity.
This is Corollary~2.7.8 of Geraghty with the twist by
\(\kappa_\ell\). The algebraic weight character of the
completed group ring uses precisely the reversed negative
\(\lambda\), with a fixed factor independent of \(\lambda\),
as in his Definition~2.6.3. Thus multiplying the three base
unit characters by \(\nu_T^2,\nu_T,\nu_T^{-3}\) defines a
continuous map on the principal-unit factors. At \(z_k\) it is
a purely algebraic weight with strictly increasing Hodge indices
and strictly dominant coefficient weight. The choice
\((\ell-1)\mid k\) keeps its prime-to-\(\ell\) diamond part
at the spherical value. At \(t_n\) it is the original dominant
algebraic weight times a finite-order character, so the control
theorem applies there as well.~\cite{VII-V-geraghty363}

The resulting module \(M\) over \(A\) is finite free and nonzero
by its adjoint specialization. Inside \(\operatorname{End}_A(M)\)
let \(B\) be the \(A\)-algebra generated by the good split
spherical Hecke operators. It is finite over \(A\) because \(A\)
is Noetherian, and it is torsion-free. On every classical fiber
these commuting operators act semisimply. Indeed the
characteristic-zero automorphic decomposition expresses the
space in terms of its irreducible constituents, on each of
which the good spherical operators act by their Satake
characters. The same holds on the dual. Every character above
one of the indicated parameters occurs on this fiber. To see
this, \(M\) is finite and faithful over \(B\), so its support
is all of \(\operatorname{Spec}B\). Localize at the prime for
the character and take the residue field. Nakayama gives a
nonzero module, also after extending that residue field.

For each eigensystem take its classical semisimple Galois
representation and twist by \(\kappa_\ell\), calling the result
\(U_y\). Geraghty's Proposition~2.7.2 gives the representation,
conjugate duality after this twist, potential semistability at
\(\ell\), and its good Frobenius polynomials. This construction
also applies to the constituents of a noncuspidal transfer.
For a good split geometric Frobenius the normalized trace
is the first spherical operator multiplied by the inverse
residue cardinality, a unit in \(O\). Pick one representation
for each good eigensystem. Its semisimplification is independent
of the choice by Chebotarev.~\cite{VII-V-geraghty363}

Evaluation at the classical characters above \(0,t_n,z_k\)
embeds \(B\) in the product of their coefficient valuation
rings. Indeed, an element evaluating to zero there acts by zero
on the corresponding fibers, by semisimplicity. Each of its
matrix entries on \(M\) is then a power series vanishing at
infinitely many \(z_k\). Any nonzero series in \(O[[T]]\) is a
power of a uniformizer times a distinguished polynomial times
a unit, by Weierstrass preparation. It has only finitely many
roots in the open unit disc. The same proof gives injectivity
if just the \(z_k\) are used. The finite \(A\)-module \(B\) is
compact with its maximal-base-ideal topology, and its image
in the product is closed. Good split Frobenius conjugacy
classes are dense in \(G_{L_0}\). Indeed the primes of \(L_0\)
of residue degree greater than one over \(F_0\) have density
zero, and one can apply Chebotarev to any finite quotient over
\(L_0\). By continuity on each finite list of classical
representations, the tuple of traces at any \(h\) is therefore
in that closed image. This defines \(t(h)\) already in \(B\).
Newton's identities give all characteristic coefficients there,
since \(2,3\) are units.

Take a minimal prime of \(B\) below the kernel of the adjoint
evaluation and pass to its quotient \(B_0\). Its contraction to
\(A\) is zero, since elements of a minimal prime of a Noetherian
ring are zero divisors and \(B\) is torsion-free over \(A\).
Lying over now supplies a point above each stated parameter,
with finite coefficient extension. A finite domain over the
complete local ring \(A\) is local, by the decomposition of a
finite algebra over a Henselian ring into local factors. We
took \(O\) large enough for the adjoint eigenvalues, so evaluation
at that point identifies this residue field with the residue
field of \(O\). All points on \(B_0\) thus have the same reduced
Frobenius polynomials. Chebotarev and Brauer--Nesbitt identify
their residual semisimplifications with that of \(U_0\).
This proves the proposition.
\end{proof}

\begin{lemma}[Ordinary flags at the finite-order parameters]
\label{lem:ordinaryflags363}
At any point \(y\) of Proposition~\ref{prop:ordinaryfamily363}
where \(U_y\) is absolutely irreducible, its local representations
at the chosen places above \(\ell\) have stable full flags with
the ordered inertial characters \eqref{eq:weights363}
specialized at \(y\).
\end{lemma}
\begin{proof}
Fix one chosen place, with decomposition group \(D\) and inertia
\(I\). Choose \(s_0\in I\) for which
\(\bar\kappa_\ell^r(s_0),1,\bar\kappa_\ell^{-r}(s_0)\)
are pairwise distinct. This is possible since rational
cyclotomic inertia is full on \(L_0\) and \(2r<\ell-1\).
Let \(a_i=\xi_i(s_0)\in A^\times\). The differences \(a_i-a_j\)
are units for \(i\ne j\). Define polynomials over \(A\) by
\[
          e_i(X)=\prod_{j\ne i}\frac{X-a_j}{a_i-a_j}.
\]
At each regular algebraic specialization \(z_k\), the ordinary
flag and its characters are given by Geraghty's
Corollary~2.7.8. The operators \(e_i(U_y(s_0))\) are the
projectors onto the three eigenlines, ordered so that the first
line and first two lines give this flag. For \(d\in D\),
\(g\in G_{L_0}\) and \(j>i\), form the trace of
\[
       e_j(U_y(s_0))\,U_y(d)\,e_i(U_y(s_0))\,U_y(g).
\]
Expanding the polynomials writes it as a combination over \(A\)
of the group traces already constructed in \(B\). It vanishes
at every \(z_k\), since \(d\) is upper triangular in the flag.
The injectivity proved in Proposition~\ref{prop:ordinaryfamily363}
makes this an identity in \(B\). The same argument gives
\(\det(X-U_y(s_0))=\prod_i(X-a_i)\) and, for \(h\in I\),
\[
             t(s_0^j h)=\sum_i a_i^j\xi_i(h)
                          \qquad (j=0,1,2)
\]
before specialization.~\cite{VII-V-geraghty363}

At the point in question the three roots still differ by units.
Work over the algebraic closure of its fraction field. By
Burnside's theorem the matrices \(U_y(g)\) span the full matrix
algebra. Its trace pairing is nondegenerate. Hence
\(e_j(U_y(s_0))U_y(d)e_i(U_y(s_0))=0\) for \(j>i\).
This gives the ordered flag on the entire decomposition group.
The three trace identities on \(I\) identify its diagonal
characters by the invertible Vandermonde matrix on the \(a_i\).
The projectors are defined over the field of the representation,
so the flag is defined there as well. At the conjugate place
conjugate duality gives the reversed inverse characters.
\end{proof}

\begin{proposition}[Irreducibility and the inertial types away from the coefficient prime]
\label{prop:awaytype363}
There is a nonzero series \(a(T)\in A\) such that at every
classical parameter of Proposition~\ref{prop:ordinaryfamily363}
where \(a\) does not vanish, all points of \(B_0\) above it are
absolutely irreducible. For such points the determinant of
\(U_y\) is one, and on inertia away from \(\ell\) one has
\[
                    U_y|_{I}\simeq\alpha\oplus1\oplus\alpha^{-1}.
\]
In particular this holds at \(t_n\) for every sufficiently large \(n\).
\end{proposition}
\begin{proof}
For any nonzero \(b\in B_0\) there is a nonzero
\(a_b\in A\cap bB_0\). Indeed \(B_0\otimes_A\operatorname{Frac}(A)\)
is a finite-dimensional domain over a field and hence is a field.
Clear the denominator of \(b^{-1}\) there. Consequently any
infinite set of distinct parameters \(t_n\), or \(z_k\), with
one point chosen over each, detects all nonzero elements of
\(B_0\), by Weierstrass preparation.

By absolute irreducibility at the adjoint, choose nine elements
\(h_i\) whose matrices there span the full matrix algebra. The
Gram determinant
\[
                 D_*=\det\big(t(h_i h_j)\big)_{1\le i,j\le9}
\]
is nonzero at that point by the trace pairing. Wherever it is
nonzero the same matrices span the full algebra, so the
representation is absolutely irreducible.

Let \(S'\) be the finite set of places of \(L_0\) away from
\(\ell\) over the chosen tame level and the conductor of \(\alpha\).
At the remaining places every \(U_y\) is unramified by
Proposition~2.7.2 of Geraghty. Fix a place in \(S'\) of residue
characteristic \(p\). Inside a fixed local closure pass to the
completion \(M'\) of \(L_0E_*\), and let \(e'\) be its degree over
the local completion of \(L_0\). Write \(q\) for the residue
cardinality of \(M'\). Residual semisimplification makes the
reduction on \(I_{M'}\) upper unipotent in a suitable stable
lattice, since \(\alpha\) and the residual cyclotomic character
are trivial there. The image of \(I_{M'}\) is pro-\(\ell\).
Indeed it lies above a finite \(\ell\)-group in a matrix group
over a coefficient DVR whose principal congruence kernel is
pro-\(\ell\). Thus wild inertia acts trivially, as do the
prime-to-\(\ell\) factors of tame inertia. If \(H\) is the
image of a generator of its tame \(\ZZ_\ell(1)\) factor, then
arithmetic Frobenius conjugates \(H\) to \(H^q\). Raising to
the \(q\)-th power permutes its at most three distinct eigenvalues.
Each eigenvalue therefore has order dividing \(q^6-1\).
Continuity gives the same bound for all inertia in \(M'\).
For inertia of the original completion the order of every
eigenvalue divides \(e'(q^6-1)\), because the \(e'\)-th power
of any such element belongs to \(I_{M'}\). The local field
extension is Galois.~\cite{VII-V-geraghty363}

For a fixed inertia element \(h\) each coefficient of its
characteristic polynomial thus takes only finitely many possible
values. Form the monic polynomial with the distinct possible
values as roots, adjoining all their conjugates. Its coefficients
lie in \(O\), since they are integral and invariant over its
fraction field. It vanishes on the corresponding function of
\(B_0\), by the preceding density argument. At the adjoint this
function takes a value \(c\in O\). Removing the simple factor
\(X-c\) leaves a factor nonzero on specialization at the
adjoint. Since \(B_0\) is a domain, the function equals \(c\)
identically. This applies separately to each \(h\), and shows
that all inertial characteristic polynomials agree with those
of the adjoint. On \(I_{M'}\) they are \((X-1)^3\).

To remove possible unipotent inertia take any arithmetic
Frobenius lift \(\phi\) over \(M'\). With \(P_\phi(X)\) its
characteristic polynomial over \(B_0\), form
\[
           b_p=\operatorname{Res}_X
                    \big(P_\phi(X),q^3 P_\phi(X/q)\big).
\]
This is nonzero at the adjoint. We verify this even at the
ramified primes of the elliptic newform. At a good prime the
Frobenius roots of the newform have equal complex absolute
values by purity. At a bad prime the two-dimensional
Weil--Deligne parameter has zero monodromy and two characters
distinct on inertia by the proof of Theorem~\ref{thm:adjointseed361}.
For the unitary local normalization the pair of quasicharacters
is stable under \(\gamma\mapsto\overline{\gamma}^{-1}\).
This follows by taking the conjugate dual of the unitary local
representation and using its Langlands parameter. That
involution fixes each finite inertial character. The two
inertial characters are distinct, hence it fixes each
quasicharacter separately. Both are unitary. The common twist
back to the classical normalization gives them equal complex
absolute values on every Frobenius lift. Passing to the local
extension raises norms to the residue degree and introduces
only roots of unity from inertia. The adjoint eigenvalues are
the two ratios and one, all of absolute value one. Multiplication
by \(q>1\) separates their multiset from itself. This proves
the asserted nonvanishing.~\cite{VII-V-VI-deligne40,VII-V-V-newforms,VII-V-I-compat123}

At any point where \(b_p\ne0\), the unipotent matrix \(H\) on
\(I_{M'}\) has a finite polynomial logarithm \(N\).
Conjugation by \(U_y(\phi)\) sends \(N\) to \(qN\). The
eigenvalues of this conjugation on endomorphisms are the ratios
of the three Frobenius eigenvalues, also for a nondiagonalizable
matrix, by triangularization. The resultant excludes \(q\)
from those ratios. Therefore \(N=0\) and \(I_{M'}\) acts
trivially. The inertia image at the original place is now
finite. Its characteristic-zero representation is semisimple
and is determined by its character. Its trace on every
element was shown to agree with the adjoint, giving
\(\alpha,1,\alpha^{-1}\).

The determinant at each classical point is everywhere
unramified. Away from \(\ell\) this follows already from the
inertial polynomials. At the chosen places above \(\ell\),
the identity that the inertial characteristic polynomial is
\(\prod_i(X-\xi_i(h))\) holds for every \(h\in I\), by
specialization at the regular \(z_k\) and injectivity as in
Lemma~\ref{lem:ordinaryflags363}. The product is one, and
conjugate duality handles the other places. The determinant
hence factors through the finite ideal class group of \(L_0\).
For each \(g\) its values at all classical points are roots
of unity of bounded order. Apply the finite-value argument
again to the determinant function of \(g\) in \(B_0\).
Its value at the adjoint is one, so it is one throughout.

Finally take \(b\) to be \(D_*\) times one resultant for every
place in \(S'\). Each factor is nonzero at the adjoint and the
product is nonzero in \(B_0\). Take \(a=a_b\). All points above
a parameter at which \(a\ne0\) have every required property.
Only finitely many of the distinct \(t_n\) can be zeros of \(a\).
\end{proof}

\begin{theorem}[Full special linear monodromy at the original weights]
\label{thm:slthree363}
For each sufficiently large \(n\), the same \(L_n\) has the
following property for every \(x,r,\alpha\) on the interval.
There is a determinant-one representation \(U\) of \(G_{L_n}\)
with Zariski closure \(\mathrm{SL}_3\) over the algebraic
closure of the coefficient field. Its residual semisimplification
is \(\bar\Psi\oplus1\oplus\bar\Psi^{-1}\). Its restriction to
inertia away from \(\ell\) is \(\alpha\oplus1\oplus\alpha^{-1}\).
At each place above \(\ell\) it is ordinary crystalline with
ordered inertial characters \(\kappa_\ell^r,1,\kappa_\ell^{-r}\).
Its ordinary characters on the full decomposition group reduce
in the order \(\bar\Psi,1,\bar\Psi^{-1}\).
The representation \(U\otimes\kappa_\ell^{-1}\) is attached to
a regular algebraic conjugate self-dual cuspidal automorphic
representation on \(\mathrm{GL}_3(\mathbb A_{L_n})\).
\end{theorem}
\begin{proof}
Choose \(y\) over \(t_n\) as in Proposition~\ref{prop:awaytype363}.
The weak base change of its definite unitary automorphic
representation to \(\mathrm{GL}_3(\mathbb A_{L_0})\) is cuspidal.
Indeed the description of the transfer used in Geraghty's
Proposition~2.7.2 expresses the Galois representation of a
transfer with several constituents as their direct sum,
and a nontrivial residual discrete-series block as a sum
of the corresponding twists. This is checked at good primes
and then by uniqueness of semisimplification. Either
decomposition would contradict \(D_*(y)\ne0\).
The cuspidal transfer is regular algebraic and conjugate
self-dual, with associated representation \(U_y\otimes
\kappa_\ell^{-1}\).~\cite{VII-V-geraghty363}

First \(U_y\) stays absolutely irreducible on every open
subgroup of \(G_{L_0}\). To see this, pass to the normal core
of any purported counterexample. On this finite-index normal
subgroup the sum of all simple submodules is nonzero and
stable under the full group, hence is the whole representation.
Its isotypic components are permuted transitively. Dimension
three leaves the scalar case or three distinct lines if that
restriction is reducible. In the scalar case the projective
image is finite, contradicting the distinct Hodge--Tate
weights, which persist on finite extensions. In the three-line
case the kernel of \(\bar\Psi\) maps into a pro-\(\ell\) subgroup
of the matrix image. Indeed the global residual representation
is triangular in a suitable lattice basis, with all diagonal
entries trivial on that kernel. The continuous permutation
action on the lines is a quotient of the matrix image. Since
\(\ell>3\), it factors through the cyclic group
\(\bar\Psi(G_{L_0})\). A transitive cyclic subgroup of \(S_3\)
is the group of order three. In a basis of the three lines the
matrices are monomial. Conjugating by a diagonal matrix with
entries \(1,\zeta_3,\zeta_3^2\), in cyclic order, therefore gives
a nontrivial cubic self-twist. Its reduction would permute
\(\{1,\bar\Psi,\bar\Psi^{-1}\}\) by multiplication by a character
of order three. The orbit of \(1\) would be the full set, whereas
\(\bar\Psi\) has order greater than three on inertia at \(\ell\)
by \(4r<\ell-1\). This rules out the last case.

The identity component of the algebraic monodromy now acts
irreducibly, since its preimage has finite index. Its unipotent
radical has fixed vectors by triangularization and normality
makes that space invariant under the component. The radical
thus acts trivially. Its connected center acts by scalars and
is trivial because the determinant is one. The resulting
connected semisimple subgroup is \(\mathrm{SL}_3\) or the
image of \(\mathrm{SL}_2\) on its symmetric square, by
Theorem~\ref{thm:orth122}. In the second case it preserves a
unique line of nondegenerate symmetric forms. Its normalizer,
and hence the entire image, acts by a character on this line.
It follows that \(U_y^\vee\simeq U_y\otimes\delta\) with
\(\delta^3=1\), taking determinants. Compare the two sides on
an ordinary inertia group using Lemma~\ref{lem:ordinaryflags363}.
The three algebraic exponents \(r,0,-r\) are distinct.
Characters with different exponents cannot agree after a
finite twist, since the cyclotomic image on inertia is open.
Consequently the middle character on the two sides must
agree, forcing \(\nu_{t_n}^{\,2}\) or its inverse to be
the restriction of \(\delta\). Its order is \(\ell^n>3\).
Thus the connected monodromy is \(\mathrm{SL}_3\), as is the
whole monodromy. Restricting to any finite extension retains
this connected Zariski closure.

It remains to check the local assertions after restriction to
\(L_n\), since the finite characters were essential over
\(L_0\). Each \(\nu_{t_n}\) is obtained from the rational
cyclotomic character via its quotient
\((1+\ell\ZZ_\ell)/(1+\ell^{n+1}\ZZ_\ell)\). This quotient
is killed by restriction to \(F_n\). The inertial characters
on the flag at \(L_n\) are therefore precisely those stated,
including at the conjugate places by conjugate duality.
Each \(U_y\) is potentially semistable at \(\ell\).
In its Weil--Deligne representation the semisimple inertia
characters are the finite parts of the three de Rham
characters on the flag. Indeed formation of the potentially
semistable module is exact on potentially semistable
representations, and a crystalline cyclotomic power has trivial
Weil--Deligne inertia. The unramified part of a flag character
also has trivial inertia. The three finite characters at \(L_0\)
are distinct, so the nilpotent monodromy operator commuting
with inertia preserves three lines and vanishes. After passage
to \(L_n\), Weil--Deligne inertia and monodromy are both
trivial. The potentially semistable representation is therefore
crystalline there, by descent of its period module. Residual
cyclotomic inertia on \(L_n\) still has order \(\ell-1\) by
Proposition~\ref{prop:auxfield361}. Its three distinct residual
characters on inertia identify the flag characters uniquely
inside the global residual triple, also on the full
decomposition group. The assertion away from \(\ell\) follows by
restriction from Proposition~\ref{prop:awaytype363}.

Cyclic base change to \(L_n\) preserves cuspidality at every
degree-\(\ell\) step. Otherwise there would be a self-twist
of order \(\ell\), which is impossible in rank three by
comparison of central characters. Regularity and conjugate
self-duality persist, and the attached representation agrees
with \(U_y|_{G_{L_n}}\otimes\kappa_\ell^{-1}\) by the good
Frobenius polynomials.~\cite{VII-V-I-ac123,VII-V-V-shin}
Finally the interval uses finitely many characters of \(E_*\)
and \(r<t_*\). Take \(n\) above the largest of the finitely
many exceptional indices from Proposition~\ref{prop:awaytype363}.
The original \(v,w,F_w\), and now \(L_n\), are common to all
these choices.
\end{proof}

\begin{corollary}[An integral inverse trace on the common auxiliary tower]
\label{cor:skewsize363}
For fixed interval, \(\ell,r,\alpha\) and the component \(B_0\)
above, there exist
\(h\in[G_{L_0},G_{L_0}]\), a nonnegative rational number \(\mu\)
and a nonnegative integer \(j\) such that for every sufficiently
large \(n\) and every \(y\) above \(t_n\),
\[
  0<v_\ell\!\left(\operatorname{tr}U_y(h)-
                   \operatorname{tr}U_y(h^{-1})\right)
     \le \mu+\frac{j}{\ell^{n-1}(\ell-1)} ,
                  \qquad v_\ell(\ell)=1 .
\]
The element \(h\) lies in every \(G_{L_n}\). At each such \(n\)
the same valuation at a given \(y\) occurs at good Frobenius
elements of rational primes splitting completely in \(L_n\).
\end{corollary}
\begin{proof}
Choose one point with full special linear monodromy from
Theorem~\ref{thm:slthree363}. The image of the commutator
subgroup of \(G_{L_0}\) there is Zariski dense in
\(\mathrm{SL}_3\). Indeed taking its closure gives an
algebraic subgroup containing the algebraic commutators of
the dense image and hence of the whole group, which is perfect.
The trace and the inverse trace are distinct regular functions
there, as is seen on \(\operatorname{diag}(2,2,1/4)\). Choose
an element \(h\) of that commutator subgroup at which they
differ. Every \(L_n/L_0\) is abelian, so \(h\in G_{L_n}\)
for every \(n\). Now \(b=t(h)-t(h^{-1})\) is nonzero in
\(B_0\). Choose \(0\ne a\in (b)\cap A\) as in the proof of
Proposition~\ref{prop:awaytype363}. At every integral
specialization \(v_\ell(b(y))\le v_\ell(a(t_n))\) when
\(a(t_n)\ne0\).

Factor \(a=\varpi^e P(T)u(T)\) by Weierstrass preparation,
with \(P\) distinguished of degree \(j\) and \(u\) a unit.
Set \(\mu=e\,v_\ell(\varpi)\). Each root of \(P\) has strictly
positive valuation, counting a zero root as valuation infinity.
Since \(v_\ell(t_n)=1/(\ell^{n-1}(\ell-1))\), for sufficiently
large \(n\) this is smaller than the valuation of every root.
The nonarchimedean triangle law then gives
\(v_\ell(a(t_n))=\mu+j/(\ell^{n-1}(\ell-1))\).
The lower strict inequality follows because all reduced inverse
traces coincide for the symmetric residual triple.

Finally fix \(n\) and \(y\). The valuation of the nonzero
difference is constant on a sufficiently small open
neighborhood of \(h\) and on its \(G_{L_n}\)-conjugates.
Chebotarev over \(L_n\) gives a positive-density set of primes
with Frobenius in this set. Prime ideals of residue degree
greater than one over \(\QQ\) have density zero. Since
\(L_n/\QQ\) is Galois, an unramified rational prime below a
degree-one prime splits completely. Finitely many bad places
can be discarded.
\end{proof}

\begin{proposition}[The first cyclotomic layer admitting crystalline descent]
\label{prop:descentlayer363}
Fix \(y\) and \(n\) as in Theorem~\ref{thm:slthree363}, and put
\(U=U_y|_{G_{L_n}}\). For \(0\le b\le n\), the unique
determinant-one extension of \(U\) to \(G_{L_b}\), up to
isomorphism, is \(U_y|_{G_{L_b}}\). At a chosen place over
\(\ell\) its projective Weil--Deligne inertia has order
\(\ell^{n-b}\). Every extension to \(G_{L_b}\) has the same
projective inertia there. In particular \(n\) is the least
layer permitting a crystalline extension, even with arbitrary
determinant. There is no representation of \(G_K\) crystalline
above \(\ell\) whose restriction to \(G_{L_n}\) is \(U\).
\end{proposition}
\begin{proof}
Identify the restrictions of two extensions over an algebraic
closure of the coefficient field. For each \(g\in G_{L_b}\)
their matrices induce the same conjugation on the image of
\(G_{L_n}\), by normality. The ratio of their matrices is scalar
by the absolute irreducibility of \(U\). The ratio is a character
of the cyclic quotient \(\operatorname{Gal}(L_n/L_b)\), of
order dividing \(\ell^{n-b}\). If both determinants are one
its cube is one and it is trivial.

Locally \(L_0\) is unramified over \(\QQ_\ell\), while the
cyclotomic degree-\(\ell^b\) layer is totally ramified. The
image of inertia of \(L_b\) in the rational degree-\(\ell^n\)
cyclotomic quotient therefore has order \(\ell^{n-b}\).
For instance the compositum of the local field \(L_0\) with the
degree-\(\ell^n\) cyclotomic layer is totally ramified over its
degree-\(\ell^b\) layer with exactly that degree. Thus the
restriction of \(\nu_{t_n}\) has this exact order. On
Weil--Deligne inertia the three characters of \(U_y\) are its
powers \(2,1,-3\). The projective image has order dividing
\(\ell^{n-b}\) and achieves it, since the ratio of the first
two characters is \(\nu_{t_n}\). Scalar twists preserve this
image. For \(b<n\) the nontrivial projective inertia rules
out crystallinity, and for \(b=n\) Theorem~\ref{thm:slthree363}
gives crystallinity. A crystalline representation on \(K\)
would stay crystalline after restriction to \(L_0\), in
contradiction to the case \(b=0\).
\end{proof}

Theorem~\ref{thm:auxdescent361} realizes the full strict module
\(C\) on each of these same \(L_n\). Corollary~\ref{cor:skewsize363}
gives a nonzero classical inverse trace on \(G_{L_n}\) and
bounds its valuation along the tower for a fixed component.
Proposition~\ref{prop:descentlayer363} excludes these points
from the crystalline descents counted on \(K\). The local
calculation distinguishes having the required weights and
types over an auxiliary field from obtaining a point in
the trace algebra on the original field. Controlling the full order in
Theorem~\ref{thm:classicaldefect360} still requires an
integral comparison with all of \(C\) on the original interval.
The corresponding comparison of the two signed periods is
specified in Section~\ref{sec:missing}.

\section{Ordinary curves and cuspidal components}\label{nav:chapter-ordinary-families}
\subsection{Global units and ordinary traces in residue characteristic}
\label{sec:equalweights364}

Keep \(L=L_0\), \(F=F_0=L^+\), \(d=[F:\QQ]\), and the chosen
places over \(\ell\) of Section~\ref{sec:ordinaryfamily363}.
Write \(k\) for the residue field of \(O\), enlarged to contain
the values of the three residual characters. Put
\[
 H=\prod_{v\mid\ell} H_v,\qquad H_v=1+\ell\mathcal O_{F_v},
 \qquad \Lambda=O[[H\times H]] .
\]
Each completion is unramified and \(\ell\) is odd, so logarithm
identifies \(H\) with a free \(\ZZ_\ell\)-module of rank \(d\).
Two variable characters \(\lambda_1,\lambda_2\) of \(H\)
specify the three weight variations with product one by taking
\(\lambda_3=(\lambda_1\lambda_2)^{-1}\). A character of \(H\)
is evaluated on a local unit by removing its Teichm\"uller part
and using the corresponding factor. For a real global unit
\(e\in\mathcal O_F^\times\), denote by \(j(e)\in H\) the
tuple of all these principal parts. On principal units we use
local reciprocity in the same convention at both places of \(L\)
over each \(v\).

\begin{theorem}[The weights of sums of three global characters]
\label{thm:unitweights364}
Let \(R\) be a complete discrete valuation ring of characteristic
\(\ell\) with finite residue field containing \(k\), with the
indicated copy of \(k\) in \(R\). Take continuous
\(\lambda_i:H\to1+\mathfrak m_R\) as above. At the chosen
places prescribe the ordered inertial characters
\[
 \bar\kappa_\ell^r\lambda_1,\qquad \lambda_2,\qquad
 \bar\kappa_\ell^{-r}\lambda_3 ,
\]
and at the conjugate places their reversed conjugate inverses.
There exist three continuous global characters of \(G_L\)
with residual characters \(\bar\Psi,1,\bar\Psi^{-1}\), with
these ordered inertial characters and with conjugate self-dual
direct sum, after a finite coefficient extension, if and only if
\begin{equation}
             (\lambda_1/\lambda_3)(j(e))=1
                      \quad\text{for all } e\in\mathcal O_F^\times .
                                                        \label{eq:units364}
\end{equation}
When they exist their finite inertial characters away from
\(\ell\) are those of the original residual triple and their
product is one. In \(\operatorname{Spec}(\Lambda/\varpi)\),
where \(\varpi\) is a uniformizer of \(O\), the reduced weight
locus defined by \eqref{eq:units364} has codimension \(d-1\).
\end{theorem}
\begin{proof}
The ratio of a global character to its residual character,
lifted through the finite constant field, takes values in
principal units. This latter group is pro-\(\ell\) and
torsion-free. In particular in characteristic \(\ell\) the
identity \((1+x)^{\ell^a}=1+x^{\ell^a}\) excludes nontrivial
\(\ell\)-power torsion. Away from \(\ell\) the local unit group
is an extension of a finite group by a pro-\(p\) group with
\(p\ne\ell\). Its image in these principal units is trivial.
The ratio is thus unramified away from \(\ell\).

The residual characters at a chosen inertia group are distinct,
since \(2r<\ell-1\) and residual cyclotomic inertia is full.
The global character reducing to \(\bar\Psi\) therefore has
variation \(\lambda_1\) at the chosen place and
\(\lambda_3^{-1}\) at the conjugate place, after identification
by complex conjugation. In fact \(\bar\Psi\) comes from
\(\QQ\), so the reversed inverse with this residual value is
the inverse of the third character. A principal idele has
trivial global Artin symbol. Applying the character ratio to
the idele of a real unit \(e\) leaves just the product on
the left of \eqref{eq:units364}. This proves necessity.

Conversely prescribe this first variation on the principal
units of \(L\) above \(\ell\), using \(\lambda_1\) on one side
and \(\lambda_3^{-1}\) on the other. It kills all global units.
Indeed for \(e'\in\mathcal O_L^\times\) the quotient
\(e'/\overline{e'}\) is an algebraic integer and every complex
conjugate has absolute value one. It is a root of unity by
Kronecker's theorem. A positive power of \(e'\) is therefore
a real unit. The given character kills that power by
\eqref{eq:units364}, and the target is torsion-free.
The corresponding local variation for the middle character
uses \(\lambda_2\) and \(\lambda_2^{-1}\). It kills real units
by cancellation and hence all global units in the same way.

For clarity, the ray class sequences at the powers of the
primes above \(\ell\), followed by their pro-\(\ell\) inverse
limit, identify the image of the local principal units in the
maximal abelian pro-\(\ell\) quotient unramified outside
\(\ell\) with their quotient by the closure of the projected
global units. The quotient of that image in the full group is
the \(\ell\)-primary ideal class group. These are the unit,
local-residue, ray-class and ideal-class terms of the ray
class exact sequence. Projection to the pro-\(\ell\) factors
removes the local prime-to-\(\ell\) roots of unity.
Thus each prescribed character defines a continuous character
on a subgroup of finite \(\ell\)-power index. If \(\ell^a\)
kills the quotient, its extension takes an element \(g\) to
the unique \(\ell^a\)-th root of the prescribed value on
\(g^{\ell^a}\). All values lie in the single finite extension
obtained by adjoining the \(\ell^a\)-th roots of the fraction field of \(R\),
which is a Laurent series field over a finite field. The
root map is multiplicative and continuous, and the roots of
principal units are principal units. Uniqueness makes the
middle character conjugate inverse to itself. Multiply the
first variation by \(\bar\Psi\), take its conjugate inverse
as the third character, and use the constructed middle
character. Their determinant is unramified everywhere,
reduces to one and has torsion-free image. Finiteness of the
class group makes it one. This proves sufficiency.

For the codimension, Leopoldt's theorem for totally real
abelian fields gives rank \(d-1\) for the closure of
\(j(\mathcal O_F^\times)\) in \(H\). Since
\(\lambda_1/\lambda_3=\lambda_1^2\lambda_2\), the subgroup
tested in \(H\times H\) is the image of that closure under
\(h\mapsto(h^2,h)\). This is an embedding with projection
back onto \(H\) in the second factor. In Smith coordinates on
\(H\times H\simeq\ZZ_\ell^{2d}\) the subgroup has generators
\(\ell^{a_i}e_i\) for \(1\le i\le d-1\). In the power series
coordinates of \(k[[H\times H]]\) the equations are
\((1+X_i)^{\ell^{a_i}}-1=X_i^{\ell^{a_i}}\). Their radical
is \((X_1,\ldots,X_{d-1})\), proving the assertion, including
the possible non-saturation of the original unit
lattice.~\cite{VII-V-brumer330}
\end{proof}

\begin{proposition}[The integral Hecke algebra on the weight space of product one]
\label{prop:bigweight364}
There is a finite local domain \(B'\) containing \(\Lambda\),
finite as a \(\Lambda\)-module and of dimension \(1+2d\), with
a characteristic-zero specialization at the cuspidal adjoint
\(U_0\). It carries a continuous trace \(t:G_L\to B'\)
interpolating the normalized ordinary Galois traces.
Let a prime of \(B'\) lie over the kernel of a surjective
continuous weight map \(\Lambda\to k[[T]]\). Over a finite
extension of the fraction field of \(B'\) modulo this prime,
the resulting trace comes from a continuous semisimple Galois
representation. It has determinant one, is conjugate self-dual,
reduces semisimply to \(\bar\Psi\oplus1\oplus\bar\Psi^{-1}\),
and has local ordinary flags with the prescribed weight
characters.
\end{proposition}
\begin{proof}
In the construction of Proposition~\ref{prop:ordinaryfamily363}
use the entire factor of the weight algebra on which the product
of the three inertial characters is one. In the relative
coordinates at the adjoint it is precisely \(\Lambda\).
Base change of the same free ordinary module gives a finite
free module over \(\Lambda\). We justify the density assertion
needed to interpolate here in several variables.

At every chosen \(\ell\)-adic embedding vary the first two
algebraic exponents by integers \(a_\tau,b_\tau\) and the last
one by \(-a_\tau-b_\tau\). Choose all these integers divisible
by the orders of the local residue unit groups, preserving
the prime-to-\(\ell\) diamonds. Impose the inequalities
making the resulting coefficient weight strictly dominant at
each embedding. These integer tuples are dense in
\(\ZZ_\ell^{2d}\). In any fixed residue classes modulo
\(\ell^m\) choose integral representatives with the required
prime-to-\(\ell\) divisibility by the Chinese remainder theorem.
With \(b_\tau\) fixed one can increase \(a_\tau\) by arbitrary
positive multiples of both moduli, ensuring both consecutive
gaps are sufficiently large.

A series in \(\Lambda\) vanishing at these weights is zero.
In logarithmic coordinates on an open neighborhood of the
adjoint, the map from embedding exponents to unit character
values has invertible Jacobian over the coefficient field.
Indeed its blocks are the matrices
\((\tau(\log h_j))_{\tau,j}\) for principal-unit bases
\(h_j\). Their logarithms are bases of the local fields
over \(\QQ_\ell\), and the embeddings identify the scalar
extension of each separable local field with a product.
Continuity and the density just proved give vanishing of the
composite series on a product of sufficiently small
\(\ZZ_\ell\)-balls about the origin. Iterating the
one-variable identity theorem in each coordinate forces its
analytic germ to be zero, also over the larger coefficient
field. The formal inverse to the change of coordinates then
forces the original Taylor series to be zero. The weight
coordinates are relative to the adjoint, so this is the
original power series itself.

Now form the good split Hecke algebra inside the endomorphisms
of the free module. The classical fibers at the regular
weights just used are semisimple for these operators. Applying
the density assertion to their matrix entries gives a faithful
embedding in the product of the classical coefficient rings,
also if the adjoint evaluation is included. The compactness
and Chebotarev argument of Proposition~\ref{prop:ordinaryfamily363}
produces \(t(h)\) for all \(h\). Continuity of \(t\) follows
because the compact Hecke algebra embeds homeomorphically onto its
image in that Hausdorff product. Take a minimal prime below
the adjoint specialization exactly as there. This gives the
finite domain \(B'\), integral over \(\Lambda\) with zero
contraction of the minimal prime. Integral extension preserves
dimension. Its residue field is that of \(O\) by evaluation
at the adjoint. The ordinary flags at the regular weights and
their character formulas are those of Geraghty, as used
in Lemma~\ref{lem:ordinaryflags363}.~\cite{VII-V-geraghty363}

Lying over supplies a prime above any specified kernel in the
statement. The quotient domain is finite over \(k[[T]]\).
Its fraction field is a finite extension of \(k((T))\) and
its normalization is a complete discrete valuation ring there.
The trace satisfies the degree-three pseudocharacter identities
since they are polynomial trace identities at every regular
classical point. Over the algebraic closure of the new fraction
field a degree-three pseudocharacter is the trace of a unique
semisimple representation up to isomorphism. Here \(3!\) is
invertible, as required by the trace formulation of this
theorem.~\cite{VII-V-chenevier364}

We give the coefficient and continuity arguments. In the
semisimple representation let \(D\) be the span of the group
matrices over the algebraic closure. It is a product of full
matrix algebras on the distinct simple constituents. This
follows by factoring the acting algebra by its radical, which
kills a semisimple faithful module, and applying the structure
theorem for finite-dimensional semisimple algebras.
The trace pairing on \(D\) is nondegenerate, since the
multiplicity of each constituent is at most three and is
nonzero in the coefficient field. Choose group matrices
forming a basis of \(D\) and adjoin their entries, a finite
coefficient extension. Every group matrix is defined over it.
Indeed its coordinates in that basis are recovered from
its traces against the basis by inverting the trace Gram
matrix. Those traces are continuous functions of the group
element by the construction of \(t\), so the representation
is continuous. The compact matrix image preserves the integral
span of the translates of any fixed full lattice.
This span is bounded by compactness and therefore is a
finite free module. The residual characteristic polynomials
are those at the adjoint by reduction in the local ring \(B'\).
Brauer--Nesbitt gives the claimed residual semisimplification.
The identity \(t(chc^{-1})=t(h^{-1})\), for complex
conjugation \(c\), interpolates from the classical fibers and
gives conjugate self-duality by semisimplicity.

At each chosen place take an inertia element with separating
Teichm\"uller values as in Lemma~\ref{lem:ordinaryflags363}.
The three projectors there have denominators which are units
over \(\Lambda\). The vanishing traces of the lower matrix
entries against all group matrices, and the three identities
recovering the diagonal on inertia, interpolate from the
regular weights by the multivariable density argument.
At the specialized representation these products of projectors
and group matrices belong to \(D\). Its trace pairing is
nondegenerate even for a reducible semisimple representation,
as just proved. The lower entries thus vanish and the
Vandermonde argument gives the ordered characters. Their
product is one. The determinant reduces to one globally and
is unramified at \(\ell\), also at the conjugate places.
Away from \(\ell\) it is unramified by the principal-unit
argument in Theorem~\ref{thm:unitweights364}. Finiteness of
the class group and torsion-freeness of the target make the
determinant one.
\end{proof}

\begin{lemma}[A common auxiliary prime of large residue degree]
\label{lem:inertaux364}
The prime \(w\) of Proposition~\ref{prop:auxfield361} can
be chosen so that every \(F_v\) for \(v\mid\ell\) contains
the unramified degree-\(\ell\) extension of \(\QQ_\ell\).
\end{lemma}
\begin{proof}
Let \(M=\QQ(\mu_{sN_*\ell})\). At every place above \(\ell\)
the normalized valuation of the rational integer \(\ell\)
is \(\ell-1\), since \(\ell\nmid sN_*\). Thus \(\ell\) is
not an \(\ell\)-th power in \(M\). The splitting field of
\(X^\ell-\ell\) over \(M\) is a cyclic extension of degree
\(\ell\) by Kummer theory, and the compositum is Galois over
\(\QQ\). Choose a prime with Frobenius in the conjugacy class
of a nonidentity element of this cyclic subgroup, excluding
the finite set of previously forbidden primes. Chebotarev
supplies such a \(w\). It splits in \(M\), so
\(w\equiv1\bmod sN_*\ell\). The integer \(\ell\) is not an
\(\ell\)-th power modulo \(w\), since the Frobenius on the roots
would then fix them all. The quotient of \(\mathbf F_w^\times\)
by its \(\ell\)-th powers has order \(\ell\) and corresponds
to \(F_w\subset\QQ(\mu_w)\). Frobenius of \(\ell\) generates
this quotient. The completion of \(F_w\) at \(\ell\) is thus
unramified of degree \(\ell\). The field \(F\) contains \(F_w\).
All the choices here only use the interval and \(\ell\).
\end{proof}

\begin{theorem}[A nonabelian ordinary trace curve in characteristic \(\ell\)]
\label{thm:rescurve364}
Choose \(w\) as in Lemma~\ref{lem:inertaux364}. On the common
\(L_0\) there is one map of weight variations
\(\Lambda\to k[[T]]\) which works for every \(r,\alpha\) on
the interval, after extending the finite constant field as
needed. On every component \(B'\) of
Proposition~\ref{prop:bigweight364} through the corresponding
adjoint, each prime over its kernel has one-dimensional
quotient. Its generic trace cannot be a sum of three
characters. The three ordered ordinary characters are
pairwise distinct even on the kernel of the full cyclotomic
character on inertia at every chosen place above \(\ell\).
The coefficient \(T\) itself lies in
the algebra generated over the constants by the specialized
group traces.

Over the algebraic closure of its coefficient field the
generic representation either has Zariski closure
\(\mathrm{SL}_3\) or is a
direct sum of a character reducing to \(1\) and an
irreducible two-dimensional representation with projective
Zariski closure \(\mathrm{PGL}_2\). In the second case the
one-dimensional summand satisfies \(\chi^c=\chi^{-1}\).
\end{theorem}
\begin{proof}
Take a real unit \(e\in F\) of infinite order. There is such a
unit since \(F\) contains the real quadratic field
\(\QQ(\sqrt{sv})\). Its principal part at every \(v\mid\ell\)
has nonzero logarithm, since an algebraic unit which is a
local root of unity is a root of unity. At each \(v\) put
\[
 u_v(z)=\frac{\log N_{F_v/\QQ_\ell}(z)}{\log(1+\ell)}
       \quad(z\in H_v).
\]
It is onto by the norm calculation in
Section~\ref{sec:ordinaryfamily363}. Its kernel has rank
\([F_v:\QQ_\ell]-1>0\) by Lemma~\ref{lem:inertaux364}.
Choose integral linear functionals \(\beta_v:H_v\to\ZZ_\ell\)
in logarithmic coordinates such that each \(\beta_v\) is
nonzero on this kernel and
\(\sum_v\beta_v(j(e)_v)\ne0\). For the latter condition,
start with any choices satisfying the former. At one place
add an integral multiple of a functional nonzero at the
local principal part of \(e\). At most one multiple can
make the new restriction to the kernel identically zero
and at most one makes the required sum zero.
We use the continuous
homomorphism \(\ZZ_\ell\to k[[T]]^\times\),
\(a\mapsto(1+T)^a\), obtained by reducing the integral
binomial coefficients. It is injective. If
\(a=\ell^m b\ne0\) with \(b\) a unit, the coefficient of
\(T^{\ell^m}\) in \((1+T)^a=(1+T^{\ell^m})^b\) is
\(\bar b\ne0\).

At \(v\) define \(\lambda_1,\lambda_2,\lambda_3\) as powers
of \(1+T\) with respective exponents
\[
                2u_v+\ell\beta_v,\qquad
                u_v,\qquad
               -3u_v-\ell\beta_v .
\]
They have product one. At any \(z\) with \(u_v(z)=1\), all
three exponents are nonzero and pairwise different modulo
\(\ell\), since \(\ell>5\). For the stronger separation in
the statement take \(z\in\ker u_v\) with \(\beta_v(z)\ne0\).
The exponents are then \(\ell\beta_v(z),0,-\ell\beta_v(z)\),
pairwise distinct in \(\ZZ_\ell\). An inertia lift of this
principal unit has cyclotomic value one by local reciprocity.
At a unit with \(u_v=1\) the second character takes value
\(1+T\). The resulting
map of completed group algebras is surjective onto \(k[[T]]\),
since its compact image contains the polynomial ring \(k[T]\).
Every domain above its kernel is finite and integral over
\(k[[T]]\), hence has dimension one.

For a real unit the sum of the \(u_v\) on its principal parts
is zero. The product of its local norms is the rational norm
of the unit, which is \(1\) or \(-1\), and taking principal
parts and logarithms kills that sign. On the selected \(e\)
the exponent of \((\lambda_1/\lambda_3)(j(e))\) is therefore
\(2\ell\sum_v\beta_v(j(e)_v)\ne0\). By injectivity of
\(a\mapsto(1+T)^a\), equation~\eqref{eq:units364} fails.
The representation furnished by Proposition~\ref{prop:bigweight364}
cannot be a sum of three characters. This argument works
after any finite coefficient extension. A decomposition over
the algebraic closure would be defined over the finite
extension generated by the change-of-basis entries.
The assertion follows over that closure as well.

To recover \(T\) from the traces alone, lift a Teichm\"uller
unit of norm generating \(\mathbf F_\ell^\times\) to local
inertia. Write this lift as \(s_0\). Its three eigenvalues
in the specialization are the distinct constants
\(\bar\kappa_\ell^r(s_0),1,\bar\kappa_\ell^{-r}(s_0)\),
since all principal-unit variations are trivial on this local unit.
Lift a principal unit with \(u_v(z)=1\) to an inertia element
\(h\). For \(j=0,1,2\) the traces at \(s_0^jh\) give the
three linear combinations of the diagonal values at \(h\)
with this constant Vandermonde matrix. Solving for the middle
value gives \(1+T\) in the algebra generated by the traces.

Suppose first that the representation is absolutely
irreducible. The normal-core argument of
Theorem~\ref{thm:slthree363} applies in this characteristic
as well. The sum of the translates of a simple submodule on
a normal subgroup of finite index is semisimple and is the
whole space. The scalar alternative is excluded because the
three eigenvalues at the principal unit just used have
ratios of infinite order. Thus they remain distinct upon
raising the matrix to any positive power. The three-line
alternative would give a transitive cyclic permutation group.
Indeed the kernel of \(\bar\Psi\) has pro-\(\ell\) image in
a stable lattice, so acts trivially on three permuted lines.
The residual image of the semisimplification is cyclic.
The resulting cubic self-twist, reduced in the residue field,
would permute the residual triple by multiplication by an
order-three character. Its orbit of \(1\) cannot contain
\(\bar\Psi\), whose inertial order exceeds four. Thus the
representation stays irreducible on every open subgroup.

The connected algebraic monodromy is consequently semisimple
and irreducible, by the unipotent radical and determinant
argument of Theorem~\ref{thm:cyclo119}. The rank-three
classification in the proof of Theorem~\ref{thm:orth122}
applies in characteristic \(\ell>5\). A connected proper
possibility would preserve a nondegenerate symmetric
form, unique up to scalar. The full group would then be
self-dual up to a cubic character. At the principal unit
\(h\) all three eigenvalues are principal units, so any
cubic scalar relating them to their inverses reduces to
one and hence is one. A reciprocal multiset of odd size
contains an eigenvalue of square one. All three eigenvalues
here have infinite order. This excludes the orthogonal case
and proves full special linear monodromy.

In the reducible case there is a unique one-dimensional
constituent, since three characters have been excluded.
Conjugate self-duality sends this constituent to its conjugate
inverse, so the two characters agree. Each of the three
residual characters comes from \(\QQ\) and is invariant under
conjugation. The only one equal to its inverse is \(1\).
It follows that the two-dimensional constituent reduces
semisimply to \(\bar\Psi\oplus\bar\Psi^{-1}\).
At a chosen place its eigencharacters are the two outer
characters, by residual separation. Their ratio at \(h\)
has infinite order, proving that its projective image is
infinite. If the identity component of its projective
closure were a proper nontrivial subgroup of
\(\mathrm{PGL}_2\), it would be solvable. Indeed a nontrivial
connected semisimple group has dimension at least three,
and \(\mathrm{PGL}_2\) has dimension three. A triangularizable
connected subgroup with nontrivial unipotent radical fixes
a unique line, which its normalizer must preserve. This
contradicts irreducibility. A nontrivial torus fixes two
lines. The normalizer must interchange them to act irreducibly,
giving a quadratic self-twist by conjugation by the matrix
with diagonal \(1,-1\) on those lines. Its reduction would
interchange \(\bar\Psi,\bar\Psi^{-1}\), forcing
\(\bar\Psi^4=1\). This contradicts the order on inertia.
Therefore the projective closure is \(\mathrm{PGL}_2\).
The unit and all weight variations above depend only on
the common field and on \(\ell\). This gives simultaneity
in \(r,\alpha\) as asserted.
\end{proof}

The curve has separated characters on cyclotomic-kernel inertia and admits no
global triangular representation with its trace. These are
two of the conditions on the one-dimensional primes in
\srcloc{Development/AutomorphicArithmetic/AutomorphicVarieties/CurveDVRSpecialization/PropagationTheoremGL2}{16}.
The local degrees satisfy \([F_v:\QQ_\ell]\ge\ell>5\)
and \(d=2\ell\), since the two real factors of \(F\) have
coprime degrees. In the original fixed-weight problem the three variations on
principal units reduce to one. No finite extension of \(L\)
can make the variations on this curve all equal, since their ratios
already have infinite order on a single inertia element.
Passing from the variable-weight points to the original
strict group requires control of the other primes in the
deformation ring, as in the propagation implication at
line~36 of the same source file. The uniform length estimate
and the signed integral period comparison in
Section~\ref{sec:missing} still enter the deduction of
Dirichlet GRH.


\subsection{Strict classes and ordinary traces on one cyclotomic layer}
\label{sec:commonframe365}

Use the polynomial-conductor interval \(V\le x\le V^{3/2}\)
of Theorem~\ref{thm:heckedivisor345}, and fix
\(\ell>\exp(V^\eps/12)\). Choose the prime \(w\) as in
Lemma~\ref{lem:inertaux364}. Write \(F=F_a^+\) and \(L=L_a=KF\)
for the fields of Proposition~\ref{prop:auxfield361}, where \(a\)
is allowed to vary. For every \(x,r,\alpha\) of that proposition
let \(k\) be a finite field containing the values of
\(\bar\Psi=\bar\alpha\bar\kappa_\ell^r\), and put
\(\beta_i=\bar\Psi^{-i}\eta\) for \(i=1,2\).
We use \(G_D\) for the absolute Galois group of a field \(D\).
Let \(\Gamma_a\) be \(G_F\) modulo the closed normal subgroup
generated by
\[
       I_u\cap\ker\eta\cap\ker\bar\Psi,\qquad u\nmid\ell
\]
at finite places. All three characters
\(\eta,\bar\Psi,\kappa_\ell\) descend to this quotient.
Its index-two subgroup \(\Gamma_a^0=\ker\eta\) is the image
of \(G_L\). The field \(L\) splits over \(F\) above \(\ell\).
We use every completion of \(F\) above \(\ell\), and both places
above it for a group over \(L\). A strict class below means
a first cohomology class restricting to zero in all these
local first cohomology groups. Throughout this section finite
coefficients carry the discrete topology.

\begin{proposition}[Two strict classes from the same prime]
\label{prop:strictclasses365}
For each of the two characters \(\beta_i\) there is a class over
\(\QQ\) whose restriction is nonzero in the strict group
of \(H^1(\Gamma_a,k(\beta_i))\) at every depth \(a\).
The same \(w\) serves all \(x,r,\alpha\).
\end{proposition}
\begin{proof}
Put \(D=k(\beta_i)\). Over \(\QQ\) impose zero at \(\ell\),
the full local group at \(w\), and the unramified local
condition at all other finite primes. At the real place the
first cohomology vanishes. Denote these conditions by
\(\mathcal L\) and their orthogonals under local Tate duality by
\(\mathcal L^\perp\). The dimension formula from Poitou--Tate
duality gives
\[
 \dim H^1_{\mathcal L}(\QQ,D)
 -\dim H^1_{\mathcal L^\perp}(\QQ,D^\vee(1))
 =h^0(\QQ,D)-h^0(\QQ,D^\vee(1))
    +\sum_u\bigl(\dim\mathcal L_u-h^0(\QQ_u,D)\bigr).
\]
Here \(h^0\) denotes dimension over \(k\). One may use the
nine-term duality sequence with a finite set containing
\(\ell,w,\infty\) and the ramified primes of \(D\). Subtracting
the dimensions in its middle terms and inserting the local
and global Euler characteristics gives the displayed formula.
The omitted unramified primes contribute zero. This also
explains the real-place term when \(\ell\) is odd.~\cite{VII-V-V-milne98}

Both global invariant spaces vanish because their characters
are nontrivial on inertia at \(s\), where only \(\eta\) among
\(\eta,\bar\alpha,\bar\kappa_\ell\) ramifies. At the real place
we have \(\beta_i(c)=-1\). At \(\ell\) the character on inertia
is \(\bar\kappa_\ell^{-ir}\), nontrivial since \(0<ir<\ell-1\).
Both local terms are therefore zero. At each other prime with
the unramified condition the dimension is
\(\dim D^{I_u}/(\mathrm{Frob}_u-1)D^{I_u}=h^0(\QQ_u,D)\).
At \(w\) both \(\beta_i\) and \(\bar\kappa_\ell\) are trivial
on the decomposition group, by \(w\equiv1\bmod sN_*\ell\).
The full local \(H^1\) has dimension two. Indeed the
uniformizer and the cyclic \(\ell\)-part of the residue units
each contribute one independent homomorphism to \(k\), and
the principal units are pro-\(w\). Thus the right hand side is
one. In particular a nonzero class exists for each \(i\).

Restriction to \(G_F\) is injective on these classes.
The field \(F\) is Galois and totally real, so \(D^{G_F}=0\)
by oddness. Inflation--restriction gives the injection even
when \(\ell\) divides \([F:\QQ]\). The local zero conditions
persist. At \(w\) a cocycle is a homomorphism on the entire
decomposition group. Its restriction to inertia factors through
the quotient of the tame \(\ZZ_\ell(1)\) by \(\ell\).
The degree-\(\ell\) field \(F_w\) is totally and tamely ramified
at \(w\). Its inertia subgroup has zero image in this quotient,
so restriction to every completion of \(F\) kills the inertia
homomorphism. Elsewhere away from \(\ell\), the restriction of
the cocycle to inertia was a coboundary. Its value on
\(\ker\eta\cap\ker\bar\Psi\) inside inertia is zero, since
this subgroup acts trivially on \(D\). The restriction of the
cocycle therefore vanishes on the closed normal subgroup used
to define \(\Gamma_a\). Indeed on the kernel of the action it
is additive, and conjugating an element of value zero there
again gives value zero. This proves descent to the quotient
and preserves the nonzero strict class.
\end{proof}

\begin{proposition}[Uniform residual dimensions along the auxiliary tower]
\label{prop:classrank365}
There are integers \(h,g\), depending on this interval,
\(\ell\) and the auxiliary fields at depth zero, such that
at every depth and for all the pairs under consideration
\begin{align*}
 \dim_k H^1_{\rm str}(\Gamma_a,k(\beta_i))&\le h
                                  &&(i=1,2),\\
 \dim_k H^1_{\rm str}(\Gamma_a^0,k(\bar\Psi^{-1}))&\le h,&\\
 \dim_k H^1(\Gamma_a,k)&\le 1+g+h.&
\end{align*}
\end{proposition}
\begin{proof}
Set \(T_a=L_a E_*(\mu_\ell)\), let \(T_\infty\) be their union,
and put \(g=\#\{t:t\mid\ell\text{ in }T_0\}\).
These are the layers of the cyclotomic \(\ZZ_\ell\)-extension
of the abelian field \(T_0\). At \(\ell\) the ramification
index of \(T_0\) is \(\ell-1\). More precisely its local field
is an unramified extension of \(\QQ_\ell(\mu_\ell)\), since
all the other factors are unramified there. The rational
cyclotomic layers are totally ramified after this unramified
base extension. Consequently each place over \(\ell\) is
totally ramified at every step, and the number \(g\) is
constant.

Let \(X\) be the Galois group of the maximal unramified abelian
pro-\(\ell\) extension of \(T_\infty\). The cyclotomic class
group theorem of Iwasawa makes \(X\) a finitely generated torsion
module over \(\ZZ_\ell[[\operatorname{Gal}(T_\infty/T_0)]]\).
Its \(\mu\)-invariant is zero by the theorem of Ferrero and
Washington, since \(T_0/\QQ\) is abelian. The structure theorem
then makes \(X\) finitely generated over \(\ZZ_\ell\). Indeed
the cyclic factors up to finite kernel and cokernel are
quotients by powers of distinguished polynomials, all finite
free over \(\ZZ_\ell\). Set
\(h=\dim_{\mathbf F_\ell} X/\ell X\). The Hilbert
\(\ell\)-class field of \(T_a\) intersects \(T_\infty\) in
\(T_a\) by total ramification. After compositum with the tower
it is unramified and abelian over \(T_\infty\). Restriction
therefore exhibits the entire \(\ell\)-class group of \(T_a\)
as a quotient of \(X\). This bounds its \(\ell\)-rank by
\(h\).~\cite{VII-V-II-fw13,VII-V-II-iwasawa15}

The extensions \(T_a/F\) and \(T_a/L\) have degrees prime to
\(\ell\). Their degrees divide
\(2[E_*:\QQ](\ell-1)\), with \([E_*:\QQ]\le N_*<\ell\).
Inflate a class in either strict group to the absolute Galois
group and restrict to \(T_a\). This is injective by
restriction and corestriction. The coefficient character is
trivial on \(G_{T_a}\), and the resulting homomorphism kills
inertia away from \(\ell\) by the definition of \(\Gamma_a\).
It is zero locally over \(\ell\) by strictness. It consequently
factors through the class group of \(T_a\). Homomorphisms from
that group to the additive group of \(k\) have \(k\)-dimension
at most \(h\).

For trivial coefficients it suffices to work over
\(\mathbf F_\ell\), since extension to \(k\) just tensors the
space of continuous homomorphisms. Injection by the same
restriction gives invariant homomorphisms on \(G_{T_a}\)
unramified outside the \(g\) places over \(\ell\). Multiplication
by a generator of \(\mu_\ell\) identifies this image with a
subspace in the minus eigenspace of
\(H^1(G_{T_a,\{\ell,\infty\}},\mu_\ell)\) for complex
conjugation. Indeed conjugation fixes \(F\) and acts on that
generator by inversion. By Kummer theory the latter cohomology
is an extension of the \(\ell\)-torsion of the class group of
the ring of \(\ell\)-integers by its units modulo \(\ell\)-th
powers. This is the sequence obtained by mapping a Kummer
element with all away valuations divisible by \(\ell\) to
the class of its away divisor divided by \(\ell\).
The class group of the ring is a quotient of \(\Cl(T_a)\)
and hence contributes at most \(h\). Away from the ordinary
units, valuations of the \(\ell\)-units give a lattice of
rank at most \(g\). For any ordinary unit \(e\), the quotient
\(e/\bar e\) is a root of unity by Kronecker's theorem.
The minus part of the ordinary units modulo \(\ell\)-th
powers therefore has dimension at most one, since minus
projection is \((1-c)/2\) and the roots of unity are cyclic.
The valuation lattice is free, so reduction modulo \(\ell\)
preserves the injection of ordinary units modulo powers in
this sequence. Adding the three bounds proves the result.
\end{proof}

\begin{lemma}[The relative Euler characteristic for the ramification quotient]
\label{lem:frameeuler365}
Let \(W\) be any finite-dimensional continuous representation
of \(\Gamma_a\) over a finite field of characteristic \(\ell\).
For each \(u\mid\ell\) let \(b_u\subset W\) be stable under
\(G_{F_u}\), and let \(H^j_C\) denote the cohomology of the
mapping fiber of restriction from global cochains with values
in \(W\) to the sum of local cochains with values in \(W/b_u\).
If \(H^0(F,W^\vee(1))=0\), these groups are finite, vanish in
degrees above two, and satisfy
\[
 \dim H^2_C+\dim H^0_C+\sum_{u\mid\ell}[F_u:\QQ_\ell]\dim b_u
       =\dim H^1_C+\sum_{\tau:F\hookrightarrow\RR}\dim W^{c_\tau}.
\]
Here \(c_\tau\) denotes a complex conjugation at \(\tau\).
\end{lemma}
\begin{proof}
Let \(M/F\) be cut out by \((\eta,\bar\Psi)\). Its Galois
group \(\Delta\) has order prime to \(\ell\). The kernel of
\(\Gamma_a\to\Delta\) is the Galois group of the maximal
extension of \(M\) unramified at finite primes away from
\(\ell\). Indeed the removed inertia groups are precisely
the inertia groups of \(M\) at those primes. Their closed
normal closure over \(M\) is already normal over \(F\),
which permutes the primes of \(M\). Write \(S\) for the places
over \(\ell\) and the infinite places, and enlarge \(S\) to
\(S'\) by the away primes ramifying in \(M/F\). The same
description with \(S'\) gives \(G_{F,S'}\) above \(G_{M,S'}\).
In all such expressions the indicated places are pulled up
to the field in the subscript. Hochschild--Serre computes
the cohomology of either group over \(F\) by taking
\(\Delta\)-invariants of the cohomology over \(M\).
Invariants are exact by averaging.

We verify that replacing \(G_{F,S'}\) by \(\Gamma_a\) preserves
the Euler characteristic. Over \(M\) the module \(W\) is
unramified at each of the deleted primes \(t\) above
\(S'\setminus S\). The difference in the two cohomology
complexes is the sum of the residue-field complexes of
\(W(-1)\) shifted by one. To see this, use the comparison of
Galois and \'etale cohomology on the two rings of integers
with the indicated places inverted. It applies to the finite
locally constant sheaf here because \(\ell\) is invertible
on both rings. Excision in the localization sequence reduces
the relative term to a henselian local ring and its punctured
spectrum. Over a strictly henselian local ring the only
positive-degree cohomology of the punctured spectrum is
\(H^1(I_t,W)=W(-1)\). The wild inertia has order prime to
\(\ell\) in all finite quotients, and the \(\ell\)-factor
of tame inertia is \(\ZZ_\ell(1)\). The sheaf is unramified
and the degree-zero restriction is an isomorphism. Descent
to the residue field gives the asserted term. These comparison
maps commute with conjugation by \(\Delta\).~\cite{VII-V-V-milne98}

For a fixed \(u\in S'\setminus S\) take \(\Delta\)-invariants
of the sum over \(t\mid u\). This amounts to invariants under
the decomposition group \(\Delta_t\) at one \(t\). Equivalently
its two cohomology spaces are the degree-zero and degree-one
groups of \(G_{F_u}/I_{M_t}\) on \(W(-1)\), by Hochschild--Serre
and the prime-to-\(\ell\) order of \(\Delta_t\). The finite
inertia in this quotient also has order prime to \(\ell\).
After taking its invariants we are left with the kernel and
cokernel of \(\mathrm{Frob}_u-1\) on the same finite vector
space. Their dimensions are equal. Each deleted prime thus
changes the Euler characteristic by zero. Global finiteness
and vanishing above degree two for \(\Gamma_a\) follow over
\(M\) for \(G_{M,S}\) and pass to invariants. The global
Euler characteristic of \(G_{F,S'}\), since \(F\) is totally
real, now gives
\[
       \chi(\Gamma_a,W)
            =-[F:\QQ]\dim W+\sum_\tau\dim W^{c_\tau}.
\]

Over \(M\) the last local map in Poitou--Tate has cokernel
\(H^0(M,W^\vee(1))^\vee\) for restriction in degree two to
the places over \(\ell\). Real places have zero positive-degree
cohomology at this odd prime. Taking \(\Delta\)-invariants,
and applying local Hochschild--Serre above \(\ell\), identifies
this cokernel over \(F\) with \(H^0(F,W^\vee(1))^\vee=0\).
Restriction from \(H^2(\Gamma_a,W)\) to the sum of
\(H^2(F_u,W)\) is therefore surjective. It remains surjective
after passing to \(W/b_u\), by local cohomological dimension
two. Thus the mapping fiber has zero third cohomology. In
all finite degrees its Euler characteristic is the global one
minus the sum of the local ones. The local Euler characteristic
on \(W/b_u\) is \(-[F_u:\QQ_\ell]\dim(W/b_u)\), and the
local degrees sum to \([F:\QQ]\). Substitution proves the
formula. The duality and Euler formulas used here are
Theorems I.2.8, I.4.10 and I.5.1 of Milne, and the comparison
for the rings of integers is in II.2.9.~\cite{VII-V-V-milne98}
\end{proof}

\begin{theorem}[The cohomology hypotheses on one common layer]
\label{thm:commonframe365}
Take \(C_0=h+g\) from Proposition~\ref{prop:classrank365}.
For every sufficiently large \(a\) the arithmetic data
\((F_a^+,L_a,\Gamma_a)\) satisfy all the cohomological,
local generator, infinity and numerical hypotheses of the
source's ordinary deformation frame for all \(x,r,\alpha\)
on the interval, with this one \(C_0\). In particular
\[
    3C_0+7<[F_a^+:\QQ],\qquad [F_{a,u}^+:\QQ_\ell]>5
                      \quad(u\mid\ell).
\]
\end{theorem}
\begin{proof}
The three real factors of \(F_a^+\) have degrees \(2,\ell,\ell^a\).
The quadratic factor is disjoint from the other two by parity
of degree. The degree-\(\ell\) factor is totally ramified at
\(w\), where the rational cyclotomic tower of conductor a power
of \(\ell\) is unramified, so these two are disjoint as well.
The total real degree is \(2\ell^{a+1}\). At \(\ell\) the
completion contains both the degree-\(\ell\) unramified field
of Lemma~\ref{lem:inertaux364} and a totally ramified field
of degree \(\ell^a\). This gives the two inequalities, with
one lower threshold in \(a\).

Use complex conjugation at each real embedding and its image
in \(\Gamma_a\). It has order two, sign \(\eta=-1\), and
\(\bar\Psi=1\). The sign kernel is open, all local maps
and \(\kappa_\ell\) are continuous, and the residual character
has open kernel. The groups are profinite. The local groups
are the full absolute Galois groups of the completions.
Choose a place of \(L_a\) over each \(u\mid\ell\). They may
all be chosen above the restriction to \(K\) of any one fixed
place over \(\ell\) of \(\overline\QQ\). In fact \(F_a^+\)
and \(K\) are disjoint and Galois over \(\QQ\), so
\(\operatorname{Gal}(L_a/K)\) acts transitively on the
places of \(F_a^+\) over \(\ell\). Lift to the algebraic
closure. The conjugate places then give the other choice at
every \(u\), by splitting in \(L_a/F_a^+\).

At every such completion the full rational residual cyclotomic
inertia survives by Proposition~\ref{prop:auxfield361}.
Use an inertia element mapping to a generator of
\(\mathbf F_\ell^\times\) at the fixed place as the
distinguished element of \(\Gamma_a^0\), and use the same
choice locally at every other place. Its squared residual
\(\bar\Psi\)-value differs from one since \(0<2r<\ell-1\).
The element of \(G_{L_a}\) in that proposition with
\(\bar\kappa_\ell=-1\) and \(\bar\Psi=1\) rules out every
identity \(\bar\kappa_\ell=\bar\Psi^j\) on \(\Gamma_a^0\).
It also witnesses nontriviality of
\(\bar\kappa_\ell\), \(\bar\kappa_\ell\bar\Psi^{\pm1}\)
and \(\bar\kappa_\ell\bar\Psi^{\pm2}\eta\) where required.

On inertia in the kernel of the full cyclotomic character,
\(\bar\Psi=\bar\kappa_\ell^r=1\). Under local reciprocity the
principal units at \(u\) form a free \(\ZZ_\ell\)-module of
rank \([F_u:\QQ_\ell]\). The logarithm gives that rank on
an open subgroup, and there is no \(\ell\)-power torsion
because the ramification index \(\ell^a\) is prime to
\(\ell-1\). The norm of the principal units has rank one,
since sufficiently deep rational principal units already
have norms given by raising them to \([F_u:\QQ_\ell]\).
Its kernel is a direct summand. Choose a basis in that kernel
and extend it by one element to a basis of the principal
units. Lift the basis to inertia by reciprocity. All lifts
kill \(\bar\Psi\), and every lift but the last kills the full
cyclotomic character. A finite abelian \(\ell\)-group character
trivial on this basis kills all inertia, since the other unit
factor has order prime to \(\ell\).

Every open subgroup of \(\Gamma_a\) has only finitely many
continuous homomorphisms to \(\mathbf F_\ell\). Pull back to
the corresponding number field. Extensions arising in this
way are unramified outside the finite set consisting of the
places over \(\ell\) and the ramification of \(M/F\) in
Lemma~\ref{lem:frameeuler365}. Adjoining \(\mu_\ell\) gives
an injection by prime-to-\(\ell\) restriction, and Kummer
theory bounds the dimension by the ranks of an \(S\)-unit group
modulo powers and a finite class group. Any abelian
\(\ell\)-group character of \(\Gamma_a^0\) unramified above
\(\ell\) is everywhere unramified, since its allowed inertia
away from \(\ell\) has order prime to \(\ell\). The class
number of \(L_a\) therefore annihilates every such character.

For the zero local subspaces used in the strict groups, a
degree-one cocycle of the mapping fiber consists of a global
cocycle and a local vector witnessing its being a coboundary
at every place over \(\ell\). Each witness is unique because
the relevant local invariant spaces are zero. Indeed the
character on inertia is \(\bar\kappa_\ell^{-ir}\) or
\(\bar\kappa_\ell^{-r}\). Quotienting out the global
coboundaries thus gives exactly the strict groups of
Propositions~\ref{prop:strictclasses365}
and~\ref{prop:classrank365}. For the full local subspaces the
local quotients in the mapping fiber vanish, leaving ordinary
global \(H^1\). These results give both required nonzero
groups and all three upper bounds, also for the doubled set of
places over \(L_a\). Lemma~\ref{lem:frameeuler365} proves the
Euler inequality for every allowed coefficient module and
local subspace. The real embeddings number \([F_a^+:\QQ]\).
This verifies the hypotheses stated in
\srcloc{Development/AutomorphicArithmetic/AutomorphicVarieties/FrobeniusSmoothVectors/ContinuousCoboundaries2}{167}.
\end{proof}

\begin{proposition}[The ordinary trace on the ramification quotient]
\label{prop:tamecurve365}
The semisimple representation of
Theorem~\ref{thm:rescurve364} restricts on \(G_{L_a}\) to a
representation of \(\Gamma_a^0\) at every depth \(a\). Its
generic trace is nonconstant, its semisimplification is
nonabelian, and its three ordinary characters remain pairwise
distinct on the kernel of full cyclotomic inertia at every
chosen place. These properties hold on the layer of
Theorem~\ref{thm:commonframe365}.
\end{proposition}
\begin{proof}
First work on the integral domain \(B'\) of
Proposition~\ref{prop:bigweight364} over \(L_0\). At every
regular arithmetic weight each point above it supplies a
classical representation. These points remain Zariski dense
on \(B'\) after removing the vanishing locus of any nonzero
\(b\in B'\). Indeed for any nonzero \(b'\) the ideal
\(bb'B'\) contains a nonzero element of \(\Lambda\).
Indeed after inverting \(\Lambda\setminus\{0\}\) the finite domain
is a field, so one can clear denominators in the inverse
of \(bb'\). Use the density of the weights from the proof of
Proposition~\ref{prop:bigweight364} to choose a weight where
this element is nonzero and then use lying over. Every
point above that weight avoids the zeros of both \(b\) and \(b'\).

Only the finite set of places above the tame level and the
conductor of \(\alpha\) needs consideration away from \(\ell\).
At any one of these places, take the completion of \(L_0E_*\)
as in Proposition~\ref{prop:awaytype363}. On its inertia
all classical characteristic polynomials are unipotent after
passing to finite index bounded independently of the classical
point. In detail the residual image on this inertia is an
\(\ell\)-group, so the image factors through tame
\(\ZZ_\ell(1)\). Conjugation by Frobenius of norm \(q\)
permutes the three eigenvalues of a generator by raising them
to \(q\). Their orders divide \(q^6-1\). Taking powers to
enter this inertia bounds the orders on the original inertia
as well. Thus for every fixed inertia element there are only
finitely many possible characteristic polynomials. Each
coefficient on \(B'\) is constant: the polynomial in that
coefficient with the finitely many possible values as roots
vanishes by density, and \(B'\) is a domain. Equivalently the
polynomial may be taken over \(O\) by including coefficient
conjugates. Its root at the adjoint belongs to \(O\), and the
factor with the other distinct roots is nonzero there.
Consequently the constant polynomial is the inertial
characteristic polynomial at the adjoint for the element in
question.

For every one of the finitely many places, retain the
Frobenius resultant of Proposition~\ref{prop:awaytype363}.
Each is nonzero at the adjoint, as verified there using the
two eigenvalues of the elliptic newform. Off their joint
product's zero locus the unipotent inertia on \(L_0E_*\)
has zero logarithm, since conjugation would otherwise have
eigenvalue \(q\), forbidden by the resultant. Thus the
original inertia there has finite image. Its characteristic
polynomials just computed agree with
\(\alpha\oplus1\oplus\alpha^{-1}\). In particular at every
remaining regular point in this dense set each
\(j\in I_u\cap\ker\bar\Psi\), \(u\nmid\ell\) in \(L_0\),
acts trivially. Indeed \(\alpha\) has order prime to \(\ell\)
and there \(\bar\Psi=\bar\alpha\) on inertia. The matrix has
finite order in characteristic zero and all eigenvalues one.

For each such \(j\) and every \(g\in G_{L_0}\) we have proved
the identity \(t(jg)=t(g)\) on a Zariski dense set. It
therefore holds in \(B'\) and after specialization to the
characteristic-\(\ell\) curve. The trace pairing on the span
of the matrices of the semisimple specialization is
nondegenerate, by the argument of
Proposition~\ref{prop:bigweight364}. It follows that
\(\rho(j)-1=0\). For a place of \(F_a^+\) the inertia
away from \(\ell\) is contained in \(G_{L_a}\), since
\(L_a/F_a^+\) is unramified. It is then contained in an
inertia group over \(L_0\). Intersecting with
\(\ker\bar\Psi\) gives exactly elements of the kind just
killed. Conjugates by \(G_F\) still belong to such inertia
groups, and the kernel of the representation is closed.
This proves factorization through \(\Gamma_a^0\).

The subgroup \(G_{L_a}\) has finite index in \(G_{L_0}\).
In the irreducible case of Theorem~\ref{thm:rescurve364}
the connected monodromy \(\mathrm{SL}_3\) is preserved, and
in its reducible case the connected projective monodromy
\(\mathrm{PGL}_2\) of the two-dimensional constituent is
preserved. This follows since a connected algebraic group
cannot be a finite union of translates of a proper closed
subgroup. The restricted representation is thus semisimple
and nonabelian in both cases. A finite-index subgroup of the
original inertia contains a positive power of every prescribed
inertia element. At each place use the element in cyclotomic-
kernel inertia supplied by Theorem~\ref{thm:rescurve364}.
Its three eigenvalues are distinct powers of \(1+T\) with
exponents distinct in \(\ZZ_\ell\), so they stay distinct on
every positive power. At least one is of infinite order.
If every trace of the subgroup were constant over an
algebraic closure of \(k\), Newton's identities for the
first three powers would make its characteristic polynomial
constant too. All nonzero elements algebraic over the finite
field have finite order, a contradiction.
\end{proof}

The bound in Proposition~\ref{prop:classrank365} controls
dimensions over the residue field on a tower for a fixed
coefficient prime. The estimate in Section~\ref{sec:missing}
counts integral lengths on the original \(E_x\) and sums
over coefficient primes. Even one cyclic module of order
\(\ell^m\) has residual dimension one for arbitrarily large
\(m\). To obtain that estimate from ordinary traces one
must still control the fixed-weight integral relations in
Theorem~\ref{thm:classicaldefect360}. The same section
requires the comparison of the signed integral projected
theta periods on those original fields.


\subsection{A polarized ordinary curve through the split residual point}
\label{sec:markedcurve366}

Fix the common layer of Theorem~\ref{thm:commonframe365}.
Write \(H=\Gamma_a^0\) and \(g^\star=cg^{-1}c\) for \(g\in H\),
where \(c\) is the chosen complex conjugation. Work with the
representation \(\rho\) of Proposition~\ref{prop:tamecurve365}
over a finite extension of \(k((T))\). Further finite extensions
are allowed. Write \(g_0\) for the distinguished inertia element
of Theorem~\ref{thm:commonframe365}, \(b=\bar\Psi(g_0)\), and
\(\delta=b-b^{-1}\ne0\). A symmetric pairing on this
representation will mean a nondegenerate symmetric bilinear form
for which \(\rho(g^\star)\) is the adjoint of \(\rho(g)\).

\begin{lemma}[The pairing on the two-dimensional constituent]
\label{lem:curvesign366}
The representation \(\rho\) admits a symmetric pairing. In
the case of a two-dimensional constituent and a character, the
restriction of the pairing to each constituent is nondegenerate
and symmetric.
\end{lemma}
\begin{proof}
First work over an algebraic closure of the fraction field.
Let \(D\) be the span of the group matrices of \(\rho\).
Conjugate self-duality gives \(t(g^\star)=t(g)\) for the
trace \(t\). The operation \(\star\), extended linearly from
the group, descends to an anti-involution of \(D\). Indeed
if a group algebra element \(x\) acts by zero, then for every
\(g\) we have \(t(g^\star x^\star)=t(xg)=0\).
The nondegenerate trace pairing on \(D\) from
Proposition~\ref{prop:bigweight364} shows that \(x^\star\)
acts by zero. If \(\rho\) is irreducible, \(D=M_3\) and
this anti-involution is adjunction for a symmetric form.
In fact composing it with transposition gives an automorphism
of the matrix algebra and hence an inner automorphism. Applying
the resulting adjunction twice shows that the form matrix
has transpose equal to itself up to a scalar of square one.
A nonsingular alternating form is impossible in odd dimension.

In the other case \(D=M_2\times M_1\), and the two unequal
matrix factors are preserved. A priori the form on the first
factor could be alternating. We determine its sign using the
integral trace algebra \(B'\) of
Proposition~\ref{prop:bigweight364}. Use group elements of
\(G_{L_a}\) as arguments of that trace via \(G_{L_a}\subset
G_{L_0}\), choosing a lift of \(g_0\). In the group algebra
form \(Z=g_0-g_0^\star\). The element
\[
             u=t(Z^2)/2\in B'
\]
is a unit because it reduces to \(\delta^2\) on the three
residual characters. Set \(e=Z^2/u\) in the group algebra over
\(B'\). This expression can be used inside any trace without
choosing matrices over \(B'\).

There is a Zariski dense set of regular classical points
on \(B'\) with absolutely irreducible restriction to \(G_{L_a}\).
At the original adjoint this restriction is absolutely
irreducible by Theorem~\ref{thm:adjointseed361}. Choose nine
group elements there whose matrices span \(M_3\). The determinant
of their trace Gram matrix is nonzero at the adjoint and
remains nonzero on a dense set of regular classical points
by the density argument of Proposition~\ref{prop:tamecurve365}.
At each such point the matrices still span \(M_3\). The
conjugate self-duality gives a symmetric form there by the
odd-dimensional argument above.

At any of these points \(Z\) is skew adjoint. Its eigenvalues
are \(0,d,-d\) with \(d^2=u\), since \(Z\) is conjugate to
\(-Z^{\mathsf T}\) and its trace of the square is \(2u\).
The roots are distinct. Thus \(e\) is a self-adjoint idempotent
of rank two. Its image carries a nondegenerate symmetric
form, for its orthogonal complement is its kernel. On a
two-dimensional symmetric space the skew adjoint operators
form a one-dimensional vector space, as is seen by identifying
them with alternating bilinear forms. Consequently for any
group elements \(g,h\) the operators
\(e(g-g^\star)e\) and \(e(h-h^\star)e\) commute. Multiply
their commutator by any group element and take the trace.
The result belongs to \(B'\) and vanishes at this dense set,
so it vanishes in \(B'\). On the specialized curve the trace
pairing on \(D\) therefore gives
\[
        [\,e(\rho(g)-\rho(g^\star))e,\,
             e(\rho(h)-\rho(h^\star))e\,]=0 ,
\]
where \(e\) now denotes the image of the idempotent. All
identities extend linearly to the span of the group matrices.

On the one-dimensional block the anti-involution is the
identity, so \(Z\) acts by zero. Its trace of the square is
still \(2u\ne0\), and on the two-dimensional block it is
traceless because the trace is invariant under \(\star\).
Cayley--Hamilton gives \(Z^2=u\) on that block. Hence \(e\)
projects precisely onto \(M_2\). If the form there were
alternating, the skew adjoint subspace would be the entire
space of trace-zero two-by-two matrices. It contains the two
off-diagonal matrix units, whose commutator is nonzero.
Every skew element is half its difference from its adjoint,
and all elements of \(M_2\) occur in the span just tested.
This contradicts the commutator identity. The form is
therefore symmetric on the two-dimensional block as well.
Take the orthogonal sum with a form on the line. The
decomposition and the finitely many entries of the form
are defined over a finite extension.
\end{proof}

\begin{proposition}[A marked self-dual lattice with split reduction]
\label{prop:curvelattice366}
There is a complete coefficient discrete valuation ring \(A\)
of characteristic \(\ell\), finite over \(k[[T]]\), and a
basis for \(\rho\) over \(A\) in which the following
statements hold. The reduction is exactly
\(\operatorname{diag}(\bar\Psi,1,\bar\Psi^{-1})\).
Writing \(J\) for the three-by-three matrix with ones on the
antidiagonal and zeros elsewhere, one has
\[
       \rho(g^\star)=J\rho(g)^{\mathsf T}J,\qquad
       \rho(g_0)-\rho(g_0^\star)=\operatorname{diag}(d,0,-d),
                 \qquad d\bmod\mathfrak m_A=\delta .
\]
The ordinary flags at the chosen completions restrict to
saturated integral flags with the ordered local characters.
\end{proposition}
\begin{proof}
Compactness of the image gives a stable lattice, by taking
the integral span of the translates of any basis as in
Proposition~\ref{prop:bigweight364}. The reduction has the
stated three composition factors. Choose a composition flag
after a finite residue extension if necessary. Let \(Z\)
denote the difference of the two matrices in the statement.
In the reduction, its actions on the three simple quotients
are \(\delta,0,-\delta\). By Hensel's lemma and skew adjunction
its three eigenvalues lift to \(d,0,-d\) with unit pairwise
differences. The spectral projectors are polynomials in \(Z\)
over the coefficient ring and split the lattice as three
free lines. Choose eigenvectors in the residue labeling
\(+,0,-\). Each step of the residual composition flag is
a sum of the corresponding eigenlines, since it is stable
under \(Z\). Thus in some permutation of this basis the
residual group matrices are all upper triangular.

We give the integral adjustments explicitly. Normalize the
valuation of a uniformizer \(\pi\) here to one. Rescale the
three eigenvectors by \(\pi^{t_i}\), taking very small rational
exponents strictly increasing in the permutation just obtained,
with their spread less than one. Such scalings are available
after finite ramified extension. The valuation of the \(ij\)
entry changes by \(t_j-t_i\). Every off-diagonal unit entry
therefore acquires positive valuation, and a previously
positive valuation stays positive. The diagonal reductions
are the original characters in the residue labeling, since
the eigenvalue of \(Z\) distinguished each quotient.
We have obtained a stable lattice with split reduction.

For the symmetric form of Lemma~\ref{lem:curvesign366}, an
eigenline of \(Z\) is orthogonal to all eigenlines except
the one of opposite eigenvalue, by skew adjunction. Write
\(i^*\) for the opposite index, with \(0^*=0\), and let
\(b_i\ne0\) be the pairing of the unscaled \(i\)-th and
\(i^*\)-th vectors. We have \(b_i=b_{i^*}\). The dual of
the lattice with exponents \(t_i\) has exponents
\[
                   t_i^\#=-v(b_i)-t_{i^*}.
\]
It is also stable and has split reduction. Indeed the group
matrices on the dual basis are transposes of the matrices at
\(g^\star\), with the order of the dual lines reversed.
Consequently all their off-diagonal entries have positive
valuation. The eigenvectors with the displayed exponents
differ from that basis only by units. Replace the \(t_i\) by
\((t_i+t_i^\#)/2\). For each nonzero matrix entry its new
valuation is the average of its valuations in the two stable
lattices. Thus split reduction and stability persist, and the
new exponents satisfy \(t_i+t_{i^*}=-v(b_i)\). A further
finite ramified extension realizes these exponents. The two
outer lines now pair to a unit and the middle norm is a unit.
Rescale one outer vector by the inverse of that unit pairing
and the middle vector by an inverse square root of its norm,
after a finite extension. The form matrix is exactly \(J\).
This proves the asserted polarization and marking.

Intersect each step of a local ordinary flag over the fraction
field with this lattice. The quotient by that intersection is
torsion-free, so both the intersection and the quotient are
free and the inclusion is a direct summand over the discrete
valuation ring. Stability and the local characters descend.
On reduction the semisimplification on the decomposition group
is the restriction of the split diagonal representation.
The three characters are distinct already on residual inertia.
Their restrictions there identify their order uniquely, hence
also identify the residual characters on the entire local
group. Finally the valuation ring of a finite extension of
\(k((T))\) is finite over \(k[[T]]\). If its residue field
is \(k'\) and a uniformizer is \(z\), it is \(k'[[z]]\),
by expansion in powers of \(z\) with coefficients in the
finite field lifted by the simple roots of \(X^{|k'|}-X\).
\end{proof}

\begin{lemma}[One split prime for the local intertwining equation]
\label{lem:frobtest366}
On the common field \(L=L_a\) there is a rational prime \(p_0\)
split completely in \(L\), outside \(2sN_*w\ell\) and the
ramification of \(L\), at which
\[
      \bar\Psi(\phi)=1,\qquad \bar\kappa_\ell(\phi)=-1
\]
for every pair on the interval and every arithmetic Frobenius
lift \(\phi\) there. For any integral matrix representation
\(U\) over a local ring with this residual semisimplification,
the only matrix satisfying
\[
          \kappa_\ell(\phi)\,U(\phi)X=XU(\phi)
\]
is zero, also after arbitrary base change of the coefficient
ring. This assertion persists after twisting \(U\) by a character.
\end{lemma}
\begin{proof}
The fields \(L\) and \(E_*(\mu_\ell)\) are disjoint by
Proposition~\ref{prop:auxfield361}. Apply Chebotarev to the
element acting trivially on the former and as complex
conjugation on the latter. Exclude the finitely many
forbidden primes. The resulting prime has residue degree
one in \(L\), and the parity condition gives both displayed
values. The characters are unramified there. Modulo the
maximal ideal \(U(\phi)\) has only the eigenvalue \(1\).
Over an algebraic closure of the residue field, triangularize
this one matrix. The eigenvalues of \(X\mapsto
\kappa_\ell(\phi)U(\phi)X-XU(\phi)\) on all matrix entries
are then all \(-2\), a unit. This follows by viewing left
and right multiplication on the tensor product with the
dual space, or by triangularizing each tensor factor. The
determinant of this linear operator is a unit in the original
local ring. Its inverse therefore exists over every base
change. A scalar twist at \(\phi\) multiplies the operator
by a unit and preserves the conclusion.
\end{proof}

\begin{theorem}[A one-dimensional prime in the marked ordinary deformation ring]
\label{thm:markedcurve366}
After finite extension of coefficients the curve of
Proposition~\ref{prop:tamecurve365} defines a prime of
residue characteristic \(\ell\) with one-dimensional quotient
in the marked ordinary deformation ring of \(H\) at the split
point. Its generic representation is neither upper nor lower
triangular in the marked basis. For every chosen place its
ordinary characters are distinct in pairs even on inertia
in \(\ker\kappa_\ell\). The following further properties
hold, with \(H'=\ker(\kappa_\ell|_H)\).

\begin{enumerate}
\item The restriction to \(H'\) is absolutely irreducible,
or is an orthogonal sum of a character and an absolutely
irreducible two-dimensional representation. In the second
case the outer marked lines span the two-dimensional
constituent.
\item In every relatively open nonempty subset of \(H'\)
there is an element whose three eigenvalues are distinct.
\item The local intertwining equation of
Lemma~\ref{lem:frobtest366} has only the zero solution at
the indicated split place. On twisting the curve by
\(\bar\kappa_\ell^{-1}\), the pairing has cyclotomic
multiplier \(\bar\kappa_\ell^{-2}\) in the conjugate-dual
identity.
\end{enumerate}
\end{theorem}
\begin{proof}
Enlarge the residue field of the mixed-characteristic
coefficient ring by an unramified extension to contain \(k'\)
of Proposition~\ref{prop:curvelattice366}, using \(k'\)
as the common residue field. The marked ordinary ring here
represents matrices lifting the split triple, with the
polarization \(J\), diagonal \(Z\), and the ordered characters
of all chosen decomposition groups. For the ordinary
condition impose, for arbitrary \(h_1,h_2,h_3\) in one such
group, the equation
\[
   (\rho(h_1)-\xi_1(h_1)I)(\rho(h_2)-\xi_2(h_2)I)
                         (\rho(h_3)-\xi_3(h_3)I)=0 .
\]
An integral ordered stable flag gives this equation, since
the three factors from right to left successively lower its
steps. This is the matrix problem in
\srcloc{Development/AutomorphicArithmetic/AutomorphicVarieties/FrobeniusSmoothVectors/LeAdicClosure}{27}.

We recall its elementary representability here. The open
kernel of the split residual representation of \(H\) has
finitely generated maximal pro-\(\ell\) quotient, by the
finiteness of its continuous homomorphisms to \(\mathbf F_\ell\)
proved in Theorem~\ref{thm:commonframe365} and the pro-\(\ell\)
basis theorem. Every lift maps this kernel to a pro-\(\ell\)
principal congruence group. Hence the relevant quotient of
\(H\) is an extension of a finite group by a finitely generated
pro-\(\ell\) group. It has finitely many topological generators.
Take a power series ring on entries of matrices for these
generators relative to fixed lifts of their residues. The
matrices are invertible since their determinants have
nonzero residues. The free profinite group on the generators
acts continuously by these matrices, as the general linear
group of the power series ring is profinite. Imposing all
relators gives exactly continuous representations of
the quotient with the fixed reduction. For each local
character use its Teichm\"uller residual lift times a
principal-unit character. The latter is determined by the
maximal pro-\(\ell\) abelian quotient of the local group,
finitely generated by local reciprocity. Include one
variable per generator and the corresponding relations.
Impose the displayed word, polarization and diagonal marking
by their matrix entries. Closed ideals in this Noetherian
complete local power series ring give the resulting ring
\(R\). This presentation proves its universal property,
including for the coefficient ring \(A=k'[[z]]\).

Proposition~\ref{prop:curvelattice366} thus supplies
\(R\to A\). Its kernel is prime. Its image contains a
nonconstant local character value of residue one, by
Proposition~\ref{prop:tamecurve365}, since on the indicated
kernel the three values there were distinct powers of \(1+T\).
Subtract one to obtain a nonzero \(s_1\in zA\) in the image.
A lift of \(s_1\) belongs to the maximal ideal of \(R\),
so substitution of arbitrary convergent formal series with
coefficient lifts gives \(k'[[s_1]]\) in the image. The ring
\(A\) is finite over \(k'[[s_1]]\). Indeed the \(s_1\)-adic
and \(z\)-adic topologies agree and
\(1,z,\ldots,z^{v_z(s_1)-1}\) span \(A/s_1A\). By successive
lifting these elements generate \(A\) over \(k'[[s_1]]\).
Every intermediate algebra is therefore a finite module over
this Noetherian one-dimensional ring. In particular
\(R/\ker(R\to A)\) is an integral extension of \(k'[[s_1]]\)
and has dimension one. The same argument works for the
subalgebra generated by all matrix entries over the closed
subalgebra of the weight-zero monomials and the local
characters used in the source. The latter already contains
all the series in the chosen local character minus one.

Consider the image over an algebraic closure of the fraction
field. By Proposition~\ref{prop:tamecurve365} it is either
irreducible with connected monodromy \(\mathrm{SL}_3\), or
is a direct sum of a two-dimensional block with full
projective monodromy \(\mathrm{PGL}_2\) and a line. In the
second case the block and line are stable under the
anti-involution, and the image of \(Z^2/u\) and its kernel
identify them in the marked basis by Lemma~\ref{lem:curvesign366}.
An upper or a lower triangular representation cannot have
an irreducible constituent of dimension greater than one.
In the irreducible case there must be a nonzero matrix entry
between the middle line and the outer plane for some element,
since otherwise both subspaces would be invariant. In the
second case every such entry is zero. This identifies the
two matrix alternatives at the constructed prime.

The group \(H'\) contains all commutators of \(H\). The
Zariski closure of its image therefore contains the derived
group of the connected monodromy, by density and continuity
of the commutator morphism. In the first case this contains
\(\mathrm{SL}_3\). In the second case it contains
\(\mathrm{SL}_2\times1\) on the block and the line. In fact
the image of the connected component in \(\mathrm{PGL}_2\)
is the whole group. Its derived group maps onto that perfect
group and lies in \(\mathrm{SL}_2\times1\). A proper closed
subgroup there has dimension less than three and cannot map
onto \(\mathrm{PGL}_2\). This proves the assertions on
irreducibility and also excludes any intertwiner between
the two unequal constituents on \(H'\).

For the density assertion take a coset in \(H'\) of an open
subgroup contained in the given open set. The Zariski
closure of the image of this subgroup contains the identity
component of the closure for \(H'\), by finite index.
In the first case the coset closure therefore contains a
translate of \(\mathrm{SL}_3\), in which matrices with three
distinct eigenvalues exist. In the second it contains all
matrices of the form \((BD,\gamma)\), \(D\in\mathrm{SL}_2\),
for one invertible \(B\) and scalar \(\gamma\).
The matrix \(BD\) can be chosen diagonal with entries
\(x,\det(B)/x\), with \(x\) avoiding the finitely many values
making two among these and \(\gamma\) equal. Vanishing of the
characteristic polynomial discriminant is closed. It therefore
cannot hold on the whole coset in either case. The separation
of the three local characters was already proved in
Proposition~\ref{prop:tamecurve365}. It also excludes any
identity between one of them and another times the residual
cyclotomic character, by restriction to that same inertia
kernel. The split-prime assertion follows from
Lemma~\ref{lem:frobtest366}. Finally complex conjugation
preserves the cyclotomic character, so twisting multiplies
the two sides inside the bilinear pairing by the product of
two copies of \(\bar\kappa_\ell^{-1}\). This gives the stated
multiplier and finishes the proof.
\end{proof}

The one-dimensional prime has the required local separation.
Its symmetric pairing on every constituent gives the
polarized representation type described in
\srcloc{Development/AutomorphicArithmetic/AutomorphicVarieties/WeilGroupCharacters/UniversalOrdinaryDeterminantRing}{224}.
For each curve its generic monodromy decides which type occurs.
The propagation assertion at line~284 of that file concerns
every irreducible component through a prime of the first
type. For the second type line~293 uses chains through several
such primes. To apply these propagation assertions one must
prove that the relevant components carry the ordinary Hecke
ideal. The integral trace quotient of
Theorem~\ref{thm:mixedtrace359} still requires control on the
original fixed-weight fiber for the length estimate of
Section~\ref{sec:missing}.


\subsection{Simultaneous Frobenius tests at every parameter precision}
\label{sec:pairedtests367}

The deformation argument uses traces to test a dual cohomology
group on an ordinary curve. When the curve has two constituents,
an off-diagonal infinitesimal matrix has trace zero even after
multiplication by any block diagonal matrix. We give the trace
calculation including the needed off-diagonal deformations, and
keep track of what it detects over the finite parameter rings.

Use \(L=L_a\) from Theorem~\ref{thm:commonframe365}. Fix a
finite set of primes of \(L\) containing the primes over
\(\ell sN_*p_0\) and any tame levels in use, and stable under
complex conjugation. Include its infinite places. Write
\(G=G_{L,S}\), \(N=\ker(\kappa_\ell|_G)\) and
\(\chi=\bar\kappa_\ell|_G\). These characters factor through
\(G\). We have \(G/N\) abelian, \(\chi|_N=1\), and some
\(\gamma\in G\) has \(\chi(\gamma)=-1\) by
Lemma~\ref{lem:frobtest366}. Denote conjugation by \(c\) on
\(G\) by \(\sigma\). Pull back an ordinary curve in
Theorem~\ref{thm:markedcurve366} to \(G\), with coefficient
ring \(A=k'[[z]]\), fraction field \(Q\), and representation
\(r\). We may use \(r\) or its cyclotomic normalization since
both have the same conjugation action on matrices.

Write \(V=\bigoplus_i V_i\) over \(Q\) for the decomposition
into one or two absolutely irreducible constituents on \(N\).
The summands are already defined over \(Q\) after the allowed
finite extension, stable under \(G\), inequivalent on \(N\),
and individually polarized. Use the block lattices of
Proposition~\ref{prop:curvelattice366}. In the reducible case
these are the outer plane and the middle line. Write
\[
 E_A=\bigoplus_i\operatorname{End}_A(V_{i,A}),\qquad
 M_A=\bigoplus_{i\ne j}\operatorname{Hom}_A(V_{j,A},V_{i,A}),
\]
and \(E,M\) after extending scalars to \(Q\). Let
\(\dagger\) denote adjunction for the symmetric pairing.
It exchanges the two off-diagonal blocks, and
\(r(\sigma g)\) is a scalar multiple of
\((r(g)^\dagger)^{-1}\). A polarized off-diagonal cocycle
\(\nu:G\to M_A\) for conjugation by \(r\) satisfies
\(\nu(\sigma g)=-\nu(g)^\dagger\). Its polarized
coboundaries are \(g\mapsto r(g)Y r(g)^{-1}-Y\) with
\(Y\in M_A\), \(Y^\dagger=-Y\).

Choose continuous integral polarized cocycles
\(\nu_1,\dots,\nu_b\) independent modulo these coboundaries
over \(Q\). One may take a basis after inverting \(z\) of
any prescribed module of ordinary off-diagonal Selmer classes,
including the local conditions in
\srcloc{Development/AutomorphicArithmetic/AutomorphicVarieties/PatchingModules/PairedDualFirstOrderDualMatrixSpace}{306}.
The finiteness needed for such a basis is proved below.
In the one-block case take \(b=0\). Introduce the finite free
\(A\)-algebra
\[
 \mathcal W_A=E_A\oplus\bigoplus_{j=1}^b M_A e_j,\qquad
 (M_A e_i)(M_A e_j)=0 ,
\]
with left and right multiplication by block diagonal matrices.
There is a representation into its units given by
\[
     R(g)=\Big(1+\sum_j\nu_j(g)e_j\Big)r(g).
\]
Indeed the cocycle identity is exactly the multiplication
rule in this square-zero extension. Write
\(\mathcal W=\mathcal W_A\otimes_A Q\).

\begin{lemma}[The algebra spanned on the cyclotomic kernel]
\label{lem:spandual367}
The \(Q\)-linear span of \(R(N)\) is \(\mathcal W\).
\end{lemma}
\begin{proof}
The span \(D\) is a subalgebra, since \(R\) is a representation.
Modulo the square-zero ideal its image is \(E\). Indeed the
summands on \(N\) are absolutely irreducible and pairwise
inequivalent. Over an algebraic closure their semisimple
image algebra acts faithfully and is the product of the
full matrix algebras on the distinct simple modules. Equality
then follows over \(Q\) by comparing dimensions. In the
one-block case this proves the lemma.

In the two-block case \(D\cap M^{\oplus b}\) is an \(E\)-bimodule.
To multiply by any \(x\in E\), lift \(x\) to \(D\). The result
on the intersection is independent of the lift because the
ideal has square zero. The two central block idempotents
separate the two directions of the off-diagonal blocks.
In a given direction \(ij\), left and right matrix units
move each single entry to any chosen position. Thus the
intersection in that direction consists of all matrices
of that block size with entries in one fixed subspace
\(U_{ij}\subset Q^b\). If \(D\) is proper then for at least
one direction there is a nonzero functional
\(\lambda=(\lambda_1,\ldots,\lambda_b)\) killing \(U_{ij}\).

Project the algebra to \(E\oplus M_{ij}\) by keeping that
direction, applying \(\lambda\), and killing the opposite
direction. This is an algebra map because both radicals
are square-zero. The projected image of \(D\) maps
isomorphically to \(E\). It is therefore the graph of a
derivation \(d:E\to M_{ij}\). Let \(P_i\) be the idempotent
of block \(i\). In \(M_{ij}\) it acts as one on the left
and as zero on the right. Differentiating \(P_i x=xP_i\)
gives
\[
                     d(x)=x d(P_i)-d(P_i)x .
\]
For \(n\in N\) the equality describing the graph reads
\(\big(\sum_t\lambda_t\nu_t(n)_{ij}\big)r(n)=d(r(n))\),
viewing that block as a matrix with the other entries zero.
Thus the indicated
linear combination in that block direction is a coboundary
on \(N\). It is a coboundary on \(G\) as well. After subtracting
this coboundary a cocycle vanishes on \(N\), and its value
at any \(g\in G\) is \(N\)-invariant by the cocycle identity
and normality of \(N\). The space \(M_{ij}^N\) is zero by
inequivalence and irreducibility.

Write \(Y_{ij}\) for this coboundary witness over \(G\).
Polarization determines the opposite block uniquely and
gives the witness \(-Y_{ij}^\dagger\) there. This follows
by applying \(-\dagger\) to the coboundary at \(\sigma g\)
and using the scalar relation between \(r(\sigma g)\)
and \((r(g)^\dagger)^{-1}\). Hence the whole combination
\(\sum_j\lambda_j\nu_j\) is a polarized coboundary, contrary
to independence. Both \(U_{ij}\) must be full, which proves
the lemma.
\end{proof}

Let \(\mathcal W_A^\vee=\operatorname{Hom}_A(\mathcal W_A,A)\)
with bimodule structure \((afb)(x)=f(bxa)\), and twist its
conjugation action by \(\chi\). For a continuous cocycle
\(y:G\to\mathcal W_A^\vee(1)\) and \(n\in N\) put
\[
                   \lambda_n(y)=y(n)(R(n)).
\]
This is an ordinary trace on diagonal matrices and includes
the paired off-diagonal terms. In coordinates a functional
on the diagonal algebra is trace against a block diagonal
matrix \(P\). A functional on the \(j\)-th summand of the
radical is trace against an off-diagonal matrix \(D_j\) in
the opposite positions. Evaluation on \(R(n)\) gives the
sum of \(\operatorname{tr}(P r(n))\) and
\(\sum_j\operatorname{tr}(D_j\nu_j(n)r(n))\), allowing
\(P,D_j\) to depend on \(n\). This identifies the functional
with the trace coefficient used in the source. It vanishes
on coboundaries since \(\chi(n)=1\) and conjugation by
\(R(n)\) fixes \(R(n)\).

\begin{theorem}[Detection of the twisted dual by evaluations]
\label{thm:derivation367}
Let \(W\) be any finite-dimensional unital algebra over a
field \(Q\), \(R:G\to W^\times\) a representation, and
\(N\triangleleft G\) such that \(R(N)\) spans \(W\).
Assume \(G/N\) is abelian and a character
\(\chi:G/N\to Q^\times\) takes a value different from one.
Give \(W^\vee\) the action \(g\cdot f=\chi(g)R(g)fR(g)^{-1}\).
If a cocycle \(y\) satisfies \(y(n)(R(n))=0\) for every
\(n\in N\), then \(y\) is a coboundary.
\end{theorem}
\begin{proof}
On the group algebra \(Q[N]\) extend \(R\) linearly and put
\(d(n)=y(n)R(n)\), extending this map linearly too. Since
\(\chi|_N=1\), the cocycle law gives the Leibniz rule with
respect to \(R\). Write \(\tau(f)=f(1)\). The hypothesis is
\(\tau(d(x))=0\) on all of \(Q[N]\). If \(R(x)=0\), then
for any \(v\in Q[N]\) the Leibniz rule gives
\[
              0=\tau(d(xv))=\tau(d(x)R(v))=d(x)(R(v)).
\]
Since \(R(Q[N])=W\), this forces \(d(x)=0\). The map \(d\)
thus descends to a derivation \(D:W\to W^\vee\).

Choose \(\gamma\in G\) with \(\chi(\gamma)\ne1\) and set
\(B=R(\gamma)\). For \(X=R(n)\), compute
\(D(BXB^{-1})\) in two ways. The derivation rule in \(W\)
and \(D(B^{-1})=-B^{-1}D(B)B^{-1}\) give
\[
             B D(X) B^{-1}+[D(B)B^{-1},BXB^{-1}].
\]
Using the cocycle at \(\gamma n\gamma^{-1}\) instead gives
\[
             \chi(\gamma)B D(X)B^{-1}+[y(\gamma),BXB^{-1}].
\]
Here the commutator uses the two bimodule actions.
Both expressions are linear in \(X\), hence agree on all
of \(W\). Conjugating back and dividing by
\(1-\chi(\gamma)\) shows that \(D(X)=[Y,X]\) for the fixed
functional
\[
          Y=\frac{B^{-1}(y(\gamma)-D(B)B^{-1})B}
                        {1-\chi(\gamma)} .
\]
The restriction of \(y\) to \(N\) is therefore the
coboundary with witness \(-Y\).

Subtract the corresponding \(\chi\)-twisted coboundary on
\(G\). The remaining cocycle factors through \(G/N\) on the
invariants of \(N\), by the argument of
Lemma~\ref{lem:spandual367}. These are the functionals
commuting with all of \(W\), since \(R(N)\) spans it.
On this space the quotient acts through the scalar \(\chi\).
For a cocycle \(f\) on the abelian quotient, commutativity
now gives
\[
   (\chi(\gamma)-1) f(g)=(\chi(g)-1) f(\gamma).
\]
This gives the final coboundary witness
\(f(\gamma)/(\chi(\gamma)-1)\) and proves the assertion.
\end{proof}

\begin{proposition}[One loss for all finite parameter rings]
\label{prop:loss367}
On any of the curves above, let \(N^{\rm reg}\) be any
relatively open dense subset of \(N\). There exist finitely
many \(n_1,\ldots,n_q\in N^{\rm reg}\) and an integer \(c_1\ge0\)
such that for every \(m\ge1\), a continuous class over
\(\mathcal W_A^\vee(1)/z^m\) whose evaluations vanish at these
elements is killed by \(z^{c_1}\). This includes all classes
over \(A/z^m\), whether or not they come from classes over
\(A\). For the polarized cohomology the coboundary witness
can be taken with the required symmetry.
\end{proposition}
\begin{proof}
Here are the finiteness details. For any finite free
\(A\)-module \(P\) with continuous action of \(G=G_{L,S}\),
the module of continuous cocycles with values in \(P\)
is finitely generated over \(A\). More precisely, the
cocycles over \(A\) and over every \(A/z^m\) are given
by one finite matrix of linear relations on one finite
list of generator values. Choose the open normal subgroup
\(U\) acting trivially on \(P/zP\). It has finite-dimensional
\(H^1(U,\mathbf F_\ell)\). Indeed it is a Galois group
with ramification allowed at a finite set over the finite
field extension of \(L\) corresponding to \(U\). After
adjoining \(\mu_\ell\), Kummer theory bounds the dimension
using the units and the class group of the ring of \(S\)-integers
as in Theorem~\ref{thm:commonframe365}. It suffices here
even to inflate to that restricted ramification group.
Consequently the maximal pro-\(\ell\) quotient of \(U\)
is topologically finitely generated. Every cocycle together
with the action gives a homomorphism to the semidirect
product of the additive module with the automorphism group.
On \(U\) its image is pro-\(\ell\), since the linear part
is in a principal congruence group and the additive part
is pro-\(\ell\). The same holds modulo every \(z^m\).
All these cocycles thus factor, with their action, through
one topologically finitely generated quotient of \(G\).

Let \(l_0\) be its number of chosen generators and
\(\mathcal F\) the free profinite group on them. Over \(A\)
any assignment of values in \(P^{l_0}\) extends to a unique
continuous cocycle on \(\mathcal F\), by the semidirect
product and its profiniteness. Evaluation at any word of
\(\mathcal F\) is an \(A\)-linear map of these variables.
For a relator it must be zero to descend. The rows of all
such maps generate a submodule of the dual of \(P^{l_0}\),
so finitely many suffice by Noetherianity. Write \(C\)
for the resulting matrix. The same rows work modulo
\(z^m\), since every tuple of generator values lifts freely
before relations are imposed. This proves the claimed
presentation. It also proves the earlier assertion on the
choice of \(\nu_j\): imposed linear local and polarization
conditions give a submodule of a finite module. The quotient
by coboundaries in it is finite over \(A\), and clearing
denominators gives representatives of a basis over \(Q\).

Apply this presentation with \(P=\mathcal W_A^\vee(1)\),
and write \(B_0:P\to P^{l_0}\) for evaluation of coboundaries
on the chosen generators. The vector space
\((\ker C\otimes_A Q)/\operatorname{im}(B_0\otimes_A Q)\)
is finite-dimensional. Lemma~\ref{lem:spandual367} and
Theorem~\ref{thm:derivation367} show that all the evaluations
on \(N\) separate it. They still separate it on \(N^{\rm reg}\)
by continuity, after multiplying a fractional cocycle by a
power of \(z\) to make it integral. Finite-dimensional
linear algebra selects \(n_1,\ldots,n_q\) with the same
property. Fix words above these elements and append their
evaluation rows to \(C\), obtaining a matrix \(D_0\).
Coboundaries satisfy these rows. We have proved
\[
          \ker(D_0\otimes_A Q)=\operatorname{im}(B_0\otimes_A Q).
\]

Let \(b_1\) be the largest valuation of a nonzero Smith
invariant of \(D_0\), taking zero if the matrix is zero.
The finitely generated torsion module
\(\ker D_0/\operatorname{im} B_0\) is killed by some
\(z^{b_2}\). If \(D_0v=0\bmod z^m\), then
\(z^{b_1}v\) belongs to \(\ker D_0+z^m P^{l_0}\).
Indeed in Smith coordinates every component outside the
kernel is sent to zero modulo \(z^m\) after this
multiplication. Hence
\(z^{b_1+b_2}v\in\operatorname{im}B_0+z^m P^{l_0}\).
A cocycle modulo \(z^m\) is determined by its generator
values, giving the result with \(c_1=b_1+b_2\).

For the polarization, the signed adjunction on coefficients
intertwines the \(G\)-action with its conjugate by \(\sigma\).
Explicitly adjunction on \(\mathcal W_A\) acts by \(\dagger\)
on every block and every \(e_j\)-coefficient. The identities
for \(r\) and \(\nu_j\) give
\(R(\sigma g)\) equal to a scalar times \((R(g)^\dagger)^{-1}\).
On the dual use \(f\mapsto-\chi(c)\,f\circ\dagger\).
It is an involution and intertwines the twisted actions
because \(\chi(\sigma g)=\chi(g)\). Thus when a cocycle
has the prescribed symmetry, the average of any coboundary
witness and its image under this involution is a witness
with that symmetry. This averaging is integral since
\(\ell\) is odd.
\end{proof}

\begin{theorem}[Common primes for the dual trace tests]
\label{thm:primes367}
Fix the interval and the coefficient prime \(\ell\). For
any finite collection of the ordinary curves of
Theorem~\ref{thm:markedcurve366} on their common field,
one can choose a single number \(q\) and a single loss \(c_1\).
For every \(m,j\ge1\) there are \(q\) distinct rational primes
outside the fixed levels, split completely in \(L\) and
congruent to one modulo \(\ell^j\), with choices of
Frobenius above them such that the following holds on every
curve. All their characteristic polynomials are separable over
the fraction field. If a continuous class with coefficients
in \(\mathcal W_A^\vee(1)/z^m\) has zero evaluations at these
Frobenius elements, multiplication by \(z^{c_1}\) kills it.
The collection may include a curve for every \(r,\alpha\)
on the entire interval.
\end{theorem}
\begin{proof}
Work over the common absolute group, or its quotient \(G_{L,S}\)
for the union of the finitely many sets of places. The
regular-semisimple locus of each curve is relatively open
and dense on \(N\) by Theorem~\ref{thm:markedcurve366}.
Pullback to the common group preserves density because the
projection is open and maps its cyclotomic kernel onto the
individual kernel. The finite intersection of these open
dense sets is again dense. Use it for \(N^{\rm reg}\)
in Proposition~\ref{prop:loss367} on each curve, take the union
of the resulting finite lists of elements and the maximum of
the losses. There are finitely many possible pairs on the
interval since \(r<t_*\) and all characters factor through
the one finite extension \(E_*\).

Fix \(m,j\). Over each finite parameter ring there are
only finitely many cocycles, by the finite list of generator
values in the preceding proof. For each selected element
prescribe its image in the joint finite quotient recording
the matrices \(R\) and all these cocycles modulo \(z^m\)
on every curve. Such a quotient exists by continuity of each
semidirect-product representation. Record also the matrices
\(r\) to a precision exceeding all the valuations of the
nonzero characteristic discriminants at the selected
elements. Lift the elements to \(G_L\). Their full
cyclotomic character equals one, so these prescriptions
together with trivial action on \(\mu_{\ell^j}\) hold in
a finite Galois extension of \(L\). Pass to its normal
closure over \(\QQ\), including \(L\). Chebotarev supplies
infinitely many rational primes in the conjugacy class of
each prescribed element. At such a prime a place of the
normal closure can be chosen so that its Frobenius is the
specified element itself. It fixes \(L\) and \(\mu_{\ell^j}\),
giving both splitting and the congruence. Choose distinct
primes and discard all ramified primes and primes at the
fixed levels. The joint finite images make every evaluation
equal to the prescribed one modulo \(z^m\), and the higher
precision keeps all the discriminants nonzero. At these
Frobenius elements \(\chi=1\), so evaluation again depends
only on the class. Proposition~\ref{prop:loss367} applies
with the same \(c_1\).
\end{proof}

The trace coefficient tested here is the one in
\srcloc{Development/AutomorphicArithmetic/AutomorphicVarieties/PatchingModules/ThetaCoefficientMapEq}{254}.
The loss is independent of both precisions \(m,j\). It is
measured in the parameter \(z\) on a curve of characteristic
\(\ell\), whereas the lengths in Section~\ref{sec:missing}
are measured in a uniformizer over \(\ZZ_\ell\) at the
original weight. A comparison with those lengths still
requires the action of the Hecke algebra on the modules at
the chosen levels and its behavior on specializing the
weights. Theorem~\ref{thm:primes367} provides the common
primes for that argument.


\subsection{The first-order Selmer group at the auxiliary levels}
\label{sec:selmerlength368}

Fix \(F=F_a^+\), \(L=KF\), and the interval of
Theorem~\ref{thm:commonframe365}. Write \(d=[F:\QQ]\) and
\(n=3\). Take one of the curves of Theorem~\ref{thm:markedcurve366}
with \(A=k'[[z]]\) and \(Q=k'((z))\). Let \(E_A,M_A,\nu_j,R\)
and \(W=\mathcal W_A\) be as in Section~\ref{sec:pairedtests367}.
Choose the \(\nu_j\) on \(\Gamma_a^0\) from the ordinary
off-diagonal Selmer group, with integral local witnesses as in
that section. Write \(n_i\) for the ranks of the constituent
blocks. The involution \(\dagger\) on \(W\) preserves the labels
\(e_j\). The action of \(G_L\) by conjugation by \(R\) extends to
\(G_F\) by letting a complex conjugation \(c\) act as
\(-\dagger\). Indeed the polarization gives
\(R(cgc^{-1})\) equal to \((R(g)^\dagger)^{-1}\) times the
scalar multiplier of \(r\). This verifies the conjugation
relation, and \(c^2=1\). Write \(A_m=A/(z^m)\) and \(P_m=P/z^mP\)
for any of our free coefficient modules \(P\). For finite
groups of coefficients or cohomology use \(h\) for length over
\(A\), adding an exponent to indicate a cohomological degree.
Length is the logarithm of cardinality to the base \(|k'|\).

At each \(u\mid\ell\) use one of the two places of \(L\) over
\(u\) to identify its local group with \(G_{F_u}\). In a basis
adapted to the ordinary flag, take the upper triangular entries
in \(E_A\) and in every \(M_Ae_j\), using the order of the
three successive characters on the flag. Such a basis can be
taken in the individual blocks. A local inertia element with
distinct residual eigenvalues has its three eigenprojectors
over \(A\) by Hensel's lemma. Each lies in a constituent block,
and a saturated stable step of the flag is the sum of the
eigenlines that it contains. If \(Y_{u,j}\) are the local
witnesses for \(\nu_j\), conjugate \(R\) by
\(1+\sum_jY_{u,j}e_j\). This replaces \(\nu_j\) by
\(\nu_j-(\operatorname{Ad}(r)Y_{u,j}-Y_{u,j})\), which is
upper triangular on the local group. In these coordinates
the upper triangular subspace is stable under \(R\).
Transport it back and call it \(B_u\subset W\). Both it
and its quotient are free. Define
\[
 w=\sum_i n_i^2+2b\sum_{i<j}n_i n_j,\qquad
 u_0=(w+n)/2,\qquad w_-=(w-n)/2 .
\]

\begin{lemma}[Local and global invariant spaces]
\label{lem:localranks368}
The rank of \(B_u\) is \(u_0\). Uniformly in \(m\ge1\), the
lengths of \(H^0(F,W_m)\), \(H^0(F_u,(W/B_u)_m)\) and
\(H^0(F_u,B_{u,m}^{\vee}(1))\) are bounded. Here the dual
is the \(A_m\)-linear dual with the residual cyclotomic twist.
At a real place \(\tau\) one has \(h^0(F_\tau,W_m)=w_-m\).
\end{lemma}
\begin{proof}
Each diagonal block contributes \(n_i(n_i+1)/2\) allowed entries
and each unordered pair of distinct blocks contributes \(n_i n_j\)
in each radical copy. These sum to \(u_0\).
On the radical the local action is ordinary triangular
conjugation by \(r\), and on the quotient by the radical it
is block diagonal conjugation by \(r\). Filter these triangular
matrix spaces by matrix units in order of row increasing and,
within each row, column decreasing. Conjugation takes a
matrix unit at \(ij\) to units at \(i'j'\) with \(i'\le i\)
and \(j'\ge j\), with coefficient \(\xi_i/\xi_j\) at \(ij\)
itself. Here \(\xi_i\) is the \(i\)-th ordinary character.
Consequently the factors of \((W/B_u)\otimes Q\) all have
\(i>j\), and those of \(B_u^\vee(1)\otimes Q\) have characters
\(\bar\kappa_\ell\xi_j/\xi_i\) with \(i\le j\).
At distinct indices these characters are nontrivial on the
kernel of full cyclotomic inertia by
Theorem~\ref{thm:markedcurve366}. At equal indices
\(\bar\kappa_\ell\) is nontrivial because \(F_u\) does not
contain \(\mu_\ell\) by Theorem~\ref{thm:commonframe365}.
The two local modules therefore have zero invariants over \(Q\).

A global invariant over \(Q\) commutes with \(R(G_L)\).
Lemma~\ref{lem:spandual367} puts it in the center of \(W\otimes Q\).
Commuting with the block idempotents kills its off-diagonal
entries. Its diagonal blocks are scalars since they commute
with all matrix units in their blocks. The operation \(c\)
negates those scalars, and \(\ell\) is odd. Thus this invariant
is zero. Whenever a finite free module has zero invariants over
\(Q\), finitely many operators \(g-1\) already have trivial
common kernel there. Stack their matrices over \(A\).
Smith normal form bounds the length of its kernel over \(A_m\)
by the sum of the valuations of its nonzero invariant factors,
independently of \(m\). This proves all three bounds.

For the real formula write \(c_\tau=hc\), \(h\in G_L\), so that
\(chc^{-1}=h^{-1}\). The cyclotomic value of \(h\) is one, as
both conjugations have cyclotomic value \(-1\). The polarization
therefore gives \(R(h)^\dagger=R(h)\). On each diagonal block
modulo the radical the action of \(c_\tau\) is negative
adjunction for a nondegenerate symmetric form. Its fixed space
is the space of skew-adjoint matrices and has rank
\(n_i(n_i-1)/2\), as is seen by multiplying by the form matrix.
On each radical copy \(c_\tau\) exchanges opposite blocks by
mutually inverse linear maps. The fixed space has rank
\(\sum_{i<j}n_i n_j\). Taking invariants is exact over \(A\)
and over \(A_m\) by the idempotent \((1+c_\tau)/2\).
The ranks sum to \(w_-\), proving the claim.
\end{proof}

Fix a finite set \(S\) over \(F\) containing the real places,
the places over \(\ell\), the ramified places of our coefficients
and the places over \(p_0\) of Lemma~\ref{lem:frobtest366}.
Include all earlier levels in use. Write \(n_1',\ldots,n_q'\)
for the group elements used by Theorem~\ref{thm:primes367},
to distinguish them from the block sizes. At stage \(N\ge1\)
choose \(q\) distinct rational primes split completely in \(L\),
all congruent to one modulo \(\ell^N\), with chosen places
giving the prescribed Frobenius elements to precision \(z^N\).
In making that Chebotarev choice, record the cocycles at
\emph{every} precision \(1\le e\le N\) in the finite quotient
in its proof. There are finitely many at each precision and
finitely many curves on this interval. This makes all the tests
valid even for classes modulo \(z^e\) that do not lift modulo
\(z^N\). Write \(\mathcal Q_N\) for the chosen places of \(F\),
one for each rational prime.

On \(G_{F,S\cup\mathcal Q_N}\) take the local conditions
\(\mathcal L\) on \(W_m\) given by the image of \(H^1(F_u,B_{u,m})\)
at \(u\mid\ell\), the full group at \(\mathcal Q_N\) and above
\(p_0\), and the unramified group at every other finite place.
At the real places \(H^1=0\). Denote the orthogonal conditions
for the dual \(W_m^\vee(1)\) by \(\mathcal L^\perp\).

\begin{proposition}[Bounds at the selected Frobenius elements]
\label{prop:testbounds368}
For every \(1\le m\le N\) the length of
\(H^1_{\mathcal L^\perp}(F,W_m^\vee(1))\) is bounded by a
constant independent of \(m,N\). At \(v\in\mathcal Q_N\)
one has \(h^0(F_v,W_m^\vee(1))=nm+O(1)\) uniformly in these
parameters. At \(v\mid p_0\) this invariant group vanishes.
\end{proposition}
\begin{proof}
A dual class is locally zero at each new place. The coefficients
are unramified there, so every representing cocycle then vanishes
on the new inertia. It descends to \(G_{F,S}\), since it kills
the closed normal subgroup generated by that inertia on which
the action is trivial. Restrict the descended class to \(G_L\).
At every chosen Frobenius its evaluation of
Section~\ref{sec:pairedtests367} vanishes. Indeed it is locally
a coboundary and the twist is one at that Frobenius.
Theorem~\ref{thm:primes367} consequently kills the restricted
class by \(z^{c_1}\). Restriction has trivial kernel since
corestriction composed with it is multiplication by two.

The proof of Proposition~\ref{prop:loss367}, applied to
\(G_{F,S}\), bounds the number of generator values for its
cocycles independently of \(m\). Its proof uses finiteness
of the \(S\)-unit rank and class group and works over \(F\)
as over \(L\). Hence the dual group in question is a
subquotient of \(A_m^{\,D}\) for a fixed \(D\). Any such
subquotient has at most \(D\) generators, by lifting
submodules to \(A^D\) and using that \(A\) is a discrete
valuation ring. Being killed by \(z^{c_1}\), it has length
at most \(Dc_1\).

At the \(j\)-th new place the action on \(W_m\) equals
conjugation by \(R(n_j')\) modulo \(z^m\) and the twist is
trivial. Over \(Q\), the matrices \(r_i(n_j')\) in the
constituent blocks have simple and mutually disjoint spectra.
For different blocks \(i,k\) the operator
\(X\mapsto r_i(n_j')X-Xr_k(n_j')\) is invertible, as one
sees over a splitting field. Conjugation by an element of
one plus the square-zero ideal thus removes the radical
part of \(R(n_j')\). Its centralizer has zero radical part
and diagonal part of dimension \(\sum_i n_i=n\), since the
centralizer of a matrix with simple spectrum has dimension
equal to its size. The dual conjugation operator has the same
fixed dimension by transposition and invertibility.
Smith normal form of this fixed operator minus the identity
gives a kernel of length \(nm+\sum_j\min(m,a_j)\), where
the \(a_j\) are the valuations of its nonzero Smith factors.
The list is fixed before \(N\) varies. At a place above
\(p_0\), the residual conjugation has only eigenvalue one,
both on the radical and on the quotient, by
Lemma~\ref{lem:frobtest366}. Twisting the dual by \(-1\)
there makes Frobenius minus one invertible modulo \(z\),
and therefore over every \(A_m\).
\end{proof}

\begin{theorem}[First-order Selmer length at all auxiliary levels]
\label{thm:selmerlength368}
There is a constant \(C\) such that for every \(1\le m\le N\)
\[
 \left|h\bigl(H^1_{\mathcal L}(F,W_m)\bigr)
                 - n(d+q)m\right|\le C .
\]
For a fixed interval and coefficient prime the same field,
the same \(\mathcal Q_N\), and one constant serve all the
curves selected for \(x,r,\alpha\) on the interval.
\end{theorem}
\begin{proof}
To identify the finite Cartier dual with \(W_m^\vee(1)\), compose the \(A_m\)-valued pairing with extraction of the
coefficient of \(z^{m-1}\) followed by
\(\operatorname{Tr}_{k'/\mathbf F_\ell}\). This pairing is
perfect since a nonzero coefficient can be shifted to degree
\(m-1\) and paired by the nondegenerate finite-field trace.
The Poitou--Tate dimension formula in the proof of
Proposition~\ref{prop:strictclasses365} applies with lengths,
by taking the logarithm to base \(|k'|\) of its order formula.
The global term \(h^0(F,W_m)-h^0(F,W_m^\vee(1))\) is bounded
by Lemma~\ref{lem:localranks368} and the vanishing at \(p_0\).
The dual Selmer term is bounded by
Proposition~\ref{prop:testbounds368}.~\cite{VII-V-V-milne98}

At \(u\mid\ell\) the cohomology sequence for \(B_{u,m}\subset W_m\)
gives exactly
\[
 h(\mathcal L_u)-h^0(F_u,W_m)
 =h^1(F_u,B_{u,m})-h^0(F_u,B_{u,m})
                    -h^0(F_u,(W/B_u)_m).
\]
The local Euler characteristic and local Tate duality identify
the right side with \([F_u:\QQ_\ell]u_0m+O(1)\) by
Lemma~\ref{lem:localranks368}. At an old finite place with the
unramified condition the contribution is zero. In fact the
unramified classes are the cokernel of Frobenius minus one on
the inertia invariants, whose length equals the length of
the kernel. At a place away from \(\ell\) with the full condition
the contribution is \(h^2=h^0(F_v,W_m^\vee(1))\), by the
local Euler characteristic and duality. This is zero above
\(p_0\) and sums to \(nqm+O(1)\) at the new places. At each
real place the contribution is \(-w_-m\). Since
\(\sum_{u\mid\ell}[F_u:\QQ_\ell]=d\) and \(u_0-w_-=n\),
the claimed estimate follows. All constants arose from fixed
matrices, fixed global groups, and finitely many selected
elements. Take their maximum over the finite list of curves.
\end{proof}

\begin{corollary}[The first numerical inequality for patching]
\label{cor:bfirst368}
The first-order polarized cocycles of the ordinary problem at
these levels, modulo its polarized coboundaries, have order at
most \(|k'|^{\,n(d+q)m+C}\) for \(1\le m\le N\).
\end{corollary}
\begin{proof}
A polarized cocycle on \(G_L\) extends to \(G_F\) with value
zero at \(c\), by the defining relation with \(-\dagger\).
A coboundary between two such representatives has a witness
fixed by \(c\), exactly the required skew-adjoint witness.
Inflation from the prescribed tame quotient is injective on
classes. At an old finite prime its allowed inertia has order
prime to \(\ell\), so the restriction in \(H^1\) is zero by
averaging. The ordinary witness at \(u\) puts the local class
in the image from \(B_{u,m}\). Hence this quotient of polarized
cocycles injects into the group of Theorem~\ref{thm:selmerlength368}.
This proves the bound in the normalization of
\srcloc{Development/AutomorphicArithmetic/AutomorphicVarieties/PatchingModules/ThetaCoefficientMapEq}{54}.
\end{proof}

The length parameter in this theorem is the precision along
the ordinary curve in characteristic \(\ell\). The weighted
sum in Proposition~\ref{prop:res23} is taken in mixed
characteristic at the original weight, as \(\ell\) and the
period fields vary. Passing between these estimates still
requires control of the specialized integral trace relations.


\subsection{The obstruction to lifting ordinary cocycles}
\label{sec:selmerlift369}

Keep the fields, curves, local flags and auxiliary sets
\(\mathcal Q_N\) of Section~\ref{sec:selmerlength368}. The
second numerical input for patching concerns the off-diagonal
module \(M=M_A\) of that section, with its conjugation action
and with \(c\) acting as negative adjunction. Write \(E=E_A\).
Let \(b\) be the rank of the integral ordinary Selmer group
for \(M\) on the original quotient \(\Gamma_a\), and use a
basis after inverting \(z\), represented by the \(\nu_1,\ldots,\nu_b\).
In each ordinary basis put \(B_u(M)=M\cap\{\text{upper
triangular matrices}\}\). Use the corresponding submodule
of \(E\) for \(B_u(E)\). For these modules define the
ordinary Selmer group at level \(N\) by the same local
conditions as for \(W_m\) in Section~\ref{sec:selmerlength368},
using their own upper triangular submodules.

In this situation the diagonal requirement on local cocycles
in the source has a simple interpretation. After subtracting
a local coboundary, an upper triangular cocycle for an upper
triangular representation has additive characters on its
diagonal. Indeed conjugating an upper triangular matrix by
another invertible upper triangular matrix fixes its diagonal,
over any commutative coefficient ring. Each diagonal character
therefore kills the kernel of topological abelianization.
By local reciprocity the abelianization of \(G_{F_u}\) is
the profinite completion of \(F_u^\times\). Its
\(\ell\)-primary torsion is trivial since \(F_u\) has no
nontrivial \(\ell\)-power roots of unity, as in
Theorem~\ref{thm:commonframe365}. Its valuation factor
\(\widehat{\ZZ}\) is torsion free. Thus the diagonal kills
the whole preimage of the \(\ell\)-primary torsion used in
\srcloc{Development/AutomorphicArithmetic/AutomorphicVarieties/WeilGroupCharacters/UniversalOrdinaryDeterminantRing}{75}.
This applies also to the commutative coefficient algebra
used below.

\begin{lemma}[The off-diagonal Selmer length]
\label{lem:odbase369}
There is a constant \(C_0\) such that, for \(1\le e\le N\),
the ordinary Selmer group with coefficients in \(M_e=M/z^eM\)
at level \(N\) has length at most \(be+C_0\).
\end{lemma}
\begin{proof}
At the level without the new primes, the classes are the
ordinary classes on \(\Gamma_a\). To verify this, at an old
finite place different from \(\ell\) where unramified
cohomology is imposed, restriction to inertia is a coboundary.
Every representative thus vanishes on the subgroup of inertia
killed in \(\Gamma_a\), since that subgroup acts trivially.
Conversely the remaining inertia in \(\Gamma_a\) has order
prime to \(\ell\) and has zero first cohomology by averaging.
At \(p_0\) the full group equals the unramified group.
Their difference in length is \(h^0(F_v,M_e^\vee(1))=0\),
by local Euler characteristic and the Frobenius calculation
in Proposition~\ref{prop:testbounds368} on the off-diagonal
blocks. These facts give the asserted identification by
inflation. Normalizing a representative to be zero at \(c\)
identifies this cohomology with polarized cocycles modulo
polarized coboundaries, since two is invertible.

We give the length calculation at this fixed level. Use the
finite presentation of cocycles by values on generators in
Proposition~\ref{prop:loss367}, applied to the quotient
\(\Gamma_a\). Its open subgroups have the stated finiteness
property by Theorem~\ref{thm:commonframe365}. Adjoin a variable
\(Y_u\in M/B_u(M)\) at each \(u\mid\ell\). For each local
element \(g\), impose that the projection of the global
cocycle at \(g\) to this quotient equal \((g-1)Y_u\).
All these equations and the global cocycle relations are
rows on one finite free \(A\)-module \(A^t\). Finitely many
generate the row module and can be arranged in a matrix \(D\).
Before imposing the relations the assignments on the free
profinite group of generators are arbitrary. Thus the same
matrix describes the solutions over \(A_e\), including
solutions that have no integral lift. Coboundaries with
their projected local witnesses are the image of a fixed
matrix \(D'\) with \(DD'=0\).

For each \(u\), the module \((M/B_u(M))\otimes Q\) has zero
invariants. All its character factors come from distinct
lines of the ordinary flag by the argument of
Lemma~\ref{lem:localranks368}. The integral witnesses
are therefore unique once the global cocycle is fixed.
It follows from the definition of \(b\) that
\(t-\operatorname{rank}_Q D-\operatorname{rank}_Q D'=b\).
Over \(A_e\), forgetting the witnesses maps
\(\ker D_e/\operatorname{im}D'_e\) onto the ordinary Selmer
group. Smith normal form gives length \(be+O(1)\) for this
quotient of matrices, since each nonzero invariant factor of
either matrix adds at most its fixed valuation to \(be\).

At level \(N\) the kernel of restriction to the quotients by
unramified cohomology at the new primes is precisely the
group just counted. Indeed \(M\) is unramified at those primes,
so vanishing of the inertia restriction makes the cocycle
descend to the original group. The length of each local
quotient is \(h^0(F_v,M_e^\vee(1))\), by the same local
Euler formula. Frobenius there acts modulo \(z^e\) as at the
fixed element \(n_j'\), with trivial twist. Its blocks have
mutually disjoint spectra over a splitting field of \(Q\).
There are no intertwiners between different blocks, in either
direction. Frobenius minus one on the off-diagonal dual thus
has full rank over \(Q\), so its kernel length modulo \(z^e\)
is bounded by a constant from its Smith factors. Summing
over the fixed number \(q\) of new places proves the lemma.
\end{proof}

To retain the witnesses in the lifting problem, use the
following cochain complex. For coefficients \(P\) at precision
\(m\), take continuous global cochains on
\(G_{F,S\cup\mathcal Q_N}\). Map them by localization to the
sum of the local complexes displayed below.
\[
 \begin{array}{ll}
 C^\bullet(F_u,P/B_u(P)) & (u\mid\ell),\\[2pt]
 C^\bullet(F_v,P)\big/
 C^\bullet(G_{F_v}/I_v,P^{I_v})
       & (v\text{ old, with unramified condition}).
 \end{array}
\]
The denominator maps in by inflation. At full places the
target is zero. Take the mapping fiber of this map and denote
its cohomology by \(\mathcal H^i_N(P)\). In degree one a
cocycle consists of a global cocycle and local zero-cochains
whose coboundaries are its images in the target. Thus its
image in global cohomology is the ordinary Selmer group.

\begin{proposition}[Reduction of the obstruction group]
\label{prop:obsquot369}
Put \(X_{N,m}=\mathcal H^2_N(M_m)\). There is \(C_1\ge0\)
such that for \(1\le e\le m\le N\)
\[
 X_{N,m}/z^eX_{N,m}\ \simeq\ X_{N,e},
 \qquad h(X_{N,e})\le be+C_1 .
\]
\end{proposition}
\begin{proof}
All the cohomology groups here are finite. For the global and
local groups this is arithmetic cohomological finiteness,
and then it follows for the fiber by its exact sequence.
The global complex has no cohomology above degree two since
\(\ell\) is odd and all places above \(\ell\) and all infinite
places are in \(S\). The last map to the sum of the local
second cohomology groups in Poitou--Tate is surjective because
its cokernel is dual to \(H^0(F,M_m^\vee(1))\), which vanishes
already at \(p_0\). Here the map in question is localization
from global second cohomology. At \(u\mid\ell\) the map to
the second cohomology of \(M_m/B_u(M_m)\) is surjective
by local cohomological dimension two. At an old unramified
condition the denominator complex has cohomological
dimension one, so the quotient has the same second
cohomology as the full local complex. The localization
onto the second cohomology of our target is thus surjective.
The fiber has zero cohomology in degrees three and higher.
We use here the local and global duality and Euler formulas
of Milne, Theorems I.2.8, I.4.10 and I.5.1.~\cite{VII-V-V-milne98}

The Euler characteristic of this fiber in lengths is zero.
In detail, put \(s=\sum_{i<j}n_i n_j\). The rank of \(M\)
is \(2s\), its upper triangular part has rank \(s\), and its
invariants at every real place have rank \(s\), since opposite
blocks are interchanged. The global Euler characteristic is
\(-2smd+smd=-smd\). The sum of the local Euler
characteristics of \(M_m/B_u(M_m)\) at \(u\mid\ell\) is also
\(-smd\). At an old unramified condition the full local complex
has Euler characteristic zero, and so does its denominator,
whose two cohomology groups are the kernel and cokernel of
Frobenius minus one on the same finite group. Subtracting
the target characteristic gives zero as asserted.

It follows that \(h(X_{N,e})\le h(\mathcal H^1_N(M_e))\).
The kernel of the map from this first cohomology to the
ordinary Selmer group is a quotient of the degree-zero
cohomology of the target. At the old unramified conditions
this vanishes since inflation is an isomorphism on invariants
and injective on \(H^1\). Above \(\ell\) the target
invariants are bounded independently of \(e\) by the
separation of the lower triangular character factors and
the Smith argument of Lemma~\ref{lem:localranks368}.
Lemma~\ref{lem:odbase369} now gives the asserted length bound.

For \(e<m\), multiplication by \(z^e\) followed in the
short exact sequence by reduction gives the maps
\[
        0\longrightarrow M_{m-e}\xrightarrow{\,z^e\,}
        M_m\longrightarrow M_e\longrightarrow0 .
\]
They give a degreewise exact sequence of the three fiber
complexes. At \(u\mid\ell\) all flag quotients lift freely.
At an old unramified condition, inertia acts through a fixed
finite group of order prime to \(\ell\), so taking invariants
is exact by averaging. Surjections of finite coefficient sets
permit continuous lifts of cochains by composing with a
section of sets. This verifies exactness also on the local
quotient complexes. The vanishing in degree three makes
\(X_{N,m}\to X_{N,e}\) surjective. It likewise makes
\(X_{N,m}\to X_{N,m-e}\) surjective by interchanging the
two exponents. The image of \(X_{N,m-e}\to X_{N,m}\) is
therefore the image of multiplication by \(z^e\), by composing
with this second surjection. Exactness gives the claimed
isomorphism. For \(e=m\) it is immediate.
\end{proof}

\begin{lemma}[A relation with a unit coefficient]
\label{lem:primrelation369}
If a finite \(A_m\)-module \(X\) satisfies
\(h(X/z^eX)\le be+C\) at every \(1\le e\le m\), then for
any \(x_0,\ldots,x_b\in X\) there are \(c_i\in A\), at
least one a unit, such that \(z^C\sum_{i=0}^b c_i x_i=0\).
\end{lemma}
\begin{proof}
Order the exponents of the nonzero cyclic summands of \(X\)
decreasingly. If there are at least \(b+1\), let \(e\) be
the \((b+1)\)-st exponent. Then
\((b+1)e\le h(X/z^eX)\le be+C\), so \(e\le C\).
Thus \(z^CX\) has at most \(b\) generators. The submodule
\(Y\) generated by \(z^Cx_0,\ldots,z^Cx_b\) also has at
most \(b\), by lifting to a free module of rank at most
\(b\) over \(A\). Their images in \(Y/zY\) are dependent.
Lift a nonzero relation over \(k'\). Its value belongs to
\(zY\), so subtract a combination of the same generators
with all coefficients divisible by \(z\) to make its value
zero. Some coefficient remains a unit.
\end{proof}

For the mixed extension introduce formal elements
\(e_1,\ldots,e_b,e'_1,\ldots,e'_b,\theta\) over \(A\).
Their only nonzero pairwise products are
\(e_j e'_i=e'_i e_j=\delta_{ij}\theta\). They and \(1\)
are a basis of a commutative algebra, since all triple
products of its added elements vanish. In matrices over
this algebra set
\[
            U=\bigoplus_{j=1}^b E e'_j\ \oplus\ M\theta .
\]
Give \(U\) the conjugation action by
\(R(g)=(1+\sum_j\nu_j(g)e_j)r(g)\) on \(G_L\), and negative
adjunction at \(c\), extended coefficientwise. These preserve
\(U\). Indeed the extra first-order action on an \(E e'_j\)
entry gives an off-diagonal commutator times \(\theta\).
There is a short exact sequence of \(G_F\)-modules
\[
              0\longrightarrow M\theta\longrightarrow U
               \xrightarrow{\ \pi\ } E^{\,b}\longrightarrow0.
\]
Locally conjugate by the same \(1+\sum Y_{u,j}e_j\) used to
triangularize \(R\) in Section~\ref{sec:selmerlength368}.
In the resulting local coordinates use the upper triangular
part for \(B_u(U)\) and transport it back. The sequence
on upper parts and the sequence on their quotients are
exact because in triangular coordinates projection reads
the \(e'_j\)-entries and the kernel reads the \(\theta\)-entries.
The change of coordinates induces the identity on the
coefficient submodule and quotient in their respective
ordinary bases.

\begin{theorem}[Uniform lifting of a combination of ordinary cocycles]
\label{thm:selmerlift369}
There exists \(C\ge0\) with the following property for all
\(1\le m\le N\). For \(0\le i\le b\), let
\((y_{ij})_{1\le j\le b}\) be a tuple of polarized ordinary
\(E_m\)-valued cocycles on the CM group at level \(N\), with
the tame condition of the original problem at old primes.
There are \(c_i\in A\), at least one a unit, and a polarized
ordinary \(U_m\)-valued cocycle \(Y\) at the same level such
that for every group element \(g\) and every \(j\),
\[
                (\pi Y)_j(g)=z^C\sum_{i=0}^b c_i y_{ij}(g).
\]
The tame condition is retained. The same \(C\) may be used
for all the selected curves on the fixed interval at the
given coefficient prime.
\end{theorem}
\begin{proof}
Extend each tuple to \(G_F\) with value zero at \(c\). Write
it as \(y_i\). Its cohomology class lifts to
\(\mathcal H^1_N(E_m^{\,b})\) because of its local conditions,
as is seen from the description of cocycles of the fiber.
Choose one such lift for each \(i\). The displayed sequence
of coefficients induces a degreewise exact sequence of the
three fiber complexes at precision \(m\). This follows at
the flags from the preceding coordinate calculation. At
the old unramified conditions inertia on all three modules
acts through the same prime-to-\(\ell\) quotient, as \(r\)
and every \(\nu_j\) were taken on \(\Gamma_a^0\). Hence
inertia invariants are exact there. These facts and the
cochain lifting argument of Proposition~\ref{prop:obsquot369}
give the required exact sequence.

Apply its connecting map to the chosen lifts and call the
images \(x_i\in X_{N,m}\), identifying \(M\theta\) with \(M\).
Proposition~\ref{prop:obsquot369} and
Lemma~\ref{lem:primrelation369}, with \(C=C_1\), kill the
obstruction of \(z^C\sum_i c_i y_i\) with some \(c_i\) a unit.
Exactness gives a lift in \(\mathcal H^1_N(U_m)\).
Choose a representative for its global cocycle. Its projection
differs from the stated linear combination by a global
coboundary. Lift its witness from \(E_m^{\,b}\) to \(U_m\)
and subtract its coboundary. The projection now agrees as
a function on the group. The remaining value \(u\) at \(c\)
belongs to \(M_m\theta\) and satisfies \(cu=-u\). Subtract
the coboundary with witness \(-u/2\) in that kernel. This
makes the value at \(c\) zero and preserves the projection.
Restriction to \(G_L\) consequently has the required
polarization.

Local coboundary changes preserve the ordinary and unramified
classes. At an old finite prime a cocycle with unramified
class vanishes on the prescribed inertia kernel because the
action there is trivial. If \(p_0\) is included in that
condition, the full local \(H^1\) on \(U_m\) equals the
unramified group there as well. On its two coefficient
subquotients the twisted dual has no Frobenius invariants
by Lemma~\ref{lem:frobtest366}, and the same follows on
the extension. These vanishings give descent through the
closed normal subgroup used in the tame condition. The
diagonal assertion at the ordinary places follows from the
local abelianization calculation above. We have obtained
the stated ordinary cocycle on the required group. Finally
there are only finitely many selected curves on the interval,
and increasing \(C\) preserves the assertion by multiplying
the lift by the additional power of \(z\).
\end{proof}

The theorem proves the condition called \(B2\) in
\srcloc{Development/AutomorphicArithmetic/AutomorphicVarieties/PatchingModules/ThetaCoefficientMapEq}{58}
on these fields and prime sets. Its coefficients are allowed
to depend on the tuple, while its exponent is uniform.
Together with Corollary~\ref{cor:bfirst368} it supplies the
two cohomological estimates at these auxiliary levels. To
obtain the bound for the integral group in
Theorem~\ref{thm:classicaldefect360} one still needs control
of the action on the automorphic modules and of the passage
to the prescribed characteristic-zero weights.


\subsection{Integral comparison of the hyperspecial and Iwahori levels}
\label{sec:levelmaps370}

Use the totally real field \(F\), its CM extension \(L\), and
the sets \(\mathcal Q_N\) of Section~\ref{sec:selmerlength368}.
At a selected place \(v\) the definite unitary group in three
variables is \(\mathrm{GL}_3(F_v)\). Write \(p_v\) for the size of
the residue field and \(\mathfrak o_v\) for the ring of integers.
Take \(K_v=\mathrm{GL}_3(\mathfrak o_v)\), its standard Iwahori
subgroup \(I_v\) (upper triangular modulo the maximal ideal),
and \(I_{1,v}\) the kernel of the diagonal map from \(I_v\) to
the maximal \(\ell\)-group quotient of
\((\mathbf F_{p_v}^{\,\times})^3\). Write \(\Delta_N\) for the
product of these diagonal quotients. In this section a
hyperspecial, Iwahori or \(I_1\) level uses the corresponding
subgroup at every selected place. All other levels are the
same in the three spaces.

\begin{proposition}[Ordinary forms at the diagonal cover]
\label{prop:diamond370}
Fix a sufficiently small tame level away from the selected
places. The dual ordinary module at \(I_1\) level is finite
free over \(\Lambda[\Delta_N]\), where \(\Lambda\) is the
completed group ring of the principal diagonal units at
\(\ell\). Its quotient by the diamond augmentation ideal
is the dual ordinary module at Iwahori level. These assertions
hold also after applying a common idempotent commuting with
the diamond action and the level maps.
\end{proposition}
\begin{proof}
The definite group for \(L/F\) and its ordinary spaces are
as in Proposition~\ref{prop:ordinaryfamily363}, now over \(F\).
The extension is unramified at all finite places and split
above \(\ell\). Fix a deep principal subgroup at one odd
split tame place outside every \(\mathcal Q_N\), for example
at \(p_0\). A sufficiently deep such subgroup is torsion-free
by the matrix logarithm. Stabilizers of the finite adelic
double cosets are finite by definiteness. At this place they
lie in conjugates of the torsion-free subgroup, and a rational
element with identity component at one place is the identity.
Thus at every fixed level at \(\ell\) the
finite set of double cosets at \(I_1\) level is a disjoint
union of free \(\Delta_N\)-orbits, one over each coset at
Iwahori level. An algebraic coefficient lattice is pulled
back along this cover, since the new places act trivially on
the coefficients. On each orbit, choosing one point identifies
the forms with a copy of the regular representation tensored
with that lattice. Invariants are functions on the quotient.
On a regular representation the norm sends the class of a
single basis vector in coinvariants to the sum of the orbit
in invariants. It is an isomorphism over the integers.
Applying the ordinary idempotent at \(\ell\) commutes with
each of these operations. This is also the finite-level
calculation of Newton and Thorne, Lemma~4.3.~\cite{VII-V-VI-nt119}

Let \(O\) be the coefficient DVR and \(k_O\) its residue
field. By integral ordinary control for a sufficiently small
level, the dual ordinary module \(M_1\) is finite free over
\(\Lambda\). Its reduction modulo the maximal ideal
\(\mathfrak m_\Lambda\) is the dual of the ordinary forms over
\(k_O\) at the first principal diagonal level at \(\ell\).
Indeed the principal diagonal invariants in the discrete
ordinary tower are the ordinary space at that level, and
taking its uniformizer torsion gives this reduction on the
dual. These are Proposition~2.5.3 and its preceding control
lemma in Geraghty. The identifications commute with the
operators away from \(\ell\).~\cite{VII-V-geraghty363}

The module \(M_1/\mathfrak m_\Lambda M_1\) is therefore a
summand of a free \(k_O[\Delta_N]\)-module by the orbit
calculation, also on taking the dual with the inverse group
action. This group ring is local because its augmentation ideal
is nilpotent. Indeed that ideal is generated by the elements \(\delta-1\)
for generators of a finite abelian \(\ell\)-group in
characteristic \(\ell\). The summand is free. Lift a basis
to \(M_1\). Its \(\Delta_N\)-translates give a basis modulo
\(\mathfrak m_\Lambda\), and hence a \(\Lambda\)-basis of
the free module \(M_1\) by Nakayama's lemma and equality of
ranks. This proves freeness over \(\Lambda[\Delta_N]\).
Finally, \(\Delta_N\)-invariants in the discrete tower give
the Iwahori tower by the orbit description and the compatibility
of finite-group invariants with filtered direct limits.
Pontryagin duality identifies their dual with the
coinvariants of \(M_1\). The augmentation submodule is
closed since it is a finite sum of images of a compact
module. A commuting idempotent gives the corresponding
summands in these identifications, and the same basis
argument proves freeness for the summand.
\end{proof}

For one of these places set \(p=p_v\) and write \(H\) for
the \(O\)-algebra of Iwahori Hecke operators, with the
Iwahori of volume one. Use the Bernstein variables, and put
\[
 S=O[X_1^{\pm1},X_2^{\pm1},X_3^{\pm1}],\qquad
 Z=S^{\mathfrak S_3},\qquad
 \delta=\prod_{i<j}(X_i-X_j).
\]
For the Bernstein normalization one may take the square root
of \(p\) congruent to one in \(\ZZ_\ell\), which exists by
Hensel's lemma. The Bernstein presentation gives
\(H=\bigoplus_{\sigma\in\mathfrak S_3}S T_\sigma\), with center
\(Z\). Use the convention in which, for an adjacent transposition
\(s_i\) and \(f\in S\),
\[
 T_i f=s_i(f)T_i+(p-1)\frac{f-s_i(f)}{1-X_{i+1}/X_i},
 \qquad (T_i-p)(T_i+1)=0 .
\]
One obtains this convention by choosing the order of the
positive roots. The divided difference here lies in \(S\).
These are the Bernstein presentation of Haines, Kottwitz
and Prasad and its integral form used by Newton and Thorne
in Section~3.~\cite{VII-V-hkp370,VII-V-VI-nt119}

Let \(J=\sum_\sigma T_\sigma\). The number of complete flags
in \(\mathbf F_p^3\) is \(P=(1+p)(1+p+p^2)\). It is a unit
in \(O\), since \(p\equiv1\bmod\ell\) and \(\ell>3\).
Put \(e=J/P\). It is the averaging idempotent of \(K_v\) on
Iwahori invariants. The identities \(T_i J=pJ\) and
\(J^2=PJ\) follow either from the finite flag double cosets
or from the quadratic relation by pairing the Weyl elements
\(\sigma,s_i\sigma\). In particular \(He=Se\), freely of rank
one over \(S\), since the expansion of \(fe\) has coefficient
\(f/P\) at \(T_{\mathrm{id}}\).

\begin{lemma}[The spherical averaging functional]
\label{lem:symmetrizer370}
For \(f\in S\) write
\[
 L(f)=\frac{\sum_{\sigma}\operatorname{sgn}(\sigma)\sigma(f)}{\delta},
 \qquad D=\prod_{i<j}(pX_i-X_j),\qquad
 D'=\prod_{i<j}(pX_j-X_i).
\]
Then \(L\) takes values in \(Z\), the pairing
\((f,g)\mapsto L(fg)\) is perfect over \(Z\), and
\(J f e=L(D f)e\).
\end{lemma}
\begin{proof}
Setting \(X_i=X_j\) pairs off all the terms of an alternating
numerator with opposite signs. Each \(X_i-X_j\) therefore
divides the numerator. These are distinct prime factors, so
their product \(\delta\) divides it as well.
For Laurent polynomials first clear denominators by a
symmetric power of \(X_1X_2X_3\). The quotient is symmetric.
The six monomials \(u_{ij}=X_1^iX_2^j\), \(0\le i<3\),
\(0\le j<2\), form a basis over \(Z\). Indeed one first
adjoins a root of the universal monic cubic, then one of
the roots of its residual monic quadratic. The third root
is determined by the sum, and all are units since the
constant term is a unit. These relations span \(S\) in
the six monomials. Their independence follows over the
fraction field of \(Z\), since the rational function field
on the three variables has six distinct permutations fixing
that field and already has degree at most six.

The alternating numerator of the basis monomial \(X_1^2X_2\)
is \(\delta\), by the three-by-three Vandermonde determinant.
The other five have degree below three and alternating
numerator zero. Hence \(L\) reads off the coefficient of
the top monomial. For a monic extension of degree \(t\),
the pairing given by the coefficient of \(X^{t-1}\) in a
product is perfect because its matrix in the two power bases,
with one reversed, is triangular with diagonal ones. Composing this
pairing in the two successive monic extensions gives the
coefficient of \(X_1^2X_2\). Products of the two successive
dual bases give a dual basis for the composite. Thus \(L\)
is perfect integrally.

The action on \(He=Se\), in the fraction field of \(S\),
has \(T_i=\alpha_i+\beta_i s_i\), where
\[
 \alpha_i=\frac{(p-1)X_i}{X_i-X_{i+1}},\qquad
 \beta_i=\frac{X_i-pX_{i+1}}{X_i-X_{i+1}} .
\]
Since \(T_i J=pJ\) and \(\beta_i\ne0\), every value of \(J\)
in this representation is symmetric. Write \(J=\sum c_\sigma\sigma\)
there as an operator. Distinct field automorphisms are linearly
independent, so \(c_\sigma=\sigma(c_{\mathrm{id}})\).
The coefficient of the longest permutation comes only from
\(T_1T_2T_1\), specifically from choosing the reflection term
three times. It is
\(\beta_1s_1(\beta_2)s_1s_2(\beta_1)
 =\prod_{i<j}(X_i-pX_j)/\delta\).
Applying the longest permutation gives
\(c_{\mathrm{id}}=D/\delta\). The action of \(J\) on \(f\)
is therefore \(L(Df)\), as asserted. This calculation also
proves faithfulness of the action of \(H\) on \(Se\).
For each \(T_\sigma\), its reduced expression has a nonzero
coefficient at \(\sigma\), and terms of shorter length from
choosing fewer reflections. A nonzero linear combination
in the Bernstein basis thus gives a nonzero combination of
field automorphisms. Their independence can be tested on
\(S\), since any rational function has a symmetric common
denominator after multiplying its denominator over its orbit.
\end{proof}

\begin{theorem}[Integral maps between the two levels]
\label{thm:iwahori370}
For every left \(H\)-module \(V\) there are natural \(S\)-linear
maps
\[
       S\otimes_Z eV\xrightarrow{\ \Phi\ } V
                  \xrightarrow{\ \Psi\ } S\otimes_Z eV
\]
such that both composites are multiplication by the central
polynomial \(C_p=\prod_{i\ne j}(X_i-pX_j)\). These identities
persist under all base changes of \(Z\).
\end{theorem}
\begin{proof}
Use the six monomials \(u_a\) in the lemma and their \(L\)-dual
basis \(v_a\). Define
\[
       \Phi(f\otimes t)=ft,\qquad
       \Psi(y)=\sum_a u_a\otimes J(v_aD' y).
\]
The factor \(J\) puts each entry in \(eV\). Both maps are
defined over \(O\). The tensor \(\sum_a u_a\otimes v_a\)
has the same product with \(f\otimes1\) and \(1\otimes f\)
for every \(f\in S\). Indeed under the perfect pairing the
tensor identifies with the identity map of \(S\), and these
two products identify with multiplication by \(f\) on its
range and domain. Applying this tensor identity to the formula
for \(\Psi\) gives its \(S\)-linearity, since \(Z\) commutes
with \(J\).

There are six ordered pairs of distinct indices, so \(DD'=C_p\).
For \(t\in eV\), the lemma in the algebra \(He\) gives
\(J v_a D'f t=C_p L(v_af)t\).
Expanding \(f\) in the dual basis proves \(\Psi\Phi=C_p\).
For the reverse order consider the element
\(\sum_a u_a J v_aD'\) of \(H\). On \(f e\in He\) it acts
as \(\sum_a u_a L(Dv_aD'f)e=C_p f e\). Faithfulness in
Lemma~\ref{lem:symmetrizer370} gives equality with \(C_p\)
inside \(H\), hence on every \(V\). These are identities
of maps, so tensoring with any \(Z\)-algebra preserves both.
\end{proof}

\begin{corollary}[The comparison factor on the cohomology curves]
\label{cor:levelcurve370}
Apply the maps at all \(v\in\mathcal Q_N\) to the definite
ordinary modules after taking diamond coinvariants. Their
composites are the product of the \(C_{p_v}\). On specializing
the local Hecke centers at the characteristic polynomials
of any one of the curves of Theorem~\ref{thm:markedcurve366},
the kernel and cokernel of the resulting comparison are
killed by \(z^{D_0}\) at every precision \(m\le N\), with
\[
       D_0=\sum_{j=1}^{q}v_z\!\left(
         \operatorname{disc}\det(T-r(n_j'))\right).
\]
The \(n_j'\) are the fixed elements of
Section~\ref{sec:selmerlength368}; the exponent is independent
of the auxiliary level. A common exponent is obtained on the
interval by taking the maximum over the finite list of curves.
\end{corollary}
\begin{proof}
At finite level the operators at different split places
commute. For the Iwahori module, averaging at a selected
place gives precisely the module with that place hyperspecial.
The commutation with the ordinary projector allows successive
application of Theorem~\ref{thm:iwahori370}. Tensoring the
resulting maps, or composing them one place at a time,
multiplies their composite factors. They commute with the
transition maps at \(\ell\). Tensoring with the free module
\(S\) over \(Z\) is a direct sum of six copies in the chosen
basis and commutes with inverse limits. Thus the same
calculation holds on the ordinary tower. On duals one uses
the contragredient Hecke action, or equivalently applies the
integral Hecke algebra identities to the dual module.
Proposition~\ref{prop:diamond370} identifies its Iwahori
term with the diamond coinvariants at \(I_1\) level.

The map on a local center sends the elementary symmetric
polynomials to the signed coefficients of the indicated monic
Frobenius polynomial. Its constant term is a unit. In
characteristic \(\ell\), \(p_v=1\), so \(C_{p_v}\) evaluates
to the discriminant up to sign. Reciprocal Frobenius or a
common unit scaling of its three roots changes that valuation
by zero. The Chebotarev choice in Theorem~\ref{thm:primes367}
prescribes each polynomial to precision strictly exceeding
the discriminant valuation at \(n_j'\), in addition to all
the required cohomology tests. The resulting product is a
unit times \(z^{D_0}\) over \(A\), with the same identity
after reduction to \(A_m\). The two composite identities
give the kernel and cokernel assertion, including after
specialization with torsion.
\end{proof}

For the scalar in
\srcloc{Development/AutomorphicArithmetic/AutomorphicVarieties/PatchingModules/ThetaCoefficientMapEq}{213},
one can multiply \(\Psi\) by \(\prod_{i\ne j}(X_i-X_j)\) at
each place. This gives precisely both factors in that scalar.
The corollary evaluates the local centers on the curves.
To bound the mixed-characteristic length in
Section~\ref{sec:missing} one still needs the action of the
global deformation ring on the ordinary module with the
prescribed ramification and the integral comparison at the
original weights.


\subsection{An automorphic prime through the curves on the deep layer}
\label{sec:basesupport371}

The comparison of levels in Section~\ref{sec:levelmaps370}
applies to the values of the local centers on any of our
curves. We now prove that the specialized spherical module
has positive rank there. Fix the deep layer
\(F=F_a^+\), \(L=L_a\) of Theorem~\ref{thm:commonframe365}.
Write \(F_0=F_0^+\), \(d_0=[F_0:\QQ]\), and use \(L_0=KF_0\).
For a pair \(r,\alpha\) on the interval let \(B'\) denote
one of the integral Hecke components through the adjoint on
\(L_0\) in Proposition~\ref{prop:bigweight364}. Its product-one
weight algebra is \(\Lambda_0\), of dimension \(1+2d_0\).
We use its normalized trace \(t\) of the representations
twisted by \(\kappa_\ell\). Removing that twist gives the
conventional multiplier \(\kappa_\ell^{-2}\) in the ordinary
determinant problem. The twist is fixed on the entire family.

\begin{lemma}[A common tame level for classical base change]
\label{lem:bcpoints371}
There is a Zariski dense set of characteristic-zero points
of \(B'\) whose automorphic representations base change
cuspidally to \(L\). Their base changes descend to ordinary
forms on the definite unitary group over \(F\) at a single
tame level. Their ordered local weights are obtained by
restriction of the lower-layer Galois characters. On inertia
away from \(\ell\) the normalized Galois representations
factor through the prescribed quotient and have the finite
inertial type \(\alpha\oplus1\oplus\alpha^{-1}\). The tame
level can be chosen for all the pairs on the interval.
\end{lemma}
\begin{proof}
Use the purely algebraic weights in the proof of
Proposition~\ref{prop:bigweight364}. Require the successive
labeled Hodge exponents to differ by at least two at every
embedding, in their ordinary order. Increasing the first
weight variation while holding the second fixed as in that
proof realizes the larger gaps in every congruence class.
These weights are still dense. At the adjoint on \(L_0\)
the two-dimensional newform has projective closure
\(\mathrm{PGL}_2\) on every open subgroup, by the proof of
Theorem~\ref{thm:adjointseed361}. The trace-zero adjoint
is absolutely irreducible on \(G_L\). Choose nine group
elements \(h_i\in G_L\) whose adjoint matrices span the
matrix algebra there. The trace Gram determinant
\(\det(t(h_i h_j))\) is nonzero at the adjoint. Take its
product with the finitely many nonzero Frobenius resultants
in the proof of Proposition~\ref{prop:tamecurve365}.
Above the algebraic weights the points avoiding the zero
locus of this product remain dense in \(B'\). Indeed the
principal ideal of any nonzero element of the finite domain
\(B'\) contains a nonzero element of \(\Lambda_0\), by clearing
denominators after inverting \(\Lambda_0\setminus\{0\}\).
The same argument applies after multiplication by any
further nonzero element. Density downstairs and lying over
give the assertion.

At each of these points the nine restricted matrices span
the full matrix algebra. The classical transfer to
\(\mathrm{GL}_3/L_0\) is cuspidal, since a transfer with
several isobaric summands gives a decomposable semisimple
Galois representation. This also excludes a repeated block
of character twists from a single summand of smaller rank.
We use the attachment and decomposition of ordinary classical
representations in Geraghty's Proposition~2.7.2.
The extension \(L/L_0\) is cyclic of \(\ell\)-power degree.
In cyclic base change a cuspidal representation can lose
cuspidality at a step of prime degree \(\ell\) only in the
presence of a self-twist of that order. The central character
of a rank-three representation then makes the cube of the
twisting character trivial, impossible since \(\ell>3\).
Thus the base change is cuspidal. Its normalized representation
at good places has the trace \(t|_{G_L}\) specialized at our
point, by unramified base change and Chebotarev. Conjugate
self-duality and regular algebraicity persist.~\cite{VII-V-geraghty363,VII-V-I-ac123}

We verify ordinarity after base change at the ramified
\(\ell\)-adic places. The lower-layer ordinary flag exists
at the chosen algebraic weights by Geraghty's
Corollary~2.7.8. Each graded character is crystalline,
since on local units it is the purely algebraic character
of that weight. Via local reciprocity it differs from
the corresponding product of labeled crystalline characters
only by an unramified character. The potentially semistable
Weil--Deligne functor is exact on these representations.
Its finite inertia on the full parameter is trivial because
it is semisimple on inertia and trivial on all graded lines.
If a local field has ramification index \(e\) and residue
degree \(f\) over \(\QQ_\ell\), the valuation of the
\(f\)-fold crystalline Frobenius on a graded line is, up
to the common sign convention, the sum of its labeled
Hodge exponents divided by \(e\). This is the equality of
Newton and Hodge numbers for a crystalline character, with
\(v_\ell(\ell)=1\). Consequently successive Frobenius
valuations differ by at least \(2f\). A nonzero monodromy
operator requires two valuations differing by exactly \(f\)
by the Frobenius relation. It follows that monodromy is zero.
Restriction to any upper completion preserves crystallinity.
It repeats the labels above each embedding and multiplies
the Frobenius valuations by the relative residue degree.

Local--global compatibility at \(\ell\) for a regular
conjugate self-dual cusp form in odd rank identifies its
Frobenius-semisimplified parameter there. It is spherical
in the present case. The normalized Iwahori operator for the
first \(j\) ordered roots has valuation equal to the sum of
their valuations minus the corresponding Hodge sum, with
the same normalization as in Proposition~\ref{prop:ordinaryfamily363}.
Both terms are multiplied by the relative residue degree
under base change, so this eigenvalue is a unit for every
\(j\). The spherical generic component admits this refinement.
In particular its weights on local units are obtained by composing
the lower characters with the local norms.
~\cite{VII-V-lgp361,VII-V-geraghty363}

Away from \(\ell\), Proposition~\ref{prop:tamecurve365}
proved that all points in the chosen dense set have the
stated finite inertia on the lower layer, and hence after
restriction as well. At a prime of the conductor the extension
\(L/F\) is split by Proposition~\ref{prop:auxfield361}.
The inertial characters extend to the local Weil group.
After Frobenius semisimplification each isotypic part is a sum
of characters, and monodromy is zero. The local Langlands
component is thus a quotient of a principal series with
the three indicated restrictions to units. Section~\ref{sec:seed361}
shows that a principal congruence level on which all three
characters are trivial supplies vectors in this quotient:
the section on the identity chart and its translates
generate the inducing representation. This bounds the level
independently of the unramified character values. At the
other primes the base change is spherical. Descent to the
definite inner form, unramified at the nonsplit finite places,
is the descent used in Proposition~\ref{prop:ordinaryfamily363}.
That result applies since \(L/F\) is unramified at finite
places, split above \(\ell\), and the rank is three.
Impose the same sufficiently small level at \(p_0\) as in
Proposition~\ref{prop:diamond370}. Take the maximum of the
conductor levels over the finite list of pairs. Smaller
subgroups keep the designated vectors, proving the last
assertion.~\cite{VII-V-geraghty363,VII-V-V-shin}
\end{proof}

At this tame level let \(\Lambda_a\) be the upper-layer weight
algebra on the chosen prime-to-\(\ell\) diamond summand, with
three principal-unit coordinates at each place over \(\ell\).
Let \(\mathbf T_a\) be generated over \(\Lambda_a\) by the
spherical Hecke operators at all good split places on the
dual ordinary module \(\mathcal M_a\). Use the Frobenius
normalization of Proposition~\ref{prop:ordinaryfamily363}.
Equivalently one may take the coefficients of each good
Frobenius polynomial, since the known powers of the norm in
the normalization are units. Inverting the Satake determinant
is allowed at these places. This algebra is finite over
\(\Lambda_a\), being a submodule of the endomorphisms of a
finite free \(\Lambda_a\)-module.

\begin{proposition}[The restricted integral family on the upper Hecke algebra]
\label{prop:bchecke371}
There is a continuous map \(\mathbf T_a\to B'\) sending
good Frobenius coefficients to those of \(t|_{G_L}\) and
weight variations to their pullbacks by the local norms.
Its image \(D_a\) is a complete local domain, finite over
a power series algebra of dimension \(1+2d_0\). Its
dimension is \(1+2d_0\), and \(B'\) is finite over \(D_a\).
The ring \(D_a\) contains all restricted traces and the
values of the ordered characters on the full local
decomposition groups at \(\ell\).
\end{proposition}
\begin{proof}
The norm on each group of principal units gives a continuous
map on the completed weight algebras, using variations
relative to the adjoint on both layers. It is the required
weight map by Lemma~\ref{lem:bcpoints371}. Assign the good
split Hecke generators their indicated values in \(B'\).
Any polynomial relation over \(\Lambda_a\) among these
operators vanishes on every upper ordinary form in that
lemma. Its proposed value therefore vanishes on the dense
set in \(B'\) and is zero. This proves well-definedness.
The map is continuous since it is linear over the continuous
weight map on the finite \(\Lambda_a\)-module \(\mathbf T_a\).
The compact source has closed image.

For each lower place \(u_0\mid\ell\), choose one upper place
over it. The norms of its principal units contain the
\([F_u:F_{0,u_0}]\)-th powers of all lower principal units,
by viewing a lower unit inside the upper field. These powers
are open in the lower principal units. After taking products
and imposing the product-one weights downstairs, the image
of the upper weight group contains an open subgroup of the
free rank-\(2d_0\) group defining \(\Lambda_0\). Its completed
group ring \(\Lambda_{\rm op}\) is a power series ring over \(O\)
contained in \(D_a\). The index is finite, so coset
representatives give a finite free basis of \(\Lambda_0\)
over \(\Lambda_{\rm op}\). Thus \(B'\) is a finite
\(\Lambda_{\rm op}\)-module, and so is its submodule \(D_a\).
Also \(B'\) is finite over \(D_a\). The finite domain
\(D_a\) over the complete local ring \(\Lambda_{\rm op}\)
is complete and local, by the idempotent decomposition of
a finite algebra over a Henselian local ring. Dimension
is preserved under integral extensions.

Good split Frobenii are dense up to conjugacy in \(G_L\).
The closed image thus contains \(t(h)\) for every \(h\in G_L\).
At an upper place \(u\mid\ell\) take a local inertia lift \(s\)
with distinct residual characters as in
Theorem~\ref{thm:commonframe365}. Its three ordered diagonal
values \(a_1,a_2,a_3\) lie in \(D_a^\times\) by the weight
map, including the fixed base characters. Their differences
are units. For any \(g\in G_{F_u}\) form the \(D_a\)-linear
combination of \(t(s^j g)\), \(0\le j\le2\), whose coefficients
are those of
\[
                    P_i(X)=\prod_{k\ne i}\frac{X-a_k}{a_i-a_k}.
\]
At every dense classical point this reads the \(i\)-th
diagonal entry of \(g\) on the ordinary flag. In fact \(s\)
is diagonalizable with those three eigenvalues and the first
\(j\) ordered eigenlines span the \(j\)-th step. Thus the
combination defines on the whole local group a continuous
function with values in \(D_a\). The character identity
holds on that dense set, hence in \(D_a\). Its value at
the identity is one, so its values are units. This proves
the assertion also on the unramified part of the local group.
\end{proof}

Let \(R_a\) be the universal ordinary polarized determinant
ring for \(\Gamma_a^0\) with the ordered local characters of
Section~\ref{sec:commonframe365}. A degree-three determinant
here records the characteristic polynomial on the group
algebra after scalar extension. Since \(3!\) is a unit, it
is determined by its trace. Use the fixed cyclotomic twist
above when writing this ring in the source's normalization.
Let \(J_{\rm cl}\subset R_a\) be the intersection of the
kernels of its characteristic-zero evaluations coming from
ordinary classical forms of the tame level just fixed that
satisfy its ramification conditions, with their ordered
characters. This definition includes the points constructed
in Lemma~\ref{lem:bcpoints371}.

\begin{theorem}[A stable automorphic prime below each curve prime]
\label{thm:stableprime371}
There is a surjection \(R_a\to D_a\). Its kernel
\(\mathfrak q\) contains \(J_{\rm cl}\), and its generic
determinant is that of an absolutely irreducible representation
over an algebraic closure of its residue field. For every
curve constructed from \(B'\) in Theorem~\ref{thm:rescurve364},
let \(\mathfrak p\) be its prime in \(R_a\) after restriction
to \(L\). Then
\[
                   J_{\rm cl}\subseteq\mathfrak q
                          \subsetneq\mathfrak p .
\]
This includes the case where the generic representation on
the curve itself has two constituents.
\end{theorem}
\begin{proof}
The restricted family gives a continuous degree-three
pseudocharacter with values in \(D_a\) because its trace
identities hold on the dense classical evaluations. The
trace--determinant correspondence with \(3!\) invertible
gives the degree-three determinant. For example, on any
group algebra element its third coefficient is computed from
the traces of its first three powers by Newton's identities,
and the identities of polynomial laws can be tested
coefficientwise on generic finite sums of group elements.
The polarization and the reduction at the closed point
are those of the adjoint. On the prescribed inertia kernel
one has \(t(hg)=t(h)\) for every \(h\) and every such \(g\),
by Proposition~\ref{prop:tamecurve365}. Consequently the
trace factors through the stated tame quotient. On each
ordinary group the ordered character word in the universal
condition has zero trace after multiplication by any group
matrix. Indeed on each classical flag its factors, starting
at the right, successively lower the last, penultimate and
first step to zero. The traces vanish identically by density.
The characters in Proposition~\ref{prop:bchecke371}
have the required residual restrictions, as seen on the
adjoint, and are trivial on the local commutator closure.
The local abelianization has no \(\ell\)-primary torsion,
as verified in Section~\ref{sec:selmerlift369}. These
properties give a pair in the ordinary determinant problem
over \(D_a\), hence a map \(R_a\to D_a\).~\cite{VII-V-chenevier364}

The image contains every good trace coefficient and every
weight coordinate by the defining evaluations. It is closed,
since \(R_a\) is a complete Noetherian local ring with finite
residue field and the map is continuous. It therefore contains
the whole image of \(\mathbf T_a\) by its generators. This
proves surjectivity. The elements of \(J_{\rm cl}\) vanish
at all points of Lemma~\ref{lem:bcpoints371}, so their images
in the domain \(D_a\subset B'\) are zero. The nine-element
trace Gram determinant of that lemma is nonzero in \(D_a\).
On the semisimple representation of its generic trace over
an algebraic closure the nine matrices are consequently
linearly independent. Their span is the full matrix algebra,
so the representation is absolutely irreducible.

Every curve in the statement is obtained by composing
\(B'\) with the quotient at its defining prime and a finite
extension of the resulting coefficient ring. On \(G_L\)
this composition gives its restricted trace. The local
characters also agree, by the projector formula of
Proposition~\ref{prop:bchecke371}. Marking a
lattice in Theorem~\ref{thm:markedcurve366} preserves these
values. Thus the evaluation \(R_a\to A=k'[[z]]\) factors
through \(D_a\). Its kernel strictly contains \(\mathfrak q\)
because \(\ell\) is nonzero in \(D_a\)
and zero in \(A\).
\end{proof}

\begin{corollary}[Nonzero ordinary modules at all selected levels]
\label{cor:curvesupp371}
For each of these curves the spherical specialization
\(\mathcal M_{a,A}=\mathcal M_a\otimes_{\mathbf T_a}A\)
has positive rank \(s\) over \(A\). At level \(\mathcal Q_N\)
specialize the operators away from the selected places and
their local Bernstein centers by the same system. Over
\(k'((z))\) the Iwahori specialization has dimension
\(6^q s\), where \(q=|\mathcal Q_N|\). Its integral comparison
with the spherical specialization after adjoining the ordered
roots has kernel and cokernel killed by \(z^{D_0}\) as in
Corollary~\ref{cor:levelcurve370}. The same conclusions
hold on taking diamond coinvariants at \(I_1\) level.
\end{corollary}
\begin{proof}
As a module over \(\mathbf T_a\), \(\mathcal M_a\) is finite
and faithful by construction. Its support is all of
\(\operatorname{Spec}\mathbf T_a\). Localize at the prime
defined by \(\mathbf T_a\to D_a\to k'((z))\) and apply
Nakayama's lemma. Its fiber is nonzero, also after extending
the residue field to \(k'((z))\). This proves \(s>0\).

Use the common commutative algebra generated formally by
the weight and good operators away from \(\mathcal Q_N\)
and by the \(q\) local centers there. Its spherical action
maps onto \(\mathbf T_a\), since these centers give all the
omitted spherical operators by the Satake isomorphism.
Its Iwahori action commutes with the comparison maps.
Specialization over this common algebra thus gives exactly
\(\mathcal M_{a,A}\) on the spherical side, with the indicated
root algebras adjoined. Each such algebra is free of rank
six over its center by Lemma~\ref{lem:symmetrizer370},
including after base change. The composite factor is
nonzero over the fraction field by
Corollary~\ref{cor:levelcurve370}. Hence the two maps there
are isomorphisms and give the stated dimension. The integral
loss follows from the same composite identity.
Proposition~\ref{prop:diamond370} gives the assertion at
the diagonal cover.
\end{proof}

The prime \(\mathfrak q\) supplies the stable specialization
under the given prime required in the second-type case of
\srcloc{Development/AutomorphicArithmetic/AutomorphicVarieties/WeilGroupCharacters/UniversalOrdinaryDeterminantRing}{288}
for the ideal \(J_{\rm cl}\). To use that implication one
must still propagate support to the other primes in its
conclusion. The fixed-weight integral lengths and the
comparison of the projected theta periods in
Section~\ref{sec:missing} involve those further arguments.


\subsection{The full determinant component through the cuspidal adjoint}
\label{sec:adjcomponent372}

Work on \(F=F_a^+\), \(L=L_a\) with the common tame level
of Lemma~\ref{lem:bcpoints371}. Put \(d=[F:\QQ]\). We allow
all three ordinary weights at each place above \(\ell\).
Thus the principal weight algebra \(\Lambda_a\) is a power
series ring over \(O\) in \(3d\) variables. Fix the
prime-to-\(\ell\) diamond characters of the adjoint. Take
\(O\) large enough for its eigenvalues and its chosen basis,
and write \(J=\operatorname{Frac}O\). Let \(x_0\) denote the
characteristic-zero evaluation of the ordinary determinant
ring \(R_a\) at this adjoint, and \(y_0\) its evaluation of
the good Hecke algebra \(\mathbf T_a\). We denote the kernels
of these evaluations by the same letters. Translate the
ordinary weight variables so that they vanish there.

\begin{proposition}[The ordinary determinant on a full Hecke component]
\label{prop:upfamily372}
Let \(\mathfrak l\) be a minimal prime of \(\mathbf T_a\)
contained in \(y_0\), and put \(B=\mathbf T_a/\mathfrak l\).
There is a surjection
\[
                         R_a\longrightarrow B
\]
matching all good Frobenius polynomials and the ordinary
weights. Its kernel contains \(J_{\rm cl}\) of
Section~\ref{sec:basesupport371}. The domain \(B\) is finite
over \(\Lambda_a\) and contains it. Its generic determinant
is absolutely irreducible.
\end{proposition}
\begin{proof}
The algebra \(\mathbf T_a\) acts faithfully on the free
\(\Lambda_a\)-module of ordinary forms and is a subalgebra
of its endomorphisms. Every nonzero element of \(\Lambda_a\)
is therefore a nonzero divisor in \(\mathbf T_a\). A minimal
prime in a Noetherian ring consists of zero divisors, so
\(\mathfrak l\cap\Lambda_a=0\). The finite domain \(B\) is
local, and \(y_0\) on \(B\) has the same residual system as
every classical point on it after reduction to an algebraic
closure of the residue field. Indeed a finite domain over
the complete local weight ring is local, and evaluation at
the integral adjoint identifies its residue field with that
of \(O\).

Purely algebraic dominant weights, with consecutive Hodge
gaps at least two, are dense in the full weight algebra
just as in Proposition~\ref{prop:bigweight364}. This time
the third exponents are independent. In any residue classes
for the exponents one can increase successive gaps by
arbitrarily large multiples of the modulus while retaining
the prescribed tame diamonds. On a sufficiently small
logarithmic neighborhood the matrix of embeddings on a
basis of each local field has nonzero determinant, so the
identity theorem in each variable gives the asserted
density. At these weights the good Hecke operators are
semisimple on the classical characteristic-zero spaces.
Their evaluations detect \(\mathbf T_a\) because its elements
are matrices over \(\Lambda_a\). Compactness and Chebotarev,
as in Proposition~\ref{prop:ordinaryfamily363}, now put the
continuous traces of all \(h\in G_L\) in \(\mathbf T_a\).
Take their images in \(B\). Above the indicated weights the
points of \(B\) avoiding any specified nonzero element are
dense, by finiteness over \(\Lambda_a\) and the denominator
argument of Lemma~\ref{lem:bcpoints371}. Each such point
occurs on the classical module, since the module is finite
and faithful over \(\mathbf T_a\). The trace Gram determinant
of the nine group elements of that lemma is nonzero at \(y_0\),
so it is nonzero in \(B\).

We check the ramification condition on this full component.
Away from the fixed finite set of old tame places and \(\ell\)
the classical systems are unramified. At an old finite place
pass to a completion of \(LE_*\), which kills
\(\bar\alpha\) on inertia. In a stable lattice the residual
image of this inertia is unipotent upper triangular, since
the residual semisimplification is the split adjoint triple.
Its integral image is pro-\(\ell\) and hence factors through
the \(\ZZ_\ell(1)\) part of tame inertia. If the residue
field here has order \(q_0\), Frobenius permutes the three
eigenvalues of a generator by raising them to \(q_0\).
Their orders divide \(q_0^6-1\). For every fixed original
inertia element a power bounded by the local extension degree
lies in this subgroup. Its characteristic coefficients thus
take finitely many possible values at the classical points.
The monic polynomial with these possible values and their
coefficient conjugates as roots vanishes on the coefficient
function in \(B\). Since \(B\) is a domain, that function
equals its value at \(y_0\), because the product of the
other distinct linear factors is nonzero there. In
particular all these inertial
polynomials coincide with those of the adjoint.

At the adjoint the three eigenvalues of a Frobenius over
the chosen completion have equal complex absolute values,
by the good and bad prime calculations of
Proposition~\ref{prop:awaytype363}, which persist under
local extension. The resultant testing whether their ratios
include \(q_0\) is nonzero there. Avoid the product of the
finitely many such resultants and the nine-element Gram
determinant. At the resulting dense classical points the
logarithm of unipotent inertia on \(LE_*\) must vanish,
since Frobenius would act on it with eigenvalue \(q_0\).
Inertia on \(L\) is now finite with exactly the adjoint
inertial characters. It acts trivially on the designated
kernel at every old finite place. These points have
absolutely irreducible trace and satisfy the tame
conditions of \(J_{\rm cl}\).

The identities \(t(hg)=t(h)\) for \(g\) in the prescribed
inertia kernel hold on that dense set and therefore in \(B\).
For each \(u\mid\ell\) the three weights in \(\Lambda_a\)
give three unit-separated values on the separating inertia
element. The projector formula of
Proposition~\ref{prop:bchecke371} gives the three full local
characters in \(B\). The ordinary word and polarization
identities follow from the same classical evaluations, as
in Theorem~\ref{thm:stableprime371}. This constructs the
ordinary determinant and the map from \(R_a\). Its image
contains the weights and the good Hecke generators, which
generate \(B\). The dense set proves the assertion on
\(J_{\rm cl}\). The nonzero Gram determinant gives absolute
irreducibility of the generic determinant.
\end{proof}

\Needspace{5\baselineskip}
\begin{lemma}[The characteristic-zero tangent space of the determinant ring]
\label{lem:detangent372}
The tangent space of \((R_a)_{x_0}\) over \(J\) has dimension
at most \(3d\). Its map to the tangent space of the weights
is injective.
\end{lemma}
\begin{proof}
Consider a lift of \(x_0\) to \(A'=J[\epsilon]/(\epsilon^2)\).
Precompose with the inverse of the fixed cyclotomic twist
on determinants and local characters so that the base
representation is \(U=\operatorname{Ad}^0(W)|_{G_L}\).
The trace and the local characters remain continuous in
the \(\ell\)-adic vector space topology of \(A'\).
To see this for an arbitrary algebraic tangent map, present
the complete ring \(R_a\) as a quotient of
\(O[[X_1,\ldots,X_h]]\). The values \(a_i=x_0(X_i)\) lie
in the maximal ideal of \(O\). Translation by the \(a_i\)
is an automorphism of the complete power series ring.
Modulo the square of the ideal of the point, the image of
any series is its value plus its first derivatives at \(a\)
times the increments. These values and derivatives depend
continuously on the series for the original adic topology,
since the \(a_i\) are topologically nilpotent. This proves
the required continuity.

The adjoint on \(G_L\) is split and absolutely irreducible
over \(J\). Over the Henselian local ring \(A'\), Chenevier's
Theorem~2.22(i) therefore identifies the lifted determinant
with the determinant of a representation lifting \(U\),
unique up to strict conjugacy. The representation \(r'\)
is continuous. The matrices of the nine \(h_i\) are an
\(A'\)-basis of \(M_3(A')\) by their trace Gram determinant,
and the coordinates of \(r'(g)\) are obtained by inverting
that fixed Gram matrix against the continuous traces at
\(h_i g\).~\cite{VII-V-chenevier364}

For \(g\) in the designated tame inertia kernel all the
traces of \(r'(h)(r'(g)-1)\) are zero. Nondegeneracy of
the matrix trace pairing makes \(r'(g)=1\). At an ordinary
place take \(s\) with distinct residual eigenvalues and
write \(a_1',a_2',a_3'\) for the lifted character values
there. The ordinary trace word multiplied by every \(r'(h)\)
has trace zero, so the word itself vanishes. Taking all
three letters equal to \(s\) gives the annihilating
polynomial \(\prod_i(X-a_i')\). Its roots are unit-separated,
so \(r'(s)\) has three eigenlines over \(A'\) by its
polynomial projectors and the basepoint ranks. Now replace
only the \(k\)-th letter by an arbitrary local \(g\).
The product to its left is invertible on each eigenline
with index at least \(k\), and the product to its right
is invertible on each eigenline with index at most \(k\).
Consequently the entries of \(r'(g)-\psi_k(g)\) in rows
\(i\ge k\) and columns \(j\le k\) are zero in this basis.
Taking all \(k\) proves upper triangularity with exactly the
given diagonal characters. Here a scalar subtracted from a
matrix denotes that scalar times the identity.

It remains to pass from the polarized determinant on \(G_L\)
to the signed representation over \(F\) used in
Theorem~\ref{thm:weightchart362}. Let \(\iota(B)=
S^{-1}B^{-\mathsf T}S\), with \(S\) the symmetric adjoint
form there. The polarization and uniqueness for the lifted
determinant give a matrix \(C\) such that
\[
           r'(cgc^{-1})=C\iota(r'(g))C^{-1}\qquad(g\in G_L).
\]
By Schur's lemma its reduction differs from \(U(c)\) by a
scalar, which we remove. Applying the formula twice shows
that \(C\iota(C)\) is a scalar, since the image matrices
span the full matrix algebra. Write that scalar as \(a\).
Then \(\iota(C)C=a\) as well, while applying \(\iota\) gives
\(\iota(C)C=a^{-1}\). Hence \(a^2=1\), and reduction gives
\(a=1\) because two is invertible. Assigning \((C,-1)\)
to \(c\) defines the required lift on \(G_F\). At finite
places away from \(\ell\) the adjoint inertia kernel is
exactly the one checked above, since \(L/F\) is unramified
and the adjoint has the finite type
\(\alpha\oplus1\oplus\alpha^{-1}\).

Two determinant tangent maps with the same weights now
give admissible lifts in that theorem with the same weights.
They are strictly conjugate, so their determinants agree.
Their full ordered characters agree by the uniqueness of
the eigenline flag, or directly by the trace projector
formula. The universal ring is topologically generated by
these functions and the image of \(O\). To verify this, work
in each quotient by a power of the maximal ideal and take
the finite \(O\)-subalgebra generated there. The determinant
descends to it by the trace formulas and \(3!\in O^\times\).
Universality on these finite local rings
and uniqueness after inclusion show that the universal map
to each quotient takes values in its generated subalgebra.
Passing to the limit gives the assertion. The two continuous
tangent maps thus agree. The space of weight tangents has
dimension \(3d\) by Theorem~\ref{thm:weightchart362}, proving
the lemma.
\end{proof}

\begin{theorem}[The unique component through the adjoint]
\label{thm:adjcomponent372}
There is a unique minimal prime \(\mathfrak q_*\) of \(R_a\)
contained in \(x_0\) and a unique minimal prime
\(\mathfrak l_*\) of \(\mathbf T_a\) contained in \(y_0\).
The weight, good Hecke and determinant evaluations identify
\[
          R_a/\mathfrak q_*\ \simeq\
                    \mathbf T_a/\mathfrak l_* .
\]
This integral domain is finite over the full weight algebra
and has dimension \(1+3d\). The local ring of \(R_a\) at
\(x_0\) is regular of dimension \(3d\), and passage to
this Hecke component induces an isomorphism on local rings
at the adjoint. For every curve from
Theorem~\ref{thm:stableprime371} one has
\[
              J_{\rm cl}\subseteq \mathfrak q_*
                         \subseteq\mathfrak q\subsetneq\mathfrak p .
\]
\end{theorem}
\begin{proof}
Choose any \(\mathfrak l\) as in Proposition~\ref{prop:upfamily372}
and let \(\mathfrak q'\) be the kernel obtained there.
The contraction of the adjoint point in \(B\) to
\(\Lambda_a\) is the characteristic-zero prime of height
\(3d\) generated by the translated weight variables.
Indeed the quotient by them is \(O\), and localizing there
inverts its uniformizer. Incomparability gives the upper
bound \(3d\) for the height in \(B\), and going down for
the integral extension of the integrally closed domain
\(\Lambda_a\) gives the lower bound. The localization of
\(R_a\) at \(x_0\) thus has a local quotient of dimension
\(3d\). It is Noetherian with residue field \(J\). Its
embedding dimension is its tangent dimension, bounded above
by \(3d\) in Lemma~\ref{lem:detangent372}. Dimension is at
most embedding dimension, so they coincide and the local
ring is regular. In particular it is a domain.

The localized kernel \(\mathfrak q'(R_a)_{x_0}\) must vanish.
If it were a nonzero prime, appending the strict inclusion
of the zero prime below it to a chain in the quotient would
give a chain longer than \(3d\). By prime correspondence
for localization, \(\mathfrak q'\) is the unique minimal prime
below \(x_0\). Call it \(\mathfrak q_*\). This also proves
the local isomorphism and the dimension assertions. Applying
the same argument to any \(\mathfrak l\subset y_0\) minimal
in \(\mathbf T_a\) identifies its quotient with the fixed
\(R_a/\mathfrak q_*\). The resulting maps from
\(\mathbf T_a\) into that ring agree on weight coordinates
and on good Hecke generators, because these come from the
same universal local characters and good Frobenius
polynomials in \(R_a\). They therefore have equal kernels,
giving uniqueness of \(\mathfrak l_*\).

Finally, the map \(\mathbf T_a\to D_a\) of
Proposition~\ref{prop:bchecke371} has prime kernel below
\(y_0\) since the restricted family includes the adjoint.
It contains a minimal prime, necessarily \(\mathfrak l_*\).
Composing the displayed isomorphism with the quotient to
\(D_a\) gives precisely the map of
Theorem~\ref{thm:stableprime371}, by its good traces and
the formula for its ordered characters. Thus
\(\mathfrak q_*\subseteq\mathfrak q\), and
\(J_{\rm cl}\subseteq\mathfrak q_*\) follows from
Proposition~\ref{prop:upfamily372}.
\end{proof}

\begin{corollary}[Classical points of the component at the prescribed weights]
\label{cor:classfiber372}
The fixed-weight algebra of this component for the original
ordered weights of the adjoint is finite over \(O\). Every
geometric characteristic-zero point of it gives a classical
ordinary Hecke system at the common tame level on \(F\).
The adjoint itself is a reduced isolated point in its
characteristic-zero weight fiber.
\end{corollary}
\begin{proof}
Specialization of the finite \(\Lambda_a\)-algebra in
Theorem~\ref{thm:adjcomponent372} gives finiteness over \(O\).
Every one of its points pulls back to the support of the
finite faithful \(\mathbf T_a\)-module of ordinary forms.
Its fiber there is nonzero by Nakayama's lemma. Ordinary
control at this dominant algebraic weight identifies the
characteristic-zero weight space with classical ordinary
forms. The nonzero fiber gives the asserted system, also
for \(r=1\), where the coefficient weight is zero.
~\cite{VII-V-geraghty363}

At the adjoint the weight tangent map of
Lemma~\ref{lem:detangent372} is an isomorphism since its
source and target both have dimension \(3d\). Therefore
the weight increments span the cotangent space of the
local ring there. By Nakayama's lemma they generate its
maximal ideal. Quotienting that local ring by the weight
increments consequently gives its residue field \(J\).
Finiteness of the characteristic-zero fiber makes its
points isolated.
\end{proof}

The corollary concerns the component containing the cuspidal
adjoint. The cotangent quotient in
Theorem~\ref{thm:mixedtrace359} is taken at the split
residual system on the original field. Bounding contributions
from other components and integral relations at their
intersections there remains necessary to control \(C_{\rm ker}\) in
Theorem~\ref{thm:classicaldefect360}.

\section{Recovering classes from three-character traces}\label{nav:chapter-three-character-traces}
\subsection{Three-character traces and the complementary strict group}
\label{sec:traceminors373}

Keep the interval, \(K,\ell,r,\alpha\), and \(L=L_a\) of
Section~\ref{sec:seed361}. Write \(G=G_K\), \(H=G_L\) and
\(\delta_1=\Psi,\ \delta_2=1,\ \delta_3=\Psi^{-1}\).
For \(U=G,H\) let \(T_U\) be the closed \(O\)-algebra of traces
on \(U\) in the original fixed-weight marked matrix ring. Write
\(t(u)\) for its trace, \(t_0(u)=\sum_i\delta_i(u)\) for its
split evaluation and \(\mathfrak a_U=\ker(T_U\to O)\).
These are complete Noetherian local rings. Indeed each is a
quotient of the continuous determinant deformation ring of the
corresponding number field with ramification confined to finitely
many places. The determinant takes values in the trace algebra
by Newton's identities since \(3!\) is a unit. Its universal
coefficient functions generate the trace algebra by compactness.
Noetherianity follows from Chenevier's Proposition 3.7, since
the groups with this ramification restriction satisfy Mazur's
finiteness condition. For an open subgroup the cyclic extensions
of degree \(\ell\) with this restricted ramification are finite
in number by class field theory. This also follows by adjoining
\(\mu_\ell\) and using the finiteness of the class group and of
the group of \(S\)-units modulo \(\ell\)-th powers.~\cite{VII-V-chenevier364}

For \(m\ge1\) let \(I_m(U)\) be the ideal of \(T_U\) generated by
\[
   \det\bigl(t(u_i v_j)\bigr)_{1\le i,j\le m},
           \qquad u_i,v_j\in U .
\]
Rows and columns may be chosen independently. We use linear
extension of \(t\) to \(O[U]\). Allowing finite linear
combinations in each row or column generates the same ideal
by multilinearity. Choose \(g\) in inertia at \(\ell\) in \(H\)
with \(\lambda=\Psi(g)\) such that \(\lambda,1,\lambda^{-1}\)
have distinct residues. This exists by
Proposition~\ref{prop:auxfield361}, since \(\alpha\) is
unramified there and \(2r<\ell-1\). At the universal marked
point the characteristic polynomial on that inertia element is
\(\prod_i(X-\delta_i(g))\), by the fixed ordinary condition.

\begin{lemma}[The ideal defining sums of three characters]
\label{lem:charideal373}
The ideal \(I_4(U)\) is the smallest ideal modulo which \(t\)
is a sum of three characters into the units of the quotient,
labeled by their values on \(g\). Furthermore
\(I_{3+j}(U)\subseteq I_4(U)^j\) for \(j\ge1\).
\end{lemma}
\begin{proof}
Put \(s_i=\delta_i(g)\). Form
\(e_i=\prod_{j\ne i}(g-s_j)/(s_i-s_j)\) in \(O[U]\).
Their matrix images are orthogonal idempotents adding to one,
by Cayley--Hamilton and polynomial interpolation. Each image
has rank one. Over the local marked ring its direct summand
is free, and its rank is read off at the split residual point.
Thus \(t(e_i e_j)\) is one for \(i=j\) and zero otherwise.
Set \(d_i(a)=t(e_i a)\) for \(a\in O[U]\). The determinant with
rows \(e_1,e_2,e_3,a\) and columns \(e_1,e_2,e_3,b\) is
\[
                  F(a,b)=t(ab)-\sum_i d_i(a)d_i(b).
\]
In any quotient by \(I_4(U)\) this vanishes for all \(a,b\).
Apply the identity to \(a=e_i h,\ b=k\). Cyclicity of trace
and \(e_i e_j=0\) for \(i\ne j\) give
\(d_i(hk)=d_i(h)d_i(k)\) there. Also \(d_i(1)=1\), so the
group inverse gives an inverse for \(d_i(h)\). Since the
idempotents sum to one, \(t=\sum_i d_i\), with
\(d_i(g)=s_i\). If in any quotient \(t=\sum_i\psi_i\) for
characters, the matrix of pairings factors through three
columns with entries \(\psi_i(u)\). Hence every
four-by-four minor vanishes in that quotient. If the
characters are labeled by their distinct values \(s_i\) on
\(g\), interpolation also gives \(\psi_i(h)=d_i(h)\).

For the last assertion write each matrix \((t(u_i v_j))\)
as the sum of the matrix \((\sum_k d_k(u_i)d_k(v_j))\)
and the matrix \((F(u_i,v_j))\) over \(T_U\).
The first factors through a free module of rank three and
each entry of the second belongs to \(I_4(U)\).
In the expansion of an \(m\)-row determinant by multilinearity
in its columns, a term using more than three columns from the
first matrix is zero by the alternating property. Every surviving
term for \(m=3+j\) uses at least \(j\) entries from \(I_4(U)\).
\end{proof}

\begin{proposition}[The augmentation ideal on the original quadratic field]
\label{prop:origideal373}
In the original fixed-weight trace ring,
\[
                          I_4(G)=\mathfrak a_G .
\]
In particular the quotient defining sums of three characters
on \(K\) is \(O\), including its integral scheme structure.
\end{proposition}
\begin{proof}
All the minors vanish under the split evaluation. In
\(A=T_G/I_4(G)\) the preceding lemma gives characters
\(d_i\) lifting \(\bar\delta_i\). The residual labeling on
every element follows by applying the projector formula to
the split residual matrices. Each character is continuous by
that same trace formula. We show that
\(d_i/\delta_i\) is everywhere unramified.

At a place over \(\ell\) choose a separating inertia element
\(g_u\), as above. The residual cyclotomic character is
surjective on inertia there, since \(K\) is unramified at
\(\ell\). The polynomial of \(g_u\) in the marked ring has
the three prescribed roots. It also equals
\(\prod_i(X-d_i(g_u))\) in \(A[X]\), by the traces of its
first three powers and Newton's identities. For fixed \(i\)
the two factors \(d_i(g_u)-\delta_j(g_u)\) with \(j\ne i\)
are units. Substitution therefore gives
\(d_i(g_u)=\delta_i(g_u)\). For any inertia element \(h\)
the ordinary upper triangular matrices give
\(t(g_u^n h)=\sum_i\delta_i(g_u)^n\delta_i(h)\) for
\(n=0,1,2\). The character decomposition gives the same
equations with \(d_i(h)\) on the right. The Vandermonde
determinant is a unit, proving equality on all of inertia.
This argument applies at each place of \(K\) above \(\ell\).

At any other finite place inertia acts in the marked problem
through a quotient of the finite image of \(\alpha\), of
order prime to \(\ell\). On the prescribed inertia kernel
the matrix is one, so every \(d_i\) is one there, by its
projector formula. The continuous character \(d_i/\delta_i\)
has values in \(1+\mathfrak m_A\), a pro-\(\ell\) group because its images modulo the powers of \(\mathfrak m_A\) have
successive quotients in additive groups of residual
characteristic \(\ell\). Thus its restriction to the
prime-to-\(\ell\) inertia quotient is trivial. The resulting
everywhere unramified character factors through the class
group of the imaginary quadratic field \(K\). Its order is
prime to \(\ell\) on the intervals under consideration.
It follows that \(d_i=\delta_i\) on \(G\). Every trace in
\(A\) is now its split scalar value. Since these traces
topologically generate and \(O\) is compact, the split
evaluation identifies \(A\) with \(O\).
\end{proof}

\begin{theorem}[Four-by-four minors and the strict group]
\label{thm:minorstrict373}
Let \(D=\operatorname{Frac}(O)/O\) and let \(C\) be the dual of
\(S(\Psi^{-1})\) into \(D\), as in Theorem~\ref{thm:mixedtrace359}. Under
its map \(\Phi:T_G\to O\oplus C\), the ideal \(I_4(H)\) already
maps onto \(C\). Every minor of size five or larger maps to
zero. A finite set of four-by-four minors suffices. The rows
and columns can be chosen at Frobenius elements of good rational
primes split completely in \(L\), with a choice of lifts.
\end{theorem}
\begin{proof}
All four-by-four minors have zero split evaluation. The case
\(C=0\) follows at once. Otherwise use the mixed representation
before its final change of frame in
Theorem~\ref{thm:mixedtrace359}. It has coefficient ring
\((O\oplus C)[y]/(y^2)\), with \(C^2=0\), and its traces
identify the square-zero summand with \(yC\).
Set \(N=H\cap\ker\Psi\) on the absolute Galois group.
The quotient of \(G_{\QQ}\) by \(N\) is abelian, since both
\(L/\QQ\) and the character \(\Psi\) are abelian. On \(N\)
the lower coordinate \(z:N\to C\) and upper coordinate
\(b:N\to O\) are additive. Their needed values persist on
this exact kernel. Indeed after dividing \(C\) by the span
of \(z(N)\), the normalized cocycle factors through the
abelian quotient. Commuting with \(g_*\), where \(z(g_*)=0\)
and \(\Psi(g_*)^{-1}-1\) is a unit, makes it zero.
All the original values span \(C\), so \(z(N)\) spans \(C\).
If the reduction of \(b\) vanished on \(N\), it too would
factor through the abelian quotient. There the central
\(g_*\) acts by \(\bar\eta(g_*)\bar\Psi(g_*)\ne1\), so
every cocycle is a coboundary by the same commutation identity.
This contradicts Lemma~\ref{lem:upperclass359}. Choose
\(k\in N\) with \(b(k)\) a unit.

For \(h\in N\) put \(u=h(chc^{-1})\) and
\(v=k(ckc^{-1})^{-1}\), where \(c\) is complex conjugation.
On \(N\) the action of \(c\) fixes the lower coordinate and
negates the upper one, by the cocycle law and the two
characters. The split matrices on \(N\) are the identity.
Consequently \(M(u)-1\) has first-degree term \(2z(h)X\)
and \(M(v)-1\) has first-degree term \(2yb(k)Y\), in the
degree filtration by the ideal \(C+(y)\).
Its cube is zero. The difference of the two products of
these matrix increments, in the two possible orders, is
therefore \(4yz(h)b(k)[X,Y]\).

Set \(a=u-1,\ b_1=(v-1)g,\ a'=v-1,\ b_1'=(u-1)g\) in
\(O[H]\). The split matrix of each of these four elements
is zero. Thus each of the numbers \(\Phi(d_i(a))\),
\(\Phi(d_i(b_1))\), and their primed counterparts lies in
\(yC\). The products of any two vanish. The four-by-four
determinants \(F\) from the lemma consequently give
\[
 \Phi\bigl(F(a,b_1)-F(a',b_1')\bigr)
   =4yz(h)b(k)\bigl(2-\lambda-\lambda^{-1}\bigr).
\]
Indeed \([X,Y]=\operatorname{diag}(-1,2,-1)\), and multiplying
a degree-two term by \(M(g)\) uses only its split diagonal.
The scalar multiplying \(yz(h)\) is a unit, since
\(2-\lambda-\lambda^{-1}=-(\lambda-1)^2/\lambda\).
As \(h\) varies these differences span \(yC\). They are
linear combinations of actual four-by-four minors on \(H\).
The containment of the larger minors in \(I_4(H)^2\)
follows from Lemma~\ref{lem:charideal373} and expansion by
minors, so they vanish after \(\Phi\).

Finiteness of \(C\) permits a finite generating list. The
image of a four-by-four minor is a continuous function of
its eight group arguments into the finite discrete group
\(C\). Compactness gives one open normal subgroup of \(H\)
such that the function depends only on the eight residues.
For each entry Chebotarev over \(L\) gives good degree-one
prime ideals with Frobenius in its conjugacy class in the
finite quotient. Primes of higher residue degree have
density zero. Choosing the lift at each prime realizes
the specified element of the class. Since \(L/\QQ\) is
Galois, the underlying unramified rational primes split
completely there. They can be taken pairwise distinct and
outside any prescribed finite set. The traces in a minor
are traces of products of these chosen group elements.
\end{proof}

\begin{corollary}[A quotient of the integral reducibility ideal]
\label{cor:minororder373}
For every auxiliary field \(L_a\) of the construction there
is an \(O\)-module surjection
\[
     I_4(H)/\bigl(\mathfrak a_H I_4(H)+I_5(H)\bigr)
                         \longrightarrow C .
\]
On the original \(K\) the corresponding quotient is
\(\mathfrak a_G/\mathfrak a_G^2\). The logarithmic order
of its indicated image is
\([k:\mathbf F_\ell]\log\ell\,\len_O S(\Psi^{-1})\).
\end{corollary}
\begin{proof}
The augmentation ideal maps into the square-zero summand
and \(I_5(H)\subset I_4(H)^2\). The theorem gives the
surjection after division by the displayed submodule.
Apply the proposition on \(K\), and use equality of the
lengths of a finite module over \(O\) and its dual into
\(\operatorname{Frac}(O)/O\).
\end{proof}

These formulas use the character of the original \(E_x\)
for each \(x\) on the common interval. They identify the
integral ideal that detects its complementary length even
after restriction to \(L_a\). A uniform upper bound for
the order of the indicated image is still needed in
Section~\ref{sec:missing}.


\subsection{Integral recovery from traces on the auxiliary field}
\label{sec:restrictedtrace374}

On the original fixed-weight marked problem use \(G=G_K\),
\(H=G_{L_a}\), \(T=T_G\), and \(B=T_H\subseteq T\) as in
Section~\ref{sec:traceminors373}. Write \(M\) for the matrix
representation of \(G\) over the universal marked coefficient
ring. The three rank-one idempotents \(e_i\) in that section
are polynomials over \(O\) in \(M(g)\), with \(g\in H\).
The span of matrices below means finite linear combinations.
No division by the index \([G:H]\) will be needed.

\begin{lemma}[The spans of the two matrix images]
\label{lem:imagespan374}
In the universal matrix algebra one has
\[
                 \sum_{u\in G}T\,M(u)=\sum_{h\in H}T\,M(h).
\]
For all \(m\), \(I_m(G)=I_m(H)T\). Both statements hold for
any normal subgroup \(H\) containing the chosen separating
element \(g\), with \(I_m(H)\) formed in its closed trace ring.
\end{lemma}
\begin{proof}
The right-hand span is an algebra since \(H\) is a group.
For \(u\in G\) write \(d_j(u)=\operatorname{tr}(e_jM(u))\).
Its split evaluation is \(\delta_j(u)\), so \(d_j(u)\) is a
unit of the local ring \(T\). Acting on the free rank-one
summand \(e_j\), the endomorphism \(e_jM(u)e_j\) is
multiplication by its trace. Thus
\(e_jM(u)e_j=d_j(u)e_j\).
Normality gives \(M(u)e_jM(u)^{-1}\) in the \(O\)-span of
\(M(H)\), since it is the same polynomial evaluated at
\(M(ugu^{-1})\). For every \(i,j\) it follows that
\[
  e_i\bigl(M(u)e_jM(u)^{-1}\bigr)e_j
             =d_j(u^{-1})\, e_i M(u)e_j .
\]
The matrix on the left belongs to the right-hand span.
Division by \(d_j(u^{-1})\) puts every block \(e_i M(u)e_j\)
in that span. Summing the nine blocks gives \(M(u)\).
This proves equality of spans. Apply these expressions to
each row and column of a trace determinant on \(G\).
Multilinearity of the determinant and linearity of trace
express it as a finite \(T\)-linear combination of the
corresponding determinants on \(H\). The reverse containment
follows from \(H\subseteq G\).
\end{proof}

\begin{theorem}[Restriction recovers the entire fixed-weight trace ring]
\label{thm:restrictedtrace374}
For every \(a\) on the chosen interval,
\[
            T_{G_{L_a}}=T_{G_K},\qquad
            I_4(G_{L_a})=\mathfrak a_G .
\]
The equalities hold in the original integral ring. The
\(O\)-algebra generated topologically by traces at good
primes of \(K\) lying above rational primes split completely
in \(L_a\) is already \(T\), if both choices over \(K\) are
allowed. Every characteristic-zero specialization of the
marked matrices has the same invariant subspaces on \(G_K\)
and on \(G_{L_a}\), after any extension of its coefficient field.
\end{theorem}
\begin{proof}
Let \(\mathfrak m_B,\mathfrak m_T\) denote the two maximal
ideals. The rings have the same finite residue field, since
both contain \(O\) and the residual matrices have diagonal
characters defined over its residue field. The split evaluation
maps \(T\) onto \(O\). If \(\varpi\) is a uniformizer, then
\(\mathfrak m_T=\varpi T+\mathfrak a_G\).
The preceding lemma and Proposition~\ref{prop:origideal373}
give
\[
                 \mathfrak a_G=I_4(H)T,\qquad
                 \mathfrak m_T=\mathfrak m_B T .
\]
For the second equality, \(I_4(H)\) has zero split value.
Consequently \(\mathfrak m_T^n=\mathfrak m_B^nT\) for
every \(n\). An element of this power modulo
\(\mathfrak m_T^{n+1}\) is a sum of elements of
\(\mathfrak m_B^n\) times scalars in the common residue field.
Lifting the scalars to \(O\) represents it by an element of
\(B\). Induction starting in the residue field gives the
surjectivity of \(B\to T/\mathfrak m_T^n\) for every \(n\).
The ring \(B\) is compact and \(T\) is Hausdorff with this
topology. Hence \(B=T\), which also removes the scalar
extension in the equality of the two ideals.

For any \(h\in H\), use continuity modulo
\(\mathfrak m_T^n\) and Chebotarev over \(L_a\) to approximate
its trace by that of a degree-one prime there outside any
finite exceptional set. Its rational prime splits completely
in \(L_a\). The Frobenius on restriction to \(K\) agrees with
the Frobenius over \(L_a\), since the relative local degree
is one. Passing through all \(n\) proves the generation
assertion. Finally specialize the span equality of the
lemma. Every vector subspace invariant under \(M(H)\) is
preserved by its scalar-linear span, hence by \(M(G)\).
This remains true over any extension of the coefficient field.
\end{proof}

\begin{proposition}[An integral finite quotient on the original field]
\label{prop:orighecke374}
Let \(R_a\) be the ordinary polarized determinant ring of
Section~\ref{sec:basesupport371}, using the fixed cyclotomic
normalization there. Work over a common finite coefficient
DVR containing the chosen adjoint eigenvalues if needed.
Restriction gives a surjection \(R_a\to T\) sending each
inertial weight to its prescribed fixed value. If
\(\mathfrak q_*\) is the minimal prime through the adjoint
from Theorem~\ref{thm:adjcomponent372}, put
\(\mathfrak b_a=\mathfrak q_* T\). Then \(T/\mathfrak b_a\)
is a nonzero finite \(O\)-algebra. Its geometric
characteristic-zero points give classical ordinary Hecke
systems on the definite group over \(F_a^+\).
\end{proposition}
\begin{proof}
All the preceding trace-ring calculations persist after a
finite coefficient extension. It is finite free on the
coefficient modules, so tensoring preserves the embedding
of the trace ring in the marked ring. Traces with coefficients
in the extension generate exactly this scalar extension by
closure and by a basis of the extension.

Restrict the universal marked matrices to \(G_{L_a}\). They
kill the prescribed inertia kernel away from \(\ell\).
Indeed on inertia at those places \(\kappa_\ell\) is
unramified and \(\ker\bar\alpha=\ker\alpha\), since
\(\alpha\) has prime-to-\(\ell\) order. The original
polarized matrices on \(G_{F_a^+}\) kill
\(I_u\cap\ker\eta\cap\ker\bar\Psi\), so restriction factors
through \(\Gamma_a^0\). The original polarization supplies
the polarized determinant on this quotient. Its coefficients
are in \(T\) by Newton's identities. This also holds on
group algebras after extension of scalars by applying the
same identities to a generic finite linear combination.

At each chosen place above \(\ell\) the original ordinary
flag restricts to the upper field. Let \(s\) be a separating
inertia element in that completion and form its three
polynomial projectors just as for \(g\). The flag is ordered
by its inertial characters, whose values on \(s\) are known
scalars in \(O\). The trace of the \(i\)-th projector times
\(M(h)\) for a decomposition group element \(h\) is the
\(i\)-th diagonal character of this flag. This follows by
writing the matrices in a basis adapted to the flag, where
the projector is upper triangular with just one diagonal
entry one. Thus these full characters take values in \(T\).
The ordinary word required in \(R_a\) vanishes on the flag
because its factors successively lower its steps. These
facts give a map \(R_a\to T\). On the unit weight group
it has the original prescribed values. Its image is closed
by compactness and contains the traces of all \(h\in G_{L_a}\).
The theorem proves surjectivity.

Let \(\mathfrak w\) be the ideal of the upper weight algebra
generated by the differences between the weight coordinates
and their prescribed values. The map induces a surjection
\[
        (R_a/\mathfrak q_*)/\mathfrak w(R_a/\mathfrak q_*)
                          \longrightarrow T/\mathfrak b_a .
\]
The left side is finite over \(O\) by
Corollary~\ref{cor:classfiber372}. The quotient on the
right has the common residual point, since the original map
is local. Every geometric characteristic-zero point of the
right side pulls back to one on the left, with the indicated
Hecke system by the same corollary.
\end{proof}

\begin{corollary}[Original-field cuspidal labels for irreducible points]
\label{cor:irreduciblelabel374}
Suppose a geometric characteristic-zero point of the original
marked matrix problem factors on traces through \(T/\mathfrak b_a\)
and is absolutely irreducible over \(K\). Its trace occurs
among the rank-three cuspidal labels on \(K\) at the level
\(N_*\) of Theorem~\ref{thm:heckedivisor345}.
\end{corollary}
\begin{proof}
The restricted representation stays absolutely irreducible
by Theorem~\ref{thm:restrictedtrace374}. Its ordered trace
at the upper field is classical by the proposition. The
classical transfer of a definite unitary system in rank three
decomposes its semisimple Galois representation according
to the isobaric summands, including the character twists
of a repeated block. This is the attachment of Geraghty's
Proposition~2.7.2, with the normalization used in
Lemma~\ref{lem:bcpoints371}. Absolute irreducibility forces
a single cuspidal summand of rank three. Its attached
representation is the restriction of the original one
twisted by \(\kappa_\ell^{-1}\), by compatibility at good
primes and equality of semisimple representations with the
same trace.~\cite{VII-V-geraghty363}

Apply cyclic descent through the intermediate fields
between \(K\) and \(L_a\). Over each of them the original
representation is absolutely irreducible since its subgroup
\(G_{L_a}\) already acts irreducibly. This is the restriction
hypothesis used in the proof of
Theorem~\ref{thm:auxdescent361}. In detail the Galois-invariant
cuspidal representation at each prime-degree step descends.
The unique descent with central character one is conjugate
self-dual. Its attached Galois representation agrees on the
upper subgroup, hence differs downstairs by a character of
the cyclic quotient by Schur's lemma. Its cube is trivial
by the determinant, and the step has degree prime to three.
The character is trivial. The crystalline and tame local
conditions at \(K\) are those of the original marked point,
so the level computation of that theorem places the resulting
cuspidal system at \(N_*\).~\cite{VII-V-I-ac123,VII-V-lgp361,VII-V-V-shin}
\end{proof}

The equality of trace rings applies before passing to the
finite quotient. Finiteness and classicality in the proposition
apply after imposing \(\mathfrak b_a\). To obtain the estimate
of Section~\ref{sec:missing} one still needs to control how
much of the image \(C\) of the trace differences is lost
under such a passage.


\subsection{Crystallinity and classical evaluation of trace specializations}
\label{sec:tracecrystal375}

Use \(G=G_K\), \(H=G_{L_a}\) and the original trace ring \(T\)
of Section~\ref{sec:restrictedtrace374}. Denote the complete
local ring of the universal marked matrices by \(A\), its
representation on \(G\) by \(M\), and the \(T\)-span of its
matrix image by \(E_M\).
Thus \(E_M\subseteq M_3(A)\). We show how the local crystalline
condition can be read at a global trace point even when no
choice of marked matrices over its valuation ring is given.

\begin{lemma}[Finite quotients of the algebra of matrices]
\label{lem:torsioncrys375}
The algebra \(E_M\) is finite as a \(T\)-module. Let \(u\mid\ell\)
be a place of \(K\). For every \(n\ge1\), the finite representation
\(E_M/\mathfrak m_T^n E_M\), with \(G_{K_u}\) acting by left
multiplication, is a subquotient of a lattice in a crystalline
representation with Hodge--Tate weights in \([-r,r]\).
\end{lemma}
\begin{proof}
Use the universal degree-three determinant with our residual
trace for the group of \(K\) unramified outside the original
finite set. Its completed Cayley--Hamilton algebra is finite
over its determinant ring by Wang-Erickson's Proposition~3.6.
The required finiteness of residual adjoint cohomology follows
by restriction to the finite extension killing the residual
characters and the finiteness condition verified in
Section~\ref{sec:traceminors373}. The universal determinant
ring surjects onto \(T\). After this scalar extension, sending
group elements to \(M(g)\) defines a map of the finite algebra
to \(M_3(A)\). Indeed continuity gives a map on the completed
group algebra, and the matrices satisfy its Cayley--Hamilton
relations by the ordinary matrix identity. Its image is a finite
\(T\)-module containing \(E_M\). Since \(T\) is Noetherian,
the submodule \(E_M\) is finite.~\cite{VII-V-wangerickson375}

Put \(A_u\) for the power series ring of the selected local
ordinary chart, with maximal ideal \(\mathfrak m_u\). Its
integral field points separate its elements, as proved in
Theorem~\ref{thm:crystalline121} by evaluating convergent
power series on sufficiently small powers of a uniformizer.
Every such representation restricted to \(G_{K_u}\) is
crystalline with the indicated weights. The proof of that
theorem covers the split and unramified quadratic completions
in the present problem. At a conjugate place the dual with
conjugation has the same weight interval.

For each \(s\), the complement of \(\mathfrak m_u^s\) in
\(A_u\) is compact. Each element in it has nonzero image
at some integral field point and therefore modulo some
power of the uniformizer at that point. Kernels of those
finite evaluations are open and closed. Taking a finite
subcover gives evaluations whose common kernel \(J_s\)
is contained in \(\mathfrak m_u^s\). The local representation
over \(A_u/J_s\) embeds by entries in the direct sum of the
corresponding truncated crystalline lattices. Passing to
\(A_u/\mathfrak m_u^s\) gives a quotient. Finite subquotients
of crystalline lattices in a fixed weight interval are closed
under finite sums, submodules and quotients. For example,
the preimage of a subgroup of a finite quotient of a lattice
is again a full lattice. This gives the asserted property
for the local chart modulo each power.

The local map \(A_u\to A\) is continuous. Thus for each \(b\)
the marked module \((A/\mathfrak m_A^b)^3\) on \(G_{K_u}\)
is obtained by extension of scalars from such a finite local
chart, up to the choice of its local basis. As a group module
it is a quotient of a finite direct sum of the chart modules, by multiplying the columns by generators of the finite target
coefficient module. The property consequently holds for
each \((A/\mathfrak m_A^b)^3\).

The finite \(T\)-module \(E_M\) is compact in its
\(\mathfrak m_T\)-adic topology. Its embedding into the
Hausdorff ring \(M_3(A)\) is continuous, as is seen by
choosing finitely many module generators. A continuous
injection from a compact space is a homeomorphism onto its
image here. For each \(n\) there is therefore a \(b\) for
which \(E_M\cap M_3(\mathfrak m_A^b)\subseteq
\mathfrak m_T^n E_M\). Modulo the intersection, \(E_M\) embeds
as a group module in three copies of
\((A/\mathfrak m_A^b)^3\), by its three columns. Taking the
further quotient gives the assertion.
\end{proof}

\begin{proposition}[The local conditions at a global trace point]
\label{prop:tracecrystal375}
Let \(x:T\to O'\) be a continuous local map to the integers
of a finite coefficient field. The associated semisimple
geometric representation \(\rho_x\) is crystalline at every
place of \(K\) above \(\ell\), with labeled cyclotomic exponents
\(r,0,-r\). At all other finite places it has the prescribed
finite inertia. Its determinant is one.

More generally, for a continuous determinant lift to a
finite local characteristic-zero coefficient algebra with
absolutely irreducible reduction at \(x\), the associated
matrix representation, regarded over \(\QQ_\ell\), is
crystalline at these same places.
\end{proposition}
\begin{proof}
Enlarge the finite coefficient field to split the semisimple
representation associated to the determinant, using
Chenevier's theorem over fields. It is a module for \(E_M\).
In fact a relation \(\sum_j a_j M(g_j)=0\) in \(E_M\) gives
\(\sum_j a_j t(g_j h)=0\) for every \(h\in G\). After applying
\(x\), the corresponding matrix pairs to zero against the
whole image algebra under the trace pairing. Over a splitting
field that algebra on a semisimple representation is a product
of matrix algebras. The trace on each factor is its ordinary
matrix trace times the multiplicity, a nonzero integer in
characteristic zero. Its pairing is nondegenerate. Thus every
such relation acts by zero.

For the assertion with infinitesimal coefficients use the
matrix lift given by Chenevier's Theorem~2.22(i), after
extending coefficients to split its reduction. Here the
determinant lift is along a continuous map from \(T\).
At the irreducible reduction choose nine group elements
whose matrices span \(M_3\). The trace Gram determinant
on these elements is a unit also in the finite local
coefficient algebra. They form a basis there by Nakayama's
lemma, and the same trace argument shows that every relation
in \(E_M\) acts by zero. The matrix coordinates of group
elements are continuous since they are recovered from this
basis and its invertible Gram matrix.~\cite{VII-V-chenevier364}

In either case the action of \(E_M\) is continuous because it is
\(T\)-linear for the specified scalar map on a finite
\(T\)-module. In the underlying finite-dimensional
\(\QQ_\ell\)-space take a full \(\ZZ_\ell\)-lattice and its
span under \(E_M\). The latter contains the former and is
bounded by compactness of the acting algebra, so it too
is a full lattice \(D\), stable under \(E_M\). For each
\(m\), the action of \(E_M\) on \(D/\ell^mD\) factors
through \(E_M/\mathfrak m_T^n E_M\) for some \(n\), by
continuity. Sending the unit to each of finitely many
additive generators expresses \(D/\ell^mD\) as a quotient
of a finite direct sum of these left regular modules.
The lemma puts every one of these truncations in the same
bounded crystalline weight interval. Liu's theorem on limits
of torsion crystalline representations now implies that
\(D[1/\ell]\) is crystalline with weights in that interval.
One can twist the interval to \([0,2r]\) in the convention
of that theorem. The prime is odd.~\cite{VII-V-liu375}

At a field point the local trace on the full decomposition
group is the sum of the three ordered flag characters,
each recovered in \(T\) by the local projector formula
of Proposition~\ref{prop:orighecke374} on \(K\) itself.
The characters on inertia are
\(\kappa_\ell^r,1,\kappa_\ell^{-r}\). Equality of traces
gives these three local composition factors up to an
unramified twist of each. Exactness of the de Rham functor
with its filtration on crystalline representations gives
their three labeled weights. Away from \(\ell\) the
prescribed inertia kernel acts trivially in \(E_M\).
The remaining inertia group has order prime to \(\ell\)
and its character is the prescribed one, as follows either
from reduction or from the universal traces on that finite
group. Finally the determinant of the marked representation
is one on inertia at \(\ell\) by the weights, and away from
\(\ell\) by the finite inertial characters. It is everywhere
unramified on \(K\) and takes values in principal units.
The class number is prime to \(\ell\), so it is one. The
same identity passes to the trace determinant.
\end{proof}

\begin{theorem}[Every non-self-dual point of the finite quotient]
\label{thm:alllabels375}
Every reducible geometric characteristic-zero point of \(T\)
given by a continuous integral specialization has self-dual
trace. Every geometric characteristic-zero point of
\(T/\mathfrak b_a\) where an inverse-trace difference is
nonzero occurs among the cuspidal labels on the original
\(K\) at level \(N_*\). This includes the first twist \(r=1\).
\end{theorem}
\begin{proof}
Work with the global semisimple representation \(\rho_x\)
of the proposition. If it has three character factors,
all four-by-four trace minors vanish, so its evaluation
is the split evaluation by Proposition~\ref{prop:origideal373}.
Otherwise a reducible representation has a unique
one-dimensional summand \(\vartheta\) and an irreducible
two-dimensional summand. The identity
\(t(chc^{-1})=t(h^{-1})\) holds in the original trace ring.
It forces \(\vartheta^c=\vartheta^{-1}\) by uniqueness
of the one-dimensional summand. Its residual character must
be one of \(\bar\Psi,1,\bar\Psi^{-1}\), by reduction of a
stable lattice and uniqueness of the residual semisimplification.
Each of these is fixed by \(c\) because it comes from \(\QQ\).
The only one equal to its inverse is \(1\), by separation
on inertia. Thus \(\bar\vartheta=1\). At a place over
\(\ell\) its restriction to inertia is one of the three
prescribed characters, and their distinct reductions force
it to be trivial. Away from \(\ell\) the prime-to-\(\ell\)
inertia image and its trivial reduction give the same
conclusion. The unramified character \(\vartheta\) takes
values in principal units and the class number of \(K\)
is prime to \(\ell\), so \(\vartheta=1\). The plane has
determinant one by the proposition. Its alternating
determinant pairing identifies it with its dual. This
proves the first assertion.

A point of the finite \(O\)-algebra \(T/\mathfrak b_a\) with
values in an algebraic closure of \(\operatorname{Frac}(O)\)
has integral image in a finite extension. The evaluation
is continuous and local on valuation rings since its
reduction lies over the unique maximal ideal of this
finite local algebra. If an inverse-trace difference
survives, \(\rho_x\) is absolutely irreducible by the first
assertion. The representation factors through \(E_M\) and
Lemma~\ref{lem:imagespan374} shows that the specialized
matrix spans on \(G\) and \(H\) agree, even for a trace
point. It therefore remains irreducible on \(H\) and on
all intermediate subgroups. Proposition~\ref{prop:orighecke374}
gives a classical system on \(F_a^+\). Its transfer is
cuspidal by the decomposition of the attached representation
used in Corollary~\ref{cor:irreduciblelabel374}. The cyclic
descent in that corollary only requires irreducibility on
the subgroups, the determinant and the polarized trace.
For the descended level and weight, the preceding proposition
supplies precisely the crystalline type on \(K\) and the
finite inertial type used in
Theorem~\ref{thm:auxdescent361}. Hence the same original
level \(N_*\) applies to this trace point.
\end{proof}

\Needspace{6\baselineskip}
\begin{corollary}[The remaining length as a kernel of evaluation]
\label{cor:componentdefect375}
For each triple over the unramified coefficient ring \(O\)
of Section~\ref{sec:classicaldefect360}, let \(P_a\) be the
kernel of the map from its original \(T\) to the reduced
generic fiber of the finite quotient in
Proposition~\ref{prop:orighecke374}, contracting from the
common coefficient extension if one was made. Put
\(C_a=\Phi(\mathfrak d\cap P_a)\), with the original
inverse-trace ideal and mixed map. Uniformly on the same
intervals and fields,
\[
 0\le \mathcal L_{>}(x)
       -\sum_{\ell>P}(\log\ell)\!
        \sum_{1\le r<t(x)}\ \sum_{[\alpha]}
           b_\alpha\,\len_O C_a
       \le \log J_* \ll_{\psi,\eps}V^{C_1}.
\]
\end{corollary}
\begin{proof}
By definition \(P_a\) is the intersection of the kernels of
all characteristic-zero geometric evaluations of this finite
quotient, restricted to \(T\). At any such point either the
whole ideal \(\mathfrak d\) vanishes or its evaluation comes
from the original list \(X\), by the theorem. Therefore an
element of \(\mathfrak d\) killed by all of \(X\) belongs to
\(P_a\). This proves \(C_{\rm ker}\subseteq C_a\subseteq C\).
Apply Theorem~\ref{thm:classicaldefect360} and additivity of
length. Taking the contraction before applying \(\Phi\) keeps
the original residue degrees \(b_\alpha\).
\end{proof}

The calculation identifies the evaluation ideal on both
fields. The sum of the lengths of \(C_a\) still requires a
uniform upper bound. In the archive, the passage from an
integral deformation model to a bound on one character
uses the finiteness and point hypotheses in
\srcloc{Development/AutomorphicArithmetic/AutomorphicVarieties/UnitaryThreeGroups/ModelUpper}{104}
and the automorphy input in
\srcloc{Development/AutomorphicArithmetic/AutomorphicVarieties/UnitaryThreeGroups/OneCharOf}{158}.


\subsection{The inverse trace at reducible characteristic-zero points}
\label{sec:reduciblefiber376}

Write \(E_M,T,G,H\) as in Section~\ref{sec:tracecrystal375}.
An inverse-trace difference is \(t(h)-t(h^{-1})\) with \(h\in G\).
At a reducible point Theorem~\ref{thm:alllabels375} shows that
all these differences have value zero in the residue field.
We determine their values in finite local algebras supported
at such a point. This requires control of extensions between
the distinct irreducible constituents.

\begin{lemma}[Extensions allowed by the integral matrix algebra]
\label{lem:matrixmodules376}
Let \(D\) be a finite local algebra over a finite extension
\(J\) of \(\QQ_\ell\), with residue field \(J\), and give it
a continuous map \(T\to D\) whose reduction \(x\) is an
integral trace point. Enlarge \(J\) to split \(\rho_x\) when
necessary. In \(E_D=E_M\otimes_T D\) the simple modules are
exactly the constituents of \(\rho_x\). Every finite
\(E_D\)-module, restricted to \(G\), is crystalline above
\(\ell\) and kills the prescribed inertia kernels elsewhere.
For two nonisomorphic simple modules \(U,V\), a nonsplit
extension of \(U\) by \(V\) in this algebra gives a nonzero
element of \(H_f^1(K,\operatorname{Hom}_J(U,V))\).
\end{lemma}
\begin{proof}
Modulo the maximal ideal of \(D\), the algebra maps onto
the image algebra of the semisimple \(\rho_x\) by
Proposition~\ref{prop:tracecrystal375}. Its image is a product
of full matrix algebras on the distinct simple constituents.
The kernel consists of nilpotent elements. Indeed for a
kernel element \(b\) the traces of \(b,b^2,b^3\) vanish.
The matrix identity of degree three expresses \(b^3\) in
terms of these traces by Newton's identities, so \(b^3=0\).
That identity persists on coefficient extension by applying
it first with indeterminate coefficients to a finite sum of
matrices. An ideal consisting of nilpotents is in the
Jacobson radical, since \(1-ab\) is a unit whenever \(b\)
belongs to the ideal. The quotient is semisimple, so the
kernel is exactly that radical. Lifting back through the
nilpotent maximal ideal of \(D\) gives the asserted simples.

On a finite \(E_D\)-module the action of the compact \(E_M\)
is continuous, by module finiteness over \(T\). As in the
proof of Proposition~\ref{prop:tracecrystal375}, its span of
a full \(\ZZ_\ell\)-lattice is a full stable lattice. Every
reduction of that lattice is a quotient of finitely many
left regular modules \(E_M/\mathfrak m_T^nE_M\). Applying
Lemma~\ref{lem:torsioncrys375} and Liu's limit theorem again
gives crystallinity of the module. At the other places the
specified inertia matrices are the identity already in \(E_M\).

A splitting of an extension as a \(G\)-module gives a unique
\(G\)-stable complement isomorphic to \(U\), because
\(\operatorname{Hom}_G(U,V)=0\). The central action of \(D\)
preserves this complement by the same vanishing. Since \(D\)
and the group matrices generate \(E_D\), this is an
\(E_D\)-splitting. Thus a nonsplit extension gives a nonzero
Galois cohomology class. Its crystalline extension at
\(\ell\) becomes split over \(B_{\rm cris}\), so it satisfies
the Bloch--Kato finite condition there. Away from \(\ell\)
the finite inertial image acts semisimply in characteristic
zero. Hence the class is unramified there.~\cite{VII-V-liu375,VII-V-VI-bk117}
\end{proof}

\begin{proposition}[An odd descent of the two-dimensional constituent]
\label{prop:odddescent376}
Suppose \(\rho_x=1\oplus P\) with \(P\) absolutely
irreducible of dimension two. Then \(P\) extends to an odd
representation \(W\) of \(G_\QQ\) with \(\det W=\eta\).
A finite coefficient extension is allowed. There is a
classical cuspidal eigenform of weight \(2r+1\) whose
arithmetic Galois representation is \(W\kappa_\ell^r\).
Moreover,
\[
                        H_f^1(K,P)=H_f^1(K,P^\vee)=0.
\]
\end{proposition}
\begin{proof}
The marked symmetry at \(c\) induces the \(T\)-linear
anti-involution \(B\mapsto B^\dagger=J_3B^{\mathsf T}J_3\)
on \(E_M\), sending \(M(h)\) to \(M(ch^{-1}c^{-1})\).
It preserves the trace and hence the radical of the trace
pairing at \(x\). It descends to the product
\(J\times M_2(J)\), the image algebra of \(1\oplus P\).
It preserves both factors since their dimensions differ.
On the plane it is adjunction by either a symmetric or an
alternating nondegenerate form. Indeed composing with
transpose gives an automorphism of the matrix algebra,
which is inner. The square of the involution shows that
the representing matrix is either symmetric or skew.

We verify the sign on this plane. For any two skew-adjoint
three-by-three matrices \(X,Y\) with respect to a nondegenerate
symmetric form one has
\[
                        2XYX=\operatorname{tr}(XY)X .
\]
In an orthonormal basis write \(Xv=a\times v\), \(Yv=b\times v\).
Then \(XY=ba^{\mathsf T}-(a\cdot b)I\),
\(a^{\mathsf T}X=0\), and \(\operatorname{tr}(XY)=-2a\cdot b\).
This proves the identity in characteristic zero. For the
fixed integral \(J_3\), it is a polynomial identity in
entries of \(Z-Z^\dagger\) and \(Z'-Z'^\dagger\) with integral
coefficients. Its validity over \(\CC\) for all entries
therefore proves it in the marked ring even if that ring
has torsion. Every skew element of the semisimple image
is the image of a skew element after scalar extension,
by applying \((1-\dagger)/2\). The identity thus holds on
the plane, with its ordinary matrix trace since its
multiplicity is one. For alternating adjunction all traceless
two-by-two matrices are skew. Take \(X=\operatorname{diag}(1,-1)\)
and \(Y=E_{12}+E_{21}\) on that plane, zero on the line.
Then \(XYX=-Y\) and \(\operatorname{tr}(XY)=0\), a contradiction.
The plane form has a symmetric matrix \(S\).

Let \(J_0\) be the standard alternating matrix in dimension
two. Since \(\det P=1\), the conjugation formula on \(G\)
can be written
\[
 P(chc^{-1})=(S^{-1}J_0)P(h)(S^{-1}J_0)^{-1}.
\]
The matrix \(S^{-1}J_0\) has trace zero, as the product of
a symmetric and a skew matrix. Its square is a nonzero scalar.
After rescaling, it has square one and determinant minus one.
Assign it to \(c\). This extends \(P\) to the required \(W\).

The representation is finitely ramified and absolutely
irreducible. It is crystalline at \(\ell\) with weights
\(r,-r\), by Proposition~\ref{prop:tracecrystal375} and
unramified descent from \(K_u/\QQ_\ell\). For this descent
the crystalline periods over an unramified extension carry
semilinear descent data on the full coefficient field of
periods. Their invariants have the same dimension downstairs.
Its residual semisimplification on \(G\) consists of
\(\bar\Psi,\bar\Psi^{-1}\). If a residual constituent over
\(\QQ\) had dimension two, its restriction to the index-two
normal subgroup would be semisimple: the translates of
any simple submodule span it. The two character lines
would each be fixed by \(c\), since both characters come
from \(\QQ\). This contradicts irreducibility. Thus the
residual representation over \(\QQ\) is reducible after
semisimplification. Pan's Theorem~1.0.2 applies to this
odd regular de Rham representation, with \(\ell>3\).
Twisting to weights \(0,2r\) gives a form of weight \(2r+1\).
Any further character twist of Hodge--Tate weight zero is
finite order: its inertia at \(\ell\) is finite by the
classification of de Rham characters, all other inertia is
finite by local reciprocity, and a finite ray class quotient
then supplies the character. Absorb it into the eigenform.
~\cite{VII-V-VI-pan40}

This form is non-CM. Otherwise \(P\) becomes a sum of
characters over an open normal subgroup. A scalar restriction
would contradict its two distinct Hodge--Tate weights. If the
characters differ their two lines are permuted transitively
by \(G\). The kernel of \(\bar\Psi\) has pro-\(\ell\) matrix
image, because it is residually unipotent in an integral
upper triangular residual basis. The permutation quotient
of order two thus factors through the cyclic image of
\(\bar\Psi\). A lift of a generator swaps the lines and
has trace zero, so a generator \(a\) of the residual
image satisfies \(a+a^{-1}=0\). Its order would be four,
whereas on inertia \((\ell-1)/\gcd(r,\ell-1)>4\).

With the arithmetic convention for \(V_f=W\kappa_\ell^r\),
Kato's \(V_\lambda(f)\) is \(V_f^\vee\) by his Euler
polynomial at good primes. Therefore \(V_\lambda(f)(r)\)
restricts to \(P^\vee\simeq P\) on \(G_K\).
Theorem~14.2(1) of Kato applies over the abelian field \(K\),
since \(1\le r\le k-1\) and \(2r\ne k=2r+1\). It gives
finiteness for the Selmer group with divisible torsion
coefficients. The image of the vector space of finite
Selmer classes in that finite group is divisible, hence zero.
It is consequently contained in the image of the finitely
generated integral cohomology of a lattice over a fixed finite
ramification set. Such an image contains no nonzero vector
subspace. The integral finite generation and the compatibility
of the local maps here are the ones proved in
Theorem~\ref{thm:fullad120}.~\cite{VII-V-VI-kato120}
\end{proof}

\begin{theorem}[Symmetry on every infinitesimal reducible fiber]
\label{thm:reduciblefiber376}
Let \(T\to D\) be as in Lemma~\ref{lem:matrixmodules376}.
If \(\rho_x\) is reducible, all inverse-trace differences
vanish in \(D\). If \(\rho_x\) is the split triple, the whole
map \(T\to D\) factors through the original split evaluation.
\end{theorem}
\begin{proof}
We recall a finite algebra calculation. Write \(R=E_D\)
and let \(J_R\) be its nilpotent Jacobson radical.
Over the splitting field \(R/J_R\) is a product of matrix
algebras. The block of \(J_R/J_R^2\) from one simple factor
to another can be nonzero only when there is a nonsplit
extension of the corresponding simple modules, with the
target as submodule. To see this, lift a primitive idempotent
\(f\) for the source factor through the nilpotent radical.
The projective module \(Rf\) has simple top and radical
\(J_R f\). A simple quotient of \(J_R f/J_R^2 f\) gives,
by pushing out \(Rf/J_R^2 f\), a length-two extension with
this quotient as submodule. It is nonsplit, since the
radical has nonzero image on the resulting module. Matrix
units in \(R/J_R\) recover the whole block from its primitive
columns. Finally \(J_R^n/J_R^{n+1}\) is an image of the
\(n\)-fold tensor power of \(J_R/J_R^2\) over \(R/J_R\),
by multiplication. Each block in it is spanned by products
that pass through \(n\) successive nonzero blocks. These
facts also hold in the presence of diagonal radical blocks.

First let the constituents be \(1,P\). Extensions between
the two distinct factors are zero in both directions by
the lemma and the proposition. The same holds for blocks
between them in every \(J_R^n/J_R^{n+1}\), by the product
description. Use the idempotent \(e\) for the eigenvalue
one of the separating \(g\) of Section~\ref{sec:traceminors373}.
It reduces to the identity on the line and zero on the plane.
The off-diagonal spaces \(eR(1-e)\) and \((1-e)Re\) vanish
modulo the radical and in every successive layer, so are
zero by nilpotence. Thus \(e\) is central. In the original
matrix ring \(eM(h)e=t(eM(h))e\) and \(t(e)=1\).
Consequently \(eRe=De\), and the action on it is given by
a character \(\theta:G\to D^\times\) reducing to one.

At each place above \(\ell\) the middle spectral idempotent
of a local separating inertia element also reduces to
identity on the line and zero on the plane. To see that it equals \(e\), use the product \(eR\times(1-e)R\). The unwanted idempotent
and the complement of the wanted one lie in nilpotent
ideals and hence are zero. The ordinary projector formula
therefore gives \(\theta(h)=1\) on that inertia. Away from
\(\ell\) its image on inertia is finite and lies in
\(1+\mathfrak m_D\), a torsion-free group by the finite
logarithm in characteristic zero. It is everywhere
unramified and hence factors through the finite class
group, so \(\theta=1\). The degree-three Cayley--Hamilton
identity on \(M(h)\), with determinant one, has polynomial
\[
               X^3-t(h)X^2+t(h^{-1})X-1 .
\]
Multiply by \(e\). The resulting equation at \(X=1\)
gives \(t(h^{-1})=t(h)\), since \(t(e)=1\).

For the split triple order the simple characters by their
cyclotomic exponents. There is no extension in \(R\) from
a larger to a smaller exponent. Indeed its class would lie
in the finite Selmer group with field coefficients for a
negative power of \(\Psi\), by the lemma. Restrict to the
finite abelian extension obtained by taking the compositum of \(K\) with the field fixed by \(\ker\alpha\).
There the coefficients are a negative cyclotomic twist
\(J(-m)\), \(m>0\). The local finite group at \(\ell\)
is zero by the Bloch--Kato dimension formula. At other
finite places the unramified first cohomology is zero
because Frobenius on the coefficients is not one.
The resulting rational class is therefore everywhere
locally trivial. Its image with divisible torsion coefficients
lies in the finite strict group bounded by
Proposition~\ref{prop:finite15}, whose proof applies to
this fixed number field and \(m\). A divisible image in
that finite group is zero, and finite generation of integral
cohomology excludes a nonzero rational subspace in the
kernel as above. Restriction in rational cohomology is
injective by corestriction, proving the claimed vanishing.

It follows from the radical-layer calculation that between
distinct idempotents only the increasing direction can
occur anywhere in \(R\). The three corners cut out by the
spectral \(e_i\) are scalar, by their original rank-one
identity. In a product \(M(h)M(k)\) the \(i,i\) corner can
therefore only receive the product of the \(i,i\) corners, since a path leaving \(i\) and returning would use both directions.
Thus each \(h\mapsto t(e_iM(h))\) is multiplicative and
the trace in \(D\) is the sum of three characters.
Proposition~\ref{prop:origideal373} forces the entire map
to factor through \(O\).
\end{proof}

\begin{corollary}[The rational support of the inverse-trace ideal]
\label{cor:rationalsupport376}
In the generic fiber of the finite quotient
\(T/\mathfrak b_a\), the inverse-trace ideal has zero
projection to every local factor with reducible geometric
residue representation, including all its nilpotents.
A geometric local factor supported at the split triple
is its residue field.
\end{corollary}
\begin{proof}
The generic fiber of the finite algebra is Artinian.
After a finite splitting extension it is a product of finite
local algebras \(D\) of the kind in the theorem. The maps
from \(T\) are continuous since the original quotient is
finite over \(O\). Each such factor is generated by the
image of \(T\) and the coefficient field, by projection
from the finite scalar extension. The two assertions follow
from the two conclusions of the theorem.
\end{proof}

\clearpage
\section{References}\label{nav:bibliography-VII}
\begingroup\def\section*#1{}\begin{thebibliography}{99}\small
\bibitem{VII-V-V-milne98} J. S. Milne, \emph{Arithmetic Duality Theorems}, 2nd ed., BookSurge, 2006. \href{https://www.jmilne.org/math/Books/ADTnot.pdf}{Part I, Sections 2 and 4}.
\bibitem{VII-V-I-ac123} J. Arthur and L. Clozel, \href{https://www.claymath.org/library/cw/arthur/pdf/30.pdf}{\emph{Simple Algebras, Base Change, and the Advanced Theory of the Trace Formula}}, Annals of Mathematics Studies 120, Princeton University Press, 1989, Chapter III.
\bibitem{VII-V-clozelrajan358} L. Clozel and C. S. Rajan, \emph{Solvable base change}, J. Reine Angew. Math., \href{https://doi.org/10.1515/crelle-2020-0023}{published paper}.
\bibitem{VII-V-taylor358} R. Taylor, \emph{The image of complex conjugation in \(\ell\)-adic representations associated to automorphic forms}, Algebra Number Theory \textbf{6} (2012), 405--435, Proposition 1. \href{https://msp.org/ant/2012/6-3/ant-v6-n3-p01-p.pdf}{Original paper}.
\bibitem{VII-V-clozel360} L. Clozel, \emph{Motifs et formes automorphes}, in \emph{Automorphic Forms, Shimura Varieties, and \(L\)-functions}, Vol.~I, Perspectives in Mathematics 10, Academic Press, 1990, Theorem 3.13. \href{https://www.jmilne.org/math/Books/AA1988a.pdf}{Original paper in the collected volume}.
\bibitem{VII-V-V-newforms} F. Diamond and J. Im, \emph{Modular forms and modular curves}, in \emph{Seminar on Fermat's Last Theorem}, CMS Conf. Proc. \textbf{17} (1995), 39--133. \href{https://wstein.org/edu/Fall2003/252/references/diamond-im/Diamond-Im-Modular_forms_and_modular_curves.pdf}{Authors' text}.
\bibitem{VII-V-billerey361} N. Billerey and R. Menares, \emph{On the modularity of reducible mod \(\ell\) Galois representations}, Math. Res. Lett. \textbf{23} (2016), 15--41, Section 1. \href{https://lmbp.uca.fr/~billerey/Research/Published/Menares_Modularity.pdf}{Published paper}.
\bibitem{VII-V-katz361} N. Katz, \emph{\(p\)-adic properties of modular schemes and modular forms}, in \emph{Modular Functions of One Variable III}, Lecture Notes in Mathematics 350, Springer, 1973. \href{https://web.math.princeton.edu/~nmk/old/padicpropMFMS.pdf}{Original paper}.
\bibitem{VII-V-rapoport361} P. Deligne and M. Rapoport, \emph{Les sch\'emas de modules de courbes elliptiques}, in \emph{Modular Functions of One Variable II}, Lecture Notes in Mathematics 349, Springer, 1973, 143--316. \href{https://publications.ias.edu/book/export/html/367}{Original paper}.
\bibitem{VII-V-V-deligne} P. Deligne, \emph{Formes modulaires et repr\'esentations \(\ell\)-adiques}, S\'eminaire Bourbaki, Exp. 355 (1969). \href{https://www.numdam.org/item/SB_1968-1969__11__139_0.pdf}{Paper}.
\bibitem{VII-V-ordinary361} C. M. Skinner and A. J. Wiles, \emph{Ordinary representations and modular forms}, Proc. Natl. Acad. Sci. USA \textbf{94} (1997), 10520--10527. \href{https://doi.org/10.1073/pnas.94.20.10520}{Published paper}.
\bibitem{VII-V-VI-saito40} T. Saito, \emph{Modular forms and \(p\)-adic Hodge theory}, Invent. Math. \textbf{129} (1997), 607--620. See also \href{https://www.ms.u-tokyo.ac.jp/~t-saito/pp/h.pdf}{\emph{Hilbert modular forms and \(p\)-adic Hodge theory}}, Compos. Math. \textbf{145} (2009), 1081--1113.
\bibitem{VII-V-I-compat123} H. Carayol, \href{https://www.numdam.org/articles/10.24033/asens.1512/}{\emph{Sur les repr\'esentations \(\ell\)-adiques associ\'ees aux formes modulaires de Hilbert}}, Ann. Sci. \'Ecole Norm. Sup. (4) \textbf{19} (1986), 409--468. T. Saito, \emph{Modular forms and \(p\)-adic Hodge theory}, Invent. Math. \textbf{129} (1997), 607--620.
\bibitem{VII-V-I-gj123} S. Gelbart and H. Jacquet, \href{https://www.numdam.org/item/ASENS_1978_4_11_4_471_0/}{\emph{A relation between automorphic representations of \(\mathrm{GL}(2)\) and \(\mathrm{GL}(3)\)}}, Ann. Sci. \'Ecole Norm. Sup. (4) \textbf{11} (1978), 471--542.
\bibitem{VII-V-V-shin} S. W. Shin, \emph{Galois representations arising from some compact Shimura varieties}, Ann. of Math. \textbf{173} (2011), 1645--1741. \href{https://annals.math.princeton.edu/2011/173-3/p09}{Theorem 1.2}.
\bibitem{VII-V-lgp361} T. Barnet-Lamb, T. Gee, D. Geraghty and R. Taylor, \emph{Local-global compatibility for \(l=p\), II}, Ann. Sci. \'Ec. Norm. Sup\'er. (4) \textbf{47} (2014), 165--179, Theorem A. \href{https://www.numdam.org/articles/10.24033/asens.2212/}{Published paper}.
\bibitem{VII-V-VI-deligne40} P. Deligne, \emph{La conjecture de Weil. I}, Publ. Math. IH\'ES \textbf{43} (1974), 273--307. \href{https://www.numdam.org/item/PMIHES_1974__43__273_0/}{Purity over finite fields}.
\bibitem{VII-V-VI-bk117} S. Bloch and K. Kato, \href{https://virtualmath1.stanford.edu/~conrad/BSDseminar/refs/BKTamagawa.pdf}{\emph{L-functions and Tamagawa numbers of motives}}, in \emph{The Grothendieck Festschrift}, Vol.~I, Progr. Math. \textbf{86} (1990), 333--400, Section~3.
\bibitem{VII-V-VI-nt119} J. Newton and J. A. Thorne, \href{https://arxiv.org/pdf/1912.11265}{\emph{Adjoint Selmer groups of automorphic Galois representations of unitary type}}, J. Eur. Math. Soc. \textbf{25} (2023), 1919--1967, Theorem~A and Section~2.5.
\bibitem{VII-V-geraghty363} D. Geraghty, \emph{Modularity lifting theorems for ordinary Galois representations}, Math. Ann. \textbf{373} (2019), 1341--1427, Section 2. \href{https://doi.org/10.1007/s00208-018-1742-4}{Published paper}. Proposition and lemma numbers refer to the \href{https://citeseerx.ist.psu.edu/document?doi=8faa2d675c94ae78d0ac8ebb3f9d4f69ede42d44\&repid=rep1\&type=pdf}{original manuscript}.
\bibitem{VII-V-brumer330} A. Brumer, \emph{On the units of algebraic number fields}, Mathematika \textbf{14} (1967), 121--124. \href{https://doi.org/10.1112/S0025579300003703}{Original paper}.
\bibitem{VII-V-chenevier364} G. Chenevier, \emph{The \(p\)-adic analytic space of pseudocharacters of a profinite group and pseudorepresentations over arbitrary rings}, in \emph{Automorphic Forms and Galois Representations}, Vol.~1, Cambridge University Press, 2014, 221--285, Introduction and Theorem A. \href{https://gaetan.chenevier.perso.math.cnrs.fr/articles/determinants.pdf}{Original paper}.
\bibitem{VII-V-II-fw13} B. Ferrero and L. C. Washington, \emph{The Iwasawa invariant \(\mu_p\) vanishes for abelian number fields}, Ann. of Math. \textbf{109} (1979), 377--395. \href{https://doi.org/10.2307/1971116}{Original paper}.
\bibitem{VII-V-II-iwasawa15} K. Iwasawa, \emph{On \(\ZZ_\ell\)-extensions of algebraic number fields}, Ann. of Math. \textbf{98} (1973), 246--326.
\bibitem{VII-V-hkp370} T. J. Haines, R. E. Kottwitz and A. Prasad, \emph{Iwahori--Hecke algebras}, J. Ramanujan Math. Soc. \textbf{25} (2010), 113--145, Sections 1 and 4. \href{https://arxiv.org/abs/math/0309168}{Original paper}.
\bibitem{VII-V-wangerickson375} C. Wang-Erickson, \emph{Algebraic families of Galois representations and potentially semi-stable pseudodeformation rings}, Math. Ann. \textbf{371} (2018), 1615--1681, Proposition 3.6. \href{https://doi.org/10.1007/s00208-017-1557-8}{Published paper}.
\bibitem{VII-V-liu375} T. Liu, \emph{Torsion \(p\)-adic Galois representations and a conjecture of Fontaine}, Ann. Sci. \'Ecole Norm. Sup. (4) \textbf{40} (2007), 633--674, Sections 1 and 7.3. \href{https://www.math.purdue.edu/~tongliu/pub/fontainel.pdf}{Original paper}.
\bibitem{VII-V-VI-pan40} L. Pan, \emph{The Fontaine--Mazur conjecture in the residually reducible case}, J. Amer. Math. Soc. \textbf{35} (2022), 1031--1169. \href{https://arxiv.org/abs/1901.07166}{Author's paper}.
\bibitem{VII-V-VI-kato120} K. Kato, \href{https://archive.numdam.org/item/AST_2004__295__117_0.pdf}{\emph{\(p\)-adic Hodge theory and values of zeta functions of modular forms}}, Ast\'erisque \textbf{295} (2004), 117--290, Theorem~14.2(1) and Section~14.10.
\end{thebibliography}\endgroup
\end{document}
